% Mesures, incertitudes et analyse dimensionnelle — Physique Tle D — Cours signé OnBuch+
% Programme officiel MINESEC. Compiler avec : tectonic P01-mesures-incertitudes-analyse-dimensionnelle.tex
\newcommand{\DOCMATIERE}{Physique}
\newcommand{\DOCNIVEAU}{Tle D}
\newcommand{\DOCMODULE}{Module 1 : Mesures et incertitudes}
\newcommand{\DOCLECON}{P01}
\newcommand{\DOCTITRE}{Mesures, incertitudes et analyse dimensionnelle}
% OnBuch+ — préambule « cours signés » ENRICHI (compilé avec tectonic / XeLaTeX)
% Système visuel « Pop » d'OnBuch+ : fond crème (jamais blanc), bordures encre
% épaisses, ombres « dures » décalées, titres Archivo Black en capitales.
% Version MATHÉMATIQUES : graphes (pgfplots/tikz), boîtes pédagogiques
% (définition, propriété, méthode, attention, le savais-tu, exercice résolu,
% exercice, TP, bilan), page de garde et signature OnBuch+ ancrée en bas.
%
% Macros attendues dans main.tex AVANT \input{preamble} :
%   \DOCMATIERE  \DOCNIVEAU  \DOCMODULE  \DOCLECON (numéro)  \DOCTITRE
\documentclass[11pt]{article}

\usepackage{fontspec}
\usepackage[french]{babel}
\usepackage[a4paper,margin=2cm,top=3.4cm,bottom=2.8cm,headheight=1.4cm,footskip=1.3cm]{geometry}
\usepackage{amsmath,amssymb,amsfonts}
\usepackage{mathtools} % \xmapsto, \coloneqq... souvent utilisés par les modèles
\usepackage{cancel} % simplifications barrées (bilans de forces)
% \abs{z} (module, valeur absolue) et \norm{u} (norme), utilisés par les modèles
\providecommand{\abs}{}\let\abs\relax\DeclarePairedDelimiter{\abs}{\lvert}{\rvert}
\providecommand{\norm}{}\let\norm\relax\DeclarePairedDelimiter{\norm}{\lVert}{\rVert}
\usepackage[table]{xcolor} % [table] : fournit aussi \rowcolors (alternance de lignes)
\usepackage{pagecolor}
\usepackage{tikz}
\usetikzlibrary{arrows.meta,positioning,calc,shapes.geometric,shapes.misc,shapes.arrows,decorations.pathmorphing,decorations.markings,patterns,shadows,mindmap,trees,fit,backgrounds,matrix,intersections,angles,quotes}
\usepackage{pgfplots}
\usepackage[version=4]{mhchem} % équations nucléaires \ce{^{235}_{92}U}
\usepackage[european,siunitx]{circuitikz} % schémas électriques
\pgfplotsset{compat=1.18}
\usepgfplotslibrary{fillbetween,groupplots} % aires entre courbes (calcul intégral)
\usepackage{enumitem}
\usepackage{fancyhdr}
% [most] charge toutes les bibliothèques tcolorbox (listings, minted, raster,
% documentation...) et gonfle inutilement la mémoire TeX sur un document long
% (~300 boîtes) : on ne charge que ce qui est réellement utilisé.
\usepackage[skins,breakable]{tcolorbox}
\usepackage{graphicx}
\usepackage{colortbl}
\usepackage{booktabs}
\usepackage{tabularx}
\usepackage{diagbox} % case d'en-tête diagonale des tableaux à double entrée
% Colonnes extensibles centrée / gauche / droite (C, L, R), souvent utilisées par les modèles
\newcolumntype{C}{>{\centering\arraybackslash}X}
\newcolumntype{L}{>{\raggedright\arraybackslash}X}
\newcolumntype{R}{>{\raggedleft\arraybackslash}X}
\usepackage{array}
\usepackage{multicol}
\usepackage{siunitx}
% parse-numbers=false : accepte aussi \SI{2,5 \times 10^{-3}}{...}
\sisetup{locale=FR,output-decimal-marker={,},group-minimum-digits=4,parse-numbers=false}
\DeclareSIUnit\ohm{\ensuremath{\symup{\Omega}}} % U+2126 absent de la police
\DeclareSIUnit\division{div} % réglages d'oscilloscope (ms/div, V/div)
\usepackage{qrcode}
% \degree : commande fréquemment utilisée par les modèles pour un angle ou une
% température en dehors de \SI{}{} (non fournie par les paquets chargés).
\providecommand{\degree}{^{\circ}}
% \qed (fin de démonstration), utilisé par les modèles sans amsthm
\providecommand{\qed}{\hfill$\square$}
% \barre : double barre des valeurs interdites dans un tableau de variations (tkz-tab)
\providecommand{\barre}{\ensuremath{\|}}
% environnement proof (amsthm non chargé) utilisé par les modèles
\ifdefined\proof\else\newenvironment{proof}[1][Démonstration]{\par\noindent\textit{#1.}\ }{\qed\par}\fi
\usepackage{lastpage}
\usepackage{hyperref}
\hypersetup{hidelinks}

% ============================================================
% Palette « Pop » OnBuch+ — mêmes hex que la classe P (plus_ui.dart)
% ============================================================
\definecolor{poporange}{HTML}{F59321}
\definecolor{popdark}{HTML}{E07A0C}
\definecolor{popcream}{HTML}{FFF6E8}
\definecolor{popink}{HTML}{1C1714}
\definecolor{poppurple}{HTML}{7A5AE0}
\definecolor{popgreen}{HTML}{2E9E6A}
\definecolor{poppink}{HTML}{E0457B}
\definecolor{popblue}{HTML}{2F7DE1}
\definecolor{popgold}{HTML}{C98A12}
\definecolor{popmuted}{HTML}{8A7F76}
\definecolor{popline}{HTML}{EADFD2}
\definecolor{popgray}{HTML}{8A7F76}
\colorlet{popgrey}{popgray}
% Alias tolérants (le modèle nomme parfois les couleurs différemment)
\colorlet{popwhite}{white}
\colorlet{popblack}{popink}
\colorlet{popred}{poppink}
\colorlet{popyellow}{popgold}
% Teintes claires pour fonds de tableaux / zones
\colorlet{popblueL}{popblue!10!white}
\colorlet{poporangeL}{poporange!14!white}
\colorlet{popgreenL}{popgreen!12!white}
\colorlet{poppurpleL}{poppurple!12!white}
\colorlet{poppinkL}{poppink!12!white}
\colorlet{popgoldL}{popgold!14!white}

\pagecolor{popcream}

% ============================================================
% Typographie
% ============================================================
\defaultfontfeatures{Ligatures=TeX}
\setmainfont{PlusJakartaSans-Medium.ttf}[Path=../../../fonts/,BoldFont=PlusJakartaSans-Bold.ttf]
% Maths en sans-serif (Fira Math) avec chiffres et lettres droites de Plus
% Jakarta, pour que les formules se fondent dans le texte.
\usepackage[math-style=ISO,bold-style=ISO]{unicode-math}
\setmathfont{FiraMath-Regular.otf}[Path=../../../fonts/]
\setmathfont{PlusJakartaSans-Medium.ttf}[Path=../../../fonts/,range={up/num,up/{Latin,latin},\mathup}]
\usepackage{icomma} % virgule décimale sans espace en mode math
% Caractères mathématiques Unicode absents des polices de texte (Plus Jakarta,
% Archivo Black) : rendus par leur équivalent mathématique, en texte comme en
% formule — sinon ils s'affichent en carré vide (ex. « Arithmétique dans ℤ »).
\usepackage{newunicodechar}
\newunicodechar{ℝ}{\ensuremath{\mathbb{R}}}
\newunicodechar{ℤ}{\ensuremath{\mathbb{Z}}}
\newunicodechar{ℂ}{\ensuremath{\mathbb{C}}}
\newunicodechar{ℕ}{\ensuremath{\mathbb{N}}}
\newunicodechar{ℚ}{\ensuremath{\mathbb{Q}}}
\newunicodechar{𝔻}{\ensuremath{\mathbb{D}}}
\newunicodechar{α}{\ensuremath{\alpha}}
\newunicodechar{β}{\ensuremath{\beta}}
\newunicodechar{γ}{\ensuremath{\gamma}}
\newunicodechar{δ}{\ensuremath{\delta}}
\newunicodechar{ε}{\ensuremath{\varepsilon}}
\newunicodechar{θ}{\ensuremath{\theta}}
\newunicodechar{λ}{\ensuremath{\lambda}}
\newunicodechar{μ}{\ensuremath{\mu}}
\newunicodechar{σ}{\ensuremath{\sigma}}
\newunicodechar{τ}{\ensuremath{\tau}}
\newunicodechar{φ}{\ensuremath{\varphi}}
\newunicodechar{ψ}{\ensuremath{\psi}}
\newunicodechar{ω}{\ensuremath{\omega}}
\newunicodechar{Σ}{\ensuremath{\Sigma}}
% majuscules produites par \MakeUppercase dans les titres (α-aminés -> Α)
\newunicodechar{Α}{\ensuremath{\alpha}}
\newunicodechar{Β}{\ensuremath{\beta}}
\newunicodechar{Γ}{\ensuremath{\Gamma}}
\newunicodechar{Δ}{\ensuremath{\Delta}}
\newunicodechar{Ε}{\ensuremath{\varepsilon}}
\newunicodechar{Θ}{\ensuremath{\Theta}}
\newunicodechar{Λ}{\ensuremath{\Lambda}}
\newunicodechar{Μ}{\ensuremath{\mu}}
\newunicodechar{Τ}{\ensuremath{\tau}}
\newunicodechar{Φ}{\ensuremath{\Phi}}
\newunicodechar{Ψ}{\ensuremath{\Psi}}
\newunicodechar{Ω}{\ensuremath{\Omega}}
\newunicodechar{↦}{\ensuremath{\mapsto}}
\newunicodechar{∈}{\ensuremath{\in}}
\newunicodechar{∉}{\ensuremath{\notin}}
\newunicodechar{∧}{\ensuremath{\wedge}}
\newunicodechar{⇔}{\ensuremath{\Leftrightarrow}}
\newunicodechar{⟺}{\ensuremath{\Longleftrightarrow}}
\newunicodechar{⇒}{\ensuremath{\Rightarrow}}
\newunicodechar{⟹}{\ensuremath{\Longrightarrow}}
\newunicodechar{∘}{\ensuremath{\circ}}
\newunicodechar{∪}{\ensuremath{\cup}}
\newunicodechar{∩}{\ensuremath{\cap}}
\newunicodechar{≡}{\ensuremath{\equiv}}
\newunicodechar{⊂}{\ensuremath{\subset}}
\newunicodechar{⊥}{\ensuremath{\perp}}
\newunicodechar{∃}{\ensuremath{\exists}}
\newunicodechar{∀}{\ensuremath{\forall}}
\newunicodechar{∛}{\ensuremath{\sqrt[3]{\,}}}
\newunicodechar{∜}{\ensuremath{\sqrt[4]{\,}}}
\newunicodechar{⃗}{\ensuremath{{}^{\to}}}
\newunicodechar{ⁿ}{\ifmmode^{n}\else\textsuperscript{\ensuremath{n}}\fi}
\newunicodechar{ˣ}{\ifmmode^{x}\else\textsuperscript{\ensuremath{x}}\fi}
\newunicodechar{ᵇ}{\ifmmode^{b}\else\textsuperscript{\ensuremath{b}}\fi}
\newunicodechar{ʳ}{\ifmmode^{r}\else\textsuperscript{\ensuremath{r}}\fi}
\newunicodechar{ᵘ}{\ifmmode^{u}\else\textsuperscript{\ensuremath{u}}\fi}
\newunicodechar{ʸ}{\ifmmode^{y}\else\textsuperscript{\ensuremath{y}}\fi}
\newunicodechar{ᵃ}{\ifmmode^{a}\else\textsuperscript{\ensuremath{a}}\fi}
\newunicodechar{ᵉ}{\ifmmode^{e}\else\textsuperscript{\ensuremath{e}}\fi}
\newunicodechar{ᵏ}{\ifmmode^{k}\else\textsuperscript{\ensuremath{k}}\fi}
\newunicodechar{ᵖ}{\ifmmode^{p}\else\textsuperscript{\ensuremath{p}}\fi}
\newunicodechar{ᴺ}{\ifmmode^{N}\else\textsuperscript{\ensuremath{N}}\fi}
\newunicodechar{ᶜ}{\ifmmode^{c}\else\textsuperscript{\ensuremath{c}}\fi}
\newunicodechar{ʰ}{\ifmmode^{h}\else\textsuperscript{\ensuremath{h}}\fi}
\newunicodechar{ᵠ}{\ifmmode^{q}\else\textsuperscript{\ensuremath{q}}\fi}
\newunicodechar{ₙ}{\ifmmode_{n}\else\textsubscript{\ensuremath{n}}\fi}
\newunicodechar{ₐ}{\ifmmode_{a}\else\textsubscript{\ensuremath{a}}\fi}
\newunicodechar{ᵢ}{\ifmmode_{i}\else\textsubscript{\ensuremath{i}}\fi}
\newunicodechar{ᵣ}{\ifmmode_{r}\else\textsubscript{\ensuremath{r}}\fi}
\newunicodechar{ₖ}{\ifmmode_{k}\else\textsubscript{\ensuremath{k}}\fi}
\newunicodechar{ᵦ}{\ifmmode_{\beta}\else\textsubscript{\ensuremath{\beta}}\fi}
\newunicodechar{ₓ}{\ifmmode_{x}\else\textsubscript{\ensuremath{x}}\fi}
\newfontfamily\popdisplay{ArchivoBlack-Regular.ttf}[Path=../../../fonts/]
\newfontfamily\poplogo{SpaceGrotesk-Bold.ttf}[Path=../../../fonts/]
\newfontfamily\popsemi{PlusJakartaSans-SemiBold.ttf}[Path=../../../fonts/]
\newfontfamily\popxbold{PlusJakartaSans-ExtraBold.ttf}[Path=../../../fonts/]
\newfontfamily\popxboldls{PlusJakartaSans-ExtraBold.ttf}[Path=../../../fonts/,LetterSpace=3.0]

\setlist[enumerate]{leftmargin=1.7em,itemsep=5pt,topsep=4pt}
\setlist[itemize]{leftmargin=1.7em,itemsep=4pt,topsep=4pt,label={\color{poporange}\textbullet}}
\setlength{\parindent}{0pt}
\setlength{\parskip}{5pt}
\renewcommand{\arraystretch}{1.25}

% pgfplots : style Pop par défaut
\pgfplotsset{
  every axis/.append style={axis line style={popink,line width=1pt},tick style={popink},
    grid style={popline,line width=0.5pt},label style={font=\small},tick label style={font=\footnotesize}},
  popcurve/.style={poporange,line width=1.8pt,no markers,smooth},
}
% Même style disponible dans un \draw TikZ simple (hors pgfplots)
\tikzset{popcurve/.style={poporange,line width=1.8pt}}
\tikzset{popgrid/.style={popline,very thin}}
\colorlet{popcurve}{poporange} % le nom du style est parfois utilisé comme couleur

% ============================================================
% Pill / squiggle
% ============================================================
\newcommand{\poppill}[3]{%
  \tikz[baseline=-0.55ex]\node[draw=popink,line width=1.2pt,rounded corners=2.6pt,%
    fill=#2,inner xsep=6.5pt,inner ysep=3pt]{%
    \popxboldls\fontsize{7}{7}\selectfont\color{#3}\MakeUppercase{#1}};%
}
\newcommand{\popsquiggle}[1]{%
  \tikz[baseline=-0.4ex]{
    \draw[#1,line width=2.6pt,line cap=round]
      plot [smooth] coordinates {(0,0.09) (0.34,0.01) (0.68,0.21) (1.02,0.01) (1.36,0.10)};
  }%
}
% Mot-clé surligné (marqueur orange sous le texte)
\newcommand{\cle}[1]{{\popxbold\color{popdark}#1}}

% ============================================================
% En-tête / pied
% ============================================================
\pagestyle{fancy}
\fancyhf{}
\renewcommand{\headrulewidth}{0pt}
\renewcommand{\footrulewidth}{0pt}
\lhead{\raisebox{-0.15ex}{\poplogo\fontsize{13}{13}\selectfont\color{popink}OnBuch\color{poporange}+}}
\rhead{\poppill{\DOCMATIERE\ \textbullet\ \DOCNIVEAU\ \textbullet\ Leçon \DOCLECON}{white}{popink}}
\lfoot{\popxbold\fontsize{7.6}{7.6}\selectfont\color{popmuted}\MakeUppercase{Cours signé OnBuch+}\ \textbullet\ {\popsemi\DOCTITRE}}
\rfoot{\tikz[baseline=-0.6ex]\node[rounded corners=8.5pt,fill=popink,minimum height=17pt,minimum width=17pt,inner xsep=5pt,inner sep=0]{\popxbold\fontsize{8}{8}\selectfont\color{popcream}\thepage\,/\,\pageref*{LastPage}};}

% ============================================================
% Titres
% ============================================================
\newcommand{\chaptitre}[2]{%
  \begin{tcolorbox}[enhanced,colback=poporange,colframe=popink,boxrule=2.6pt,arc=11pt,
      drop shadow=popink,left=15pt,right=15pt,top=12pt,bottom=11pt,before skip=3pt,after skip=18pt]
    \poppill{#2}{popink}{popcream}\par\vspace{9pt}
    {\raggedright\hyphenpenalty=10000\exhyphenpenalty=10000\popdisplay\fontsize{22}{24}\selectfont\color{popcream}\MakeUppercase{#1}\par}
  \end{tcolorbox}
}
% Section numérotée : pastille orange avec le numéro + titre Archivo
\newcounter{popsec}
\newcommand{\coursec}[1]{%
  \stepcounter{popsec}\setcounter{popsub}{0}%
  \par\vspace{16pt}\noindent
  \begin{minipage}{\linewidth}
  \tikz[baseline=-0.7ex]\node[draw=popink,line width=1.6pt,fill=poporange,rounded corners=4pt,minimum size=22pt,inner sep=0]{\popdisplay\fontsize{12}{12}\selectfont\color{popcream}\thepopsec};%
  \hspace{8pt}{\popdisplay\fontsize{15}{17}\selectfont\color{popink}\MakeUppercase{#1}}\par
  \vspace{4pt}\noindent{\color{poporange}\rule{36pt}{3.6pt}}
  \end{minipage}\par\nopagebreak\vspace{8pt}
}
\newcounter{popsub}
\newcommand{\courssub}[1]{%
  \stepcounter{popsub}%
  \par\vspace{9pt}\noindent{\popxbold\fontsize{11.5}{13}\selectfont\color{popink}{\color{poporange}\thepopsec.\thepopsub}\hspace{6pt}#1}\par\nopagebreak\vspace{4pt}
}

% ============================================================
% Boîtes pédagogiques — même recette : fond blanc, bordure épaisse colorée,
% ombre dure encre, badge plein. Titre optionnel [#1].
% ============================================================
\tcbset{popbase/.style={breakable,enhanced,colback=white,boxrule=2.3pt,arc=9pt,
  shadow={3pt}{-3pt}{0pt}{popink},left=14pt,right=14pt,top=11pt,bottom=11pt,before skip=11pt,after skip=15pt,
  fontupper=\popsemi\fontsize{10.6}{14.5}\selectfont\color{popink},
  fontlower=\popsemi\fontsize{10.6}{14.5}\selectfont\color{popink}}}
% Coche dessinée (le glyphe ✓ n'existe pas dans les polices du document)
\newcommand{\popcheck}[1]{\tikz[baseline=-0.5ex]\draw[#1,line width=1.6pt,line cap=round,line join=round](-0.13,0.0)--(-0.04,-0.09)--(0.13,0.1);}
\providecommand{\coche}{\popcheck{popgreen}}
\providecommand{\resultat}[1]{\par\smallskip\noindent\fcolorbox{popgreen}{popgreenL}{\parbox{\dimexpr\linewidth-2\fboxsep-2\fboxrule\relax}{\textbf{Résultat.} #1}}\par\smallskip}
\newcommand{\popboxhead}[3]{\poppill{#1}{#2}{white}\ifx\relax#3\relax\else\hspace{6pt}{\popxbold\fontsize{10.6}{12}\selectfont\color{popink}#3}\fi\par\vspace{7pt}}

\newtcolorbox{aretenir}[1][]{popbase,colframe=popgreen,before upper={\popboxhead{À retenir}{popgreen}{#1}}}
\newtcolorbox{exemplebox}[1][]{popbase,colframe=poporange,before upper={\popboxhead{Exemple}{poporange}{#1}}}
\newtcolorbox{definition}[1][]{popbase,colframe=popblue,before upper={\popboxhead{Définition}{popblue}{#1}}}
\newtcolorbox{propriete}[1][]{popbase,colframe=poppurple,before upper={\popboxhead{Propriété}{poppurple}{#1}}}
\newtcolorbox{methode}[1][]{popbase,colframe=popgold,before upper={\popboxhead{Méthode}{popgold}{#1}}}
\newtcolorbox{attention}[1][]{popbase,colframe=poppink,before upper={\popboxhead{Attention}{poppink}{#1}}}
\newtcolorbox{experience}[1][]{popbase,colframe=popblue,colback=popblueL,before upper={\popboxhead{Expérience}{popblue}{#1}}}
% « Le savais-tu ? » : carte violette pleine (contexte, culture, Cameroun)
\newtcolorbox{savaistu}[1][]{popbase,colback=poppurple,colframe=popink,
  fontupper=\popsemi\fontsize{10.4}{14.2}\selectfont\color{white},
  before upper={\poppill{Le savais-tu\,?}{white}{poppurple}\ifx\relax#1\relax\else\hspace{6pt}{\popxbold\color{white}#1}\fi\par\vspace{7pt}}}
% Exercice résolu : énoncé en haut, \tcblower puis la solution sur fond crème
\newcounter{popexo}\newcounter{popexr}
\newtcolorbox{exoresolu}[1][]{popbase,colframe=poporange,colbacklower=popcream,
  segmentation style={popink,line width=1.2pt,dashed},
  before upper={\stepcounter{popexr}\popboxhead{Exercice résolu \thepopexr}{poporange}{#1}},
  before lower={\poppill{Solution}{popink}{popcream}\par\vspace{6pt}}}
% Exercice (non résolu en place) — la correction est dans la partie Corrigés
\newtcolorbox{exercice}[1][]{popbase,colframe=popink,
  before upper={\stepcounter{popexo}\popboxhead{Exercice \thepopexo}{popink}{#1}}}
% Corrigé
\newtcolorbox{corrige}[1][]{popbase,colframe=popgreen,colback=popgreenL,
  before upper={\popboxhead{Corrigé}{popgreen}{#1}}}
% Objectifs / prérequis / bilan
\newtcolorbox{objectifs}{popbase,colframe=popink,colback=poporangeL,
  before upper={\popboxhead{Objectifs de la leçon}{popink}{}}}
\newtcolorbox{prerequis}{popbase,colframe=popink,colback=popblueL,
  before upper={\popboxhead{Prérequis}{popblue}{}}}
\newtcolorbox{bilan}{popbase,colframe=popink,colback=white,boxrule=2.8pt,
  before upper={\popboxhead{Fiche bilan}{poporange}{L'essentiel à connaître pour l'examen}}}
% Figure encadrée avec légende
\newtcolorbox{popfigure}[1][]{enhanced,breakable=false,colback=white,colframe=popline,boxrule=1.4pt,arc=8pt,
  left=10pt,right=10pt,top=10pt,bottom=8pt,before skip=10pt,after skip=12pt,halign=center,
  lower separated=false,halign lower=center,
  fontlower=\popsemi\fontsize{9}{11.5}\selectfont\color{popmuted},
  #1}
\newcommand{\legende}[1]{\par\vspace{6pt}{\popsemi\fontsize{9}{11.5}\selectfont\color{popmuted}#1}}

% ============================================================
% Signature OnBuch+
% ============================================================
\newcommand{\poplogoinline}[1]{{\poplogo\fontsize{#1}{#1}\selectfont\color{popink}OnBuch\color{poporange}+}}
\newcommand{\signatureonbuch}{%
  \begin{tcolorbox}[enhanced,colback=popink,colframe=popink,boxrule=2.2pt,arc=10pt,
      drop shadow=poporange,left=14pt,right=14pt,top=10pt,bottom=10pt,before skip=0pt,after skip=0pt]
    \tikz[baseline=-0.7ex]\node[circle,fill=popgreen,draw=popcream,line width=1.3pt,minimum size=22pt,inner sep=0]{\popcheck{white}};%
    \hspace{9pt}\begin{minipage}[c]{0.62\linewidth}
      {\popxbold\fontsize{12}{13}\selectfont\color{popcream}Signé\ \textbullet\ {\poplogo OnBuch\color{poporange}+}}\par\vspace{2pt}
      {\raggedright\popsemi\fontsize{8}{10.5}\selectfont\color{popline}\MakeUppercase{Équipe pédagogique OnBuch+}\\\MakeUppercase{Conforme au programme officiel MINESEC}\par}
    \end{minipage}\hfill
    {\poplogo\fontsize{22}{22}\selectfont\color{popcream}OnBuch\color{poporange}+}
  \end{tcolorbox}%
}

% ============================================================
% Page de garde
% #1 titre · #2 matière · #3 niveau · #4 module · #5 n° leçon · #6 durée ·
% #7 description / objectifs (texte ou itemize)
% ============================================================
\newcommand{\pagegarde}[7]{%
  \thispagestyle{empty}
  \newgeometry{a4paper,margin=1.7cm,top=1.6cm,bottom=1.4cm}%
  \begin{tikzpicture}[remember picture,overlay]
    \fill[poporange!18] ([xshift=-3.2cm,yshift=-3.6cm]current page.north east) circle (4.6cm);
    \fill[poppurple!14] ([xshift=2.4cm,yshift=7.6cm]current page.south west) circle (3.8cm);
  \end{tikzpicture}%
  {\poplogo\fontsize{30}{30}\selectfont\color{popink}OnBuch\color{poporange}+}\hfill
  \raisebox{0.6ex}{\poppill{Cours signé \textbullet\ Premium}{popink}{popcream}}\par
  \vspace{1pt}\popsquiggle{poppurple}\par
  \vspace{8pt}
  \begin{tcolorbox}[enhanced,colback=poporange,colframe=popink,boxrule=2.8pt,arc=13pt,
      drop shadow=popink,left=18pt,right=18pt,top=12pt,bottom=13pt,after skip=10pt]
    \poppill{#2 \textbullet\ #3}{popink}{popcream}\hspace{5pt}\poppill{#4}{white}{popink}\par\vspace{8pt}
    {\popdisplay\fontsize{14}{14}\selectfont\color{popink}LEÇON #5}\par\vspace{4pt}
    {\raggedright\hyphenpenalty=10000\exhyphenpenalty=10000\popdisplay\fontsize{25}{27}\selectfont\color{popcream}\MakeUppercase{#1}\par}
    \vspace{8pt}
    {\popxbold\fontsize{9.5}{9.5}\selectfont\color{popink}\MakeUppercase{Durée conseillée : #6}}
  \end{tcolorbox}
  \begin{tcolorbox}[enhanced,colback=white,colframe=popink,boxrule=2.2pt,arc=10pt,
      drop shadow=popink,left=16pt,right=16pt,top=10pt,bottom=10pt,after skip=8pt]
    \poppill{Dans cette leçon}{poppurple}{white}\par\vspace{5pt}
    \popsemi\fontsize{9.3}{12.6}\selectfont\color{popink}#7
  \end{tcolorbox}
  \noindent\begin{minipage}[t]{0.487\linewidth}
    \begin{tcolorbox}[enhanced,colback=popgreen,colframe=popink,boxrule=2.2pt,arc=10pt,
        drop shadow=popink,left=11pt,right=11pt,top=10pt,bottom=10pt,height=3.1cm,valign=center]
      \tcbox[colback=white,colframe=popink,boxrule=1.3pt,arc=5pt,boxsep=0pt,nobeforeafter,
        left=3pt,right=3pt,top=3pt,bottom=3pt,valign=center]{\qrcode[height=1.9cm]{https://play.google.com/store/apps/details?id=com.luvvix.onbuch&hl=fr}}%
      \hspace{8pt}\begin{minipage}[c]{0.5\linewidth}
        {\popdisplay\fontsize{11}{12.5}\selectfont\color{white}\MakeUppercase{Télécharge OnBuch}}\par\vspace{4pt}
        {\raggedright\popsemi\fontsize{8.4}{11}\selectfont\color{white}Gratuit \textbullet\ Google Play\\Tuteur IA, annales, concours, résultats\par}
      \end{minipage}
    \end{tcolorbox}
  \end{minipage}\hfill
  \begin{minipage}[t]{0.487\linewidth}
    \begin{tcolorbox}[enhanced,colback=poppurple,colframe=popink,boxrule=2.2pt,arc=10pt,
        drop shadow=popink,left=11pt,right=11pt,top=10pt,bottom=10pt,height=3.1cm,valign=center]
      \tikz[baseline=-0.7ex]\node[circle,fill=white,draw=popink,line width=1.3pt,minimum size=1.6cm,inner sep=0]{\popdisplay\fontsize{24}{24}\selectfont\color{poppurple}+};%
      \hspace{8pt}\begin{minipage}[c]{0.6\linewidth}
        {\popdisplay\fontsize{11}{12.5}\selectfont\color{white}\MakeUppercase{Passe à OnBuch+}}\par\vspace{4pt}
        {\raggedright\popsemi\fontsize{8.4}{11}\selectfont\color{white}Tous les cours signés, Tuteur IA illimité, fiches PDF — onglet OnBuch+ de l'app\par}
      \end{minipage}
    \end{tcolorbox}
  \end{minipage}
  \vspace{8pt}
  \popcontenu
  \vfill
  \signatureonbuch
  \restoregeometry
  \newpage
}

% Rangée « Ce cours contient » (page de garde)
\newcommand{\poptile}[3]{% #1 couleur #2 gros texte #3 légende
  \begin{tcolorbox}[enhanced,colback=#1,colframe=popink,boxrule=2pt,arc=9pt,shadow={2.5pt}{-2.5pt}{0pt}{popink},
      left=6pt,right=6pt,top=9pt,bottom=9pt,halign=center,valign=center,height=1.75cm,width=\linewidth,nobeforeafter]
    {\popdisplay\fontsize{12.5}{13}\selectfont\color{white}\MakeUppercase{#2}}\par\vspace{3pt}
    {\centering\popsemi\fontsize{7.8}{9.5}\selectfont\color{white}#3\par}
  \end{tcolorbox}}
\newcommand{\popcontenu}{%
  \par\noindent{\popxboldls\fontsize{7.5}{7.5}\selectfont\color{popmuted}CE COURS SIGNÉ CONTIENT}\par\vspace{7pt}
  \noindent\begin{minipage}[t]{0.235\linewidth}\poptile{popblue}{Cours}{complet, illustré, pas à pas}\end{minipage}\hfill
  \begin{minipage}[t]{0.235\linewidth}\poptile{poporange}{Méthodes}{et exercices résolus}\end{minipage}\hfill
  \begin{minipage}[t]{0.235\linewidth}\poptile{poppink}{Exercices}{type examen + corrigés}\end{minipage}\hfill
  \begin{minipage}[t]{0.235\linewidth}\poptile{popgreen}{Fiche bilan}{l'essentiel à retenir}\end{minipage}\par}

% Sommaire
\newtcolorbox{sommaire}{enhanced,colback=white,colframe=popink,boxrule=2.2pt,arc=10pt,
  drop shadow=popink,left=18pt,right=18pt,top=13pt,bottom=13pt,before skip=6pt,after skip=20pt,
  before upper={\poppill{Sommaire}{poppurple}{white}\par\vspace{9pt}\popsemi\fontsize{10.6}{14.5}\selectfont\color{popink}}}

% Signature de fin de cours — poussée tout en bas de la dernière page
\newcommand{\findecours}{%
  \par\vfill\nopagebreak
  \begin{center}\popsquiggle{poporange}\end{center}\vspace{4pt}
  \signatureonbuch
}
\AtBeginDocument{\def\llbracket{\mathopen{[\mkern-3.5mu[}}\def\rrbracket{\mathclose{]\mkern-3.5mu]}}} % intervalles d'entiers (glyphe ⟦ absent de la police)
% Glyphes absents de Fira Math : redéfinis à partir de glyphes disponibles
\AtBeginDocument{%
  \def\setminus{\mathbin{\backslash}}\let\smallsetminus\setminus
  \def\vdots{\mathord{\vbox{\baselineskip3.2pt\lineskiplimit0pt\kern4pt\hbox{.}\hbox{.}\hbox{.}}}}%
  \def\ddots{\mathinner{\mkern1mu\raise6.5pt\hbox{.}\mkern2mu\raise3.6pt\hbox{.}\mkern2mu\raise0.7pt\hbox{.}\mkern1mu}}%
  \def\triangle{\mathord{\scalebox{0.9}{$\symup{\Delta}$}}}%
  \def\nabla{\mathord{\raisebox{\depth}{\scalebox{1}[-1]{$\symup{\Delta}$}}}}%
  \def\checkmark{\popcheck{popdark}}%
  \def\star{\mathbin{\ast}}\def\bigstar{\mathord{\ast}}%
  \def\diamond{\mathbin{\scalebox{0.6}{$\Diamond$}}}%
}

%% FIGFIX-BEGIN v5 — correctifs globaux des figures (docs/onbuch-generation/tools/figfix)
\usepackage{adjustbox}\usepackage{etoolbox}\tcbuselibrary{raster}
% toute figure TikZ/pgfplots plus large que la ligne est réduite à la largeur disponible
\BeforeBeginEnvironment{tikzpicture}{\begin{adjustbox}{max width=\linewidth}}
\AfterEndEnvironment{tikzpicture}{\end{adjustbox}}
% légendes pgfplots lisibles par-dessus les courbes
\pgfplotsset{every axis/.append style={legend style={fill=white,fill opacity=0.92,text opacity=1,draw=black!15,inner sep=3pt}}}
% étiquettes d'arêtes (syntaxe edge["..."]) avec fond blanc
\tikzset{every edge quotes/.append style={fill=white,inner sep=1.2pt}}
%% FIGFIX-END
\begin{document}

\pagegarde{Mesures, incertitudes et analyse dimensionnelle}{Physique}{Tle D}{Module 1 : Mesures et incertitudes}{P01}{6 h}{Cette leçon fondatrice donne à l'élève de Terminale C les outils du physicien pour mesurer, estimer les incertitudes et exprimer rigoureusement un résultat expérimental. Elle introduit également l'analyse dimensionnelle, instrument puissant de vérification et de découverte des lois physiques, utilisée tout au long de l'année.
\vspace{4pt}
\begin{multicols}{2}\small
\begin{itemize}
\item À la fin de cette leçon, je sais distinguer erreur et incertitude sur une mesure
\item À la fin de cette leçon, je sais caractériser la qualité d'un instrument (sensibilité, précision, fidélité, justesse)
\item À la fin de cette leçon, je sais estimer l'incertitude d'une mesure unique (lecture, constructeur)
\item À la fin de cette leçon, je sais traiter une série de mesures : moyenne, écart-type, incertitude-type
\item À la fin de cette leçon, je sais exprimer un résultat sous la forme valeur ± incertitude avec les chiffres significatifs corrects et l'incertitude relative
\item À la fin de cette leçon, je sais utiliser le Système international : unités de base, unités dérivées, multiples et sous-multiples
\end{itemize}\end{multicols}}

\begin{sommaire}
\begin{multicols}{2}
\begin{enumerate}
\item Mesure, erreur et qualité d'un instrument
\item Incertitude d'une mesure unique
\item Série de mesures : moyenne, écart-type, incertitude-type
\item Expression d'un résultat de mesure
\item Le Système international d'unités (SI)
\item Analyse dimensionnelle et homogénéité
\item Activité d'intégration
\item Méthodes et exercices résolus
\item Exercices
\item Corrigés des exercices
\item Fiche bilan
\end{enumerate}
\end{multicols}
\end{sommaire}

\begin{objectifs}
\begin{itemize}
\item À la fin de cette leçon, je sais distinguer erreur et incertitude sur une mesure
\item À la fin de cette leçon, je sais caractériser la qualité d'un instrument (sensibilité, précision, fidélité, justesse)
\item À la fin de cette leçon, je sais estimer l'incertitude d'une mesure unique (lecture, constructeur)
\item À la fin de cette leçon, je sais traiter une série de mesures : moyenne, écart-type, incertitude-type
\item À la fin de cette leçon, je sais exprimer un résultat sous la forme valeur ± incertitude avec les chiffres significatifs corrects et l'incertitude relative
\item À la fin de cette leçon, je sais utiliser le Système international : unités de base, unités dérivées, multiples et sous-multiples
\item À la fin de cette leçon, je sais établir la dimension d'une grandeur physique du programme (accélération, champ, capacité, inductance, impédance, activité...)
\item À la fin de cette leçon, je sais vérifier l'homogénéité d'une relation et l'utiliser pour valider une expression ou proposer une hypothèse
\end{itemize}
\end{objectifs}

\begin{prerequis}
\begin{itemize}
\item Grandeurs physiques et unités usuelles vues en Première (longueur, masse, temps, intensité, tension)
\item Notions de moyenne et d'écart en mathématiques (statistiques de Seconde/Première)
\item Écriture scientifique et puissances de 10
\item Lecture d'appareils de mesure simples (règle, balance, voltmètre, chronomètre)
\end{itemize}
\end{prerequis}

\newpage

\coursec{Mesure, erreur et qualité d'un instrument}

Au marché Mokolo de Yaoundé, deux commerçants pèsent le même sac de riz : l'un annonce \SI{25,2}{kg}, l'autre \SI{24,8}{kg}. Qui a raison ? Les deux balances affichent pourtant des chiffres bien nets. À Maroua, un technicien mesure le débit d'un forage avant que la commune n'investisse des millions de FCFA dans un château d'eau : s'il se trompe de \SI{20}{\%}, le projet entier est compromis. Au stade de Bafoussam, le chrono officiel donne \SI{10,42}{s} sur le \SI{100}{m} : peut-on départager deux athlètes à un centième près ?

Dans tous ces cas, un nombre seul ne suffit pas. Une mesure sans indication de \cle{fiabilité} ne vaut rien. La problématique de cette leçon est donc la suivante : \emph{comment un scientifique exprime-t-il un résultat de mesure pour qu'il soit crédible, comparable et vérifiable partout dans le monde ?} Pour y répondre, il faut d'abord comprendre ce qu'est une mesure, pourquoi elle est toujours entachée d'erreur, et comment juger la qualité de l'instrument utilisé.

\courssub{Mesure et mesurande : la valeur vraie est inaccessible}

\begin{definition}[Mesure et mesurande]
Le \cle{mesurande} est la grandeur physique que l'on cherche à déterminer (une masse, une durée, une tension, une température...). \textbf{Mesurer}, c'est comparer expérimentalement le mesurande à une \textbf{unité de référence} à l'aide d'un instrument, afin de lui attribuer une \textbf{valeur} accompagnée d'une unité :
\[
\text{grandeur} = \text{valeur numérique} \times \text{unité}, \qquad \text{ex. } m = \SI{25,2}{kg}.
\]
La \cle{valeur vraie} $x_{\text{vraie}}$ est la valeur que prendrait le mesurande si la mesure était parfaite. Elle existe, mais elle est \textbf{inaccessible} : toute mesure réelle n'en donne qu'une estimation.
\end{definition}

Pourquoi la valeur vraie reste-t-elle hors d'atteinte ? Parce qu'aucune chaîne de mesure n'est parfaite : l'instrument a une résolution limitée, l'opérateur a un temps de réaction, la température ambiante fait varier les composants, la méthode elle-même perturbe parfois ce qu'elle mesure (un voltmètre branché sur un circuit prélève un courant, aussi faible soit-il). Le physicien ne cherche donc pas « la » valeur, mais un \textbf{encadrement raisonné} de celle-ci.

\begin{savaistu}[Le kilogramme, de Sèvres à la constante de Planck]
Jusqu'en mai 2019, le kilogramme était défini par la masse d'un cylindre en platine-iridié conservé sous trois cloches de verre au Bureau international des poids et mesures, à Sèvres (France). Problème : ce prototype perdait ou gagnait quelques microgrammes au fil des décennies, faisant dériver silencieusement toutes les balances du monde ! Depuis le 20 mai 2019, le kilogramme est défini en fixant la valeur de la \textbf{constante de Planck} $h = \SI{6,62607015e-34}{J.s}$, constante fondamentale immuable de la nature (rappel : $\SI{1}{J.s} = \SI{1}{kg.m^2.s^{-1}}$). La valeur de l'unité n'est ainsi plus un objet, mais une loi physique : un progrès immense pour la fiabilité des mesures, y compris dans les laboratoires camerounais qui étalonnent les balances commerciales.
\end{savaistu}

\courssub{L'erreur sur la mesure : systématique ou aléatoire}

\begin{definition}[Erreur sur la mesure]
L'\cle{erreur} sur la mesure, notée $\varepsilon$, est l'écart entre la valeur mesurée $x_{\text{mes}}$ et la valeur vraie $x_{\text{vraie}}$ :
\[
\varepsilon = x_{\text{mes}} - x_{\text{vraie}}.
\]
Comme $x_{\text{vraie}}$ est inaccessible, l'erreur est \textbf{inconnue} : on ne peut pas la calculer exactement, on peut seulement l'\emph{estimer} ou la \emph{borner}. On distingue deux familles d'erreurs :
\begin{itemize}
\item les \cle{erreurs systématiques} : elles se reproduisent \textbf{identiquement} à chaque mesure, toujours dans le même sens (balance mal tarée qui ajoute \SI{0,4}{kg} à chaque pesée, zéro décalé d'un voltmètre, règle dilatée par la chaleur). On les combat en \textbf{identifiant et corrigeant leur cause} (étalonnage, réglage du zéro) ;
\item les \cle{erreurs aléatoires} : elles varient de façon imprévisible d'une mesure à l'autre, tantôt en excès, tantôt en défaut (temps de réaction de l'opérateur, fluctuations de tension du réseau ENEO, vibrations du sol). On ne peut pas les supprimer, mais on en \textbf{réduit l'effet en répétant la mesure} et en faisant la moyenne.
\end{itemize}
\end{definition}

L'idée intuitive est simple : une erreur systématique est un \emph{biais} (la balance « ment » toujours de la même façon), tandis qu'une erreur aléatoire est du \emph{bruit} (la mesure « tremble » autour de la valeur). Répéter cent fois une pesée sur une balance mal étalonnée ne corrigera jamais le biais : la moyenne ne chasse que le bruit.

\courssub{De l'erreur à l'incertitude : encadrer le résultat}

Puisque l'erreur est inconnue, le scientifique raisonne autrement : il déclare un \textbf{intervalle} dans lequel la valeur vraie se trouve avec une grande probabilité.

\begin{definition}[Incertitude sur la mesure]
L'\cle{incertitude} sur la mesure, notée $\Delta x$ (ou $U$), est un paramètre, associé au résultat, qui \textbf{caractérise la dispersion des valeurs que l'on peut raisonnablement attribuer au mesurande}. Le résultat s'écrit :
\[
x = x_{\text{mes}} \pm \Delta x \quad \Longleftrightarrow \quad x_{\text{vraie}} \in [\,x_{\text{mes}} - \Delta x \;;\; x_{\text{mes}} + \Delta x\,].
\]
L'\cle{incertitude relative} $\dfrac{\Delta x}{x_{\text{mes}}}$, souvent exprimée en \%, compare l'incertitude à la valeur mesurée : c'est elle qui juge la qualité de la mesure.
\end{definition}

\begin{attention}[Ne pas confondre erreur et incertitude]
C'est \textbf{la} confusion la plus fréquente au Baccalauréat. Retiens bien :
\begin{itemize}
\item l'\textbf{erreur} $\varepsilon = x_{\text{mes}} - x_{\text{vraie}}$ est un écart \emph{inconnu} : sans valeur de référence, elle \textbf{ne se calcule pas}. Dire « j'ai fait une erreur de \SI{0,3}{kg} » n'a de sens que si l'on connaît déjà une valeur de référence (étalon) ;
\item l'\textbf{incertitude} $\Delta x$ est une \emph{estimation raisonnée} de l'ensemble des erreurs plausibles : elle \textbf{se détermine} à partir de l'instrument, de la méthode et de la dispersion des mesures.
\end{itemize}
On ne dit donc jamais « l'erreur est de \SI{0,3}{kg} », mais « le résultat est $m = (25,2 \pm 0,3)\,\text{kg}$ ». Autre piège : une incertitude n'est pas une faute ! Mesurer avec incertitude est la démarche scientifique normale ; cacher l'incertitude, voilà la faute.
\end{attention}

\courssub{Les sources d'incertitude}

D'où vient cette dispersion ? Quatre familles de causes, à citer systématiquement dans une discussion de TP :

\begin{itemize}
\item \textbf{l'opérateur} : temps de réaction au chronométrage ($\approx \SI{0,2}{s}$), erreur de parallaxe lors de la lecture d'une aiguille, appréciation du maximum d'une graduation ;
\item \textbf{l'instrument} : résolution limitée (affichage au dixième, graduation au millimètre), dérive du zéro, vieillissement des composants ;
\item \textbf{la méthode} : l'instrument perturbe le mesurande (voltmètre non idéal, thermomètre qui refroidit le bain), hypothèses simplificatrices (frottements négligés) ;
\item \textbf{l'environnement} : température, humidité, vibrations, fluctuations de la tension secteur, champ magnétique parasite.
\end{itemize}

Une bonne analyse d'incertitude consiste à identifier la source \emph{dominante} : c'est elle qui fixe l'ordre de grandeur de $\Delta x$.

\courssub{Les qualités d'un instrument de mesure}

Pour choisir et juger un instrument, le métrologue utilise quatre qualités distinctes. Ne les confonds jamais : un instrument peut en posséder certaines et pas d'autres.

\begin{definition}[Sensibilité, précision, fidélité, justesse]
\begin{itemize}
\item La \cle{sensibilité} est la plus petite variation du mesurande que l'instrument peut détecter. Elle est liée à la résolution : un chronomètre au $1/100$ de seconde est plus sensible qu'un chronomètre au dixième.
\item La \cle{fidélité} (ou répétabilité) est l'aptitude de l'instrument à donner des indications \textbf{très proches les unes des autres} lors de mesures répétées du même mesurande, dans les mêmes conditions. Un instrument fidèle donne peu de dispersion.
\item La \cle{justesse} est l'aptitude de l'instrument à donner des indications \textbf{proches de la valeur vraie}, c'est-à-dire exemptes d'erreur systématique.
\item La \cle{précision} est la qualité globale d'un instrument \textbf{à la fois juste et fidèle} : indications groupées \emph{et} centrées sur la valeur vraie.
\end{itemize}
\end{definition}

L'analogie du \textbf{tir à la cible} rend ces notions inoubliables : chaque tir est une mesure, le centre de la cible est la valeur vraie. Des impacts groupés traduisent la fidélité ; des impacts centrés traduisent la justesse. Un tireur peut être fidèle sans être juste (impacts groupés mais décalés : erreur systématique, comme une lunette mal réglée) ou juste sans être fidèle (impacts dispersés mais centrés en moyenne).

\begin{popfigure}
\begin{tikzpicture}[scale=0.62]
% ---- Cible 1 : juste et fidèle ----
\begin{scope}[shift={(0,0)}]
\fill[popblueL] (0,0) circle (1.5);
\fill[white] (0,0) circle (1.0);
\fill[popblue!50] (0,0) circle (0.5);
\fill[popink] (0,0) circle (0.12);
\foreach \p in {(0.15,0.2),(-0.2,0.1),(0.05,-0.18),(-0.1,-0.05),(0.22,-0.08)}
  \fill[popdark] \p circle (0.07);
\node[below, font=\small\bfseries] at (0,-1.7) {juste et fidèle};
\node[below, font=\footnotesize\itshape] at (0,-2.15) {(précis)};
\end{scope}
% ---- Cible 2 : juste, non fidèle ----
\begin{scope}[shift={(4.2,0)}]
\fill[popblueL] (0,0) circle (1.5);
\fill[white] (0,0) circle (1.0);
\fill[popblue!50] (0,0) circle (0.5);
\fill[popink] (0,0) circle (0.12);
\foreach \p in {(0.9,0.6),(-1.0,0.5),(0.4,-1.0),(-0.6,-0.8),(0.1,1.1)}
  \fill[popdark] \p circle (0.07);
\node[below, font=\small\bfseries] at (0,-1.7) {juste, non fidèle};
\node[below, font=\footnotesize\itshape] at (0,-2.15) {(dispersé mais centré)};
\end{scope}
% ---- Cible 3 : non juste, fidèle ----
\begin{scope}[shift={(8.4,0)}]
\fill[popblueL] (0,0) circle (1.5);
\fill[white] (0,0) circle (1.0);
\fill[popblue!50] (0,0) circle (0.5);
\fill[popink] (0,0) circle (0.12);
\foreach \p in {(0.95,0.75),(1.1,0.55),(0.8,0.9),(1.05,0.9),(0.85,0.55)}
  \fill[poporange] \p circle (0.07);
\node[below, font=\small\bfseries] at (0,-1.7) {non juste, fidèle};
\node[below, font=\footnotesize\itshape] at (0,-2.15) {(erreur systématique)};
\end{scope}
% ---- Cible 4 : non juste, non fidèle ----
\begin{scope}[shift={(12.6,0)}]
\fill[popblueL] (0,0) circle (1.5);
\fill[white] (0,0) circle (1.0);
\fill[popblue!50] (0,0) circle (0.5);
\fill[popink] (0,0) circle (0.12);
\foreach \p in {(1.2,-0.9),(0.6,1.2),(-1.1,0.9),(1.0,0.2),(-0.4,-1.2)}
  \fill[poporange] \p circle (0.07);
\node[below, font=\small\bfseries] at (0,-1.7) {non juste, non fidèle};
\node[below, font=\footnotesize\itshape] at (0,-2.15) {(mauvais instrument)};
\end{scope}
\end{tikzpicture}
\legende{L'analogie du tir à la cible : le centre noir est la valeur vraie, chaque impact est une mesure. La fidélité se lit sur le groupement des impacts, la justesse sur leur position par rapport au centre.}
\end{popfigure}

\begin{exemplebox}[Chronométrage au jugé sur un 100 m]
Au championnat scolaire de Douala, un professeur chronomètre au jugé un élève sur \SI{100}{m} avec un chronomètre affichant le centième de seconde. Il répète l'expérience sur dix courses du même athlète et obtient des temps très groupés : \SI{11,42}{s}, \SI{11,45}{s}, \SI{11,40}{s}, \SI{11,43}{s}... La \textbf{fidélité} de l'ensemble est excellente (dispersion de quelques centièmes). Pourtant, le temps de réaction de l'opérateur, de l'ordre de \SI{0,2}{s} au départ comme à l'arrivée, introduit une erreur bien supérieure au centième affiché : la \textbf{justesse} est limitée par l'opérateur, pas par l'instrument. Afficher « \SI{11,42}{s} » est donc trompeur : le résultat honnête s'écrit plutôt $t = (11,4 \pm 0,2)\,\text{s}$, soit une incertitude relative de $\dfrac{0,2}{11,4} \approx \num{0,018}$, c'est-à-dire environ \SI{1,8}{\%}. Moralité : la résolution d'un appareil ne garantit jamais, à elle seule, la justesse de la mesure.
\end{exemplebox}

\begin{popfigure}
\begin{tikzpicture}[scale=1.0]
% Cadran du voltmètre
\draw[thick, popdark] (-3.4,-0.4) rectangle (3.4,3.1);
\draw[popdark] (0,0) circle (0.12);
\fill[popdark] (0,0) circle (0.08);
% Graduations : échelle 0 à 15 V, de -60° à +60°
\foreach \v in {0,1,...,15}{
  \pgfmathsetmacro{\ang}{-60 + 8*\v}
  \draw[popdark] (\ang:2.3) -- (\ang:2.55);
}
\foreach \v/\lab in {0/0, 5/5, 10/10, 15/15}{
  \pgfmathsetmacro{\ang}{-60 + 8*\v}
  \node[font=\footnotesize] at (\ang:2.9) {\lab};
}
\node[font=\footnotesize\bfseries] at (0,2.0) {V};
% Zone d'incertitude de lecture autour de 9 V (entre 8.5 et 9.5 V)
\pgfmathsetmacro{\aone}{-60 + 8*8.5}
\pgfmathsetmacro{\atwo}{-60 + 8*9.5}
\fill[poporange!45] (\aone:1.55) arc (\aone:\atwo:1.55) -- (\atwo:2.25) arc (\atwo:\aone:2.25) -- cycle;
% Aiguille sur 9 V
\pgfmathsetmacro{\aaig}{-60 + 8*9}
\draw[very thick, poppink] (0,0) -- (\aaig:2.2);
% Annotations
\draw[-{Stealth}, poporange, thick] (2.9,1.15) -- (2.35,1.35);
\node[right, font=\footnotesize, text=poporange, align=left] at (2.95,1.1) {zone\\ d'incertitude\\ $\pm 0{,}5$ V};
\node[below, font=\footnotesize] at (0,-0.15) {lecture : $U = (9{,}0 \pm 0{,}5)$ V};
\end{tikzpicture}
\legende{Voltmètre à aiguille (calibre \SI{15}{V}). L'aiguille indique \SI{9}{V}, mais la lecture au jugé entre deux graduations impose une incertitude de lecture de l'ordre d'une demi-graduation, ici $\pm \SI{0,5}{V}$, matérialisée par la zone orange.}
\end{popfigure}

\begin{exoresolu}[Pesées au marché Mokolo]
Deux commerçants pèsent le même sac de riz. La balance A, au gramme près mais jamais étalonnée, affiche \SI{25,2}{kg} de façon très stable. La balance B, récemment étalonnée, affiche successivement \SI{24,6}{kg}, \SI{25,0}{kg}, \SI{24,8}{kg}, \SI{24,9}{kg}. (1) Comparer la fidélité et la justesse des deux balances. (2) Exprimer le résultat de la balance B avec une incertitude de lecture de $\pm \SI{0,2}{kg}$ et calculer l'incertitude relative.
\tcblower
(1) La balance A donne des indications très stables : elle est \textbf{fidèle}. Mais n'ayant jamais été étalonnée, elle peut souffrir d'une erreur systématique : sa \textbf{justesse} est douteuse. La balance B donne des indications dispersées (moins fidèle), mais étalonnée, sa moyenne
\[
\bar{m} = \frac{24,6 + 25,0 + 24,8 + 24,9}{4} = \SI{24,8}{kg}
\]
est vraisemblablement proche de la valeur vraie : elle est plus \textbf{juste}. (2) Résultat : $m = (24,8 \pm 0,2)\,\text{kg}$. Incertitude relative : $\dfrac{\Delta m}{m} = \dfrac{0,2}{24,8} \approx \num{0,008}$, soit environ \SI{0,8}{\%} : une mesure de bonne qualité pour une pesée commerciale.
\end{exoresolu}

\begin{aretenir}[L'essentiel de la section]
\begin{itemize}
\item Mesurer, c'est attribuer une valeur et une unité à un mesurande ; la \textbf{valeur vraie est inaccessible}.
\item L'\textbf{erreur} $\varepsilon = x_{\text{mes}} - x_{\text{vraie}}$ est inconnue : systématique (biais reproductible, à corriger) ou aléatoire (bruit, réduit par la moyenne).
\item L'\textbf{incertitude} $\Delta x$ s'estime : le résultat s'écrit $x = x_{\text{mes}} \pm \Delta x$, avec une incertitude relative $\Delta x / x_{\text{mes}}$ qui juge la qualité.
\item Quatre qualités d'instrument : \textbf{sensibilité} (plus petite variation détectable), \textbf{fidélité} (indications groupées), \textbf{justesse} (proximité de la valeur vraie), \textbf{précision} (juste \emph{et} fidèle) — à visualiser avec le tir à la cible.
\end{itemize}
\end{aretenir}

\begin{exercice}[Balance de laboratoire]
Un élève pèse cinq fois la même masse marquée de \SI{100,0}{g} sur une balance de laboratoire : \SI{100,4}{g} ; \SI{100,5}{g} ; \SI{100,3}{g} ; \SI{100,4}{g} ; \SI{100,4}{g}. (1) La balance est-elle fidèle ? Juste ? (2) Quel type d'erreur met-on en évidence, et comment y remédier ? (3) Exprimer le résultat avec une incertitude de lecture de $\pm \SI{0,1}{g}$ et donner l'incertitude relative.
\end{exercice}

\begin{corrige}[Exercice 1]
(1) Les cinq indications sont groupées dans un intervalle de \SI{0,2}{g} : la balance est \textbf{fidèle}. Mais la moyenne
\[
\bar{m} = \frac{100,4 + 100,5 + 100,3 + 100,4 + 100,4}{5} = \SI{100,4}{g}
\]
s'écarte systématiquement de la masse marquée \SI{100,0}{g} : la balance n'est \textbf{pas juste}. (2) Il s'agit d'une \textbf{erreur systématique} d'environ $+\SI{0,4}{g}$ (zéro décalé ou mauvais étalonnage) ; on y remédie en réglant le zéro ou en étalonnant la balance avec une masse étalon — la répétition des mesures ne corrige rien. (3) $m = (100,4 \pm 0,1)\,\text{g}$ ; incertitude relative $\dfrac{0,1}{100,4} \approx \num{0,001}$, soit environ \SI{0,1}{\%}, ce qui traduit la bonne fidélité, sans préjuger de la justesse.
\end{corrige}

Nous savons désormais qualifier une mesure et un instrument. Reste à apprendre, dans la prochaine section, à \emph{quantifier} l'incertitude d'une mesure unique à partir des caractéristiques de l'instrument.

\coursec{Incertitude d'une mesure unique}

Dans la section précédente, nous avons vu que toute mesure est entachée d'une erreur, que cette erreur reste inaccessible, et que la qualité d'un instrument (fidélité, justesse, sensibilité, précision) conditionne la confiance que l'on peut accorder au résultat. Reste une question pratique et fondamentale : \emph{comment chiffrer cette confiance ?} Autrement dit, lorsqu'on écrit $L = \SI{23,4}{\centi\meter}$, quel intervalle autour de \num{23,4} peut honnêtement contenir la valeur vraie ? C'est toute l'affaire de l'\cle{incertitude de mesure}. Nous commençons par le cas le plus simple et le plus fréquent au laboratoire : celui où l'on n'effectue qu'\textbf{une seule mesure}.

\courssub{Pourquoi une mesure unique exige une incertitude élargie}

\begin{experience}[Une situation de paillasse]
Au TP de physique du lycée de Nkongsamba, un groupe mesure la tension aux bornes d'une lampe avec un voltmètre : l'aiguille s'immobilise entre les graduations \num{4} et \num{5} volts. Un élève lit \SI{4,5}{\volt}, sa camarade lit \SI{4,4}{\volt}, le troisième hésite entre \SI{4,5}{\volt} et \SI{4,6}{\volt}. Une seule mesure, trois lectures plausibles. Qui a raison ? Personne ne peut le savoir : la valeur vraie est inaccessible. En revanche, tous peuvent convenir d'un \emph{intervalle} dans lequel la valeur vraie se trouve presque sûrement : par exemple $[\SI{4,4}{\volt}\,;\,\SI{4,6}{\volt}]$.
\end{experience}

Lorsqu'on répète une mesure de nombreuses fois, la dispersion des résultats renseigne sur l'incertitude (ce sera l'objet de la section suivante). Mais très souvent — au Baccalauréat comme dans la vie professionnelle — on ne dispose que d'\textbf{une seule mesure}. Il faut alors estimer l'incertitude \emph{a priori}, en analysant l'instrument, sa notice et les conditions d'utilisation.

\begin{definition}[Incertitude élargie $\Delta$ pour une mesure unique]
Pour une mesure unique d'une grandeur $G$, on associe au résultat $g$ une \cle{incertitude élargie} $\Delta g$, de sorte que la valeur vraie $g_{\text{vraie}}$ appartienne, avec un niveau de confiance élevé (en pratique environ $95\ \%$), à l'intervalle :
\[
g - \Delta g \ \leq\ g_{\text{vraie}}\ \leq\ g + \Delta g
\qquad\text{que l'on écrit}\qquad
G = g \pm \Delta g.
\]
L'intervalle $[g-\Delta g\,;\,g+\Delta g]$ est l'\cle{intervalle de confiance} de la mesure.
\end{definition}

\begin{attention}[Ce que $\Delta$ n'est pas]
L'incertitude $\Delta g$ n'est \textbf{pas} l'erreur commise (celle-ci est inconnue par définition), ni une marge arbitraire « pour être sûr ». C'est une estimation raisonnée, construite à partir des caractéristiques de l'instrument et des conditions de mesure. Deux mesures de la même grandeur avec deux instruments différents n'ont pas la même incertitude.
\end{attention}

\courssub{Incertitude de lecture sur un appareil analogique}

Un appareil \cle{analogique} (ampèremètre ou voltmètre à aiguille, règle graduée, thermomètre à colonne de liquide, balance de Roberval) fournit une indication repérée sur une échelle graduée. La lecture est limitée par deux choses : la finesse des graduations et l'œil de l'opérateur qui doit apprécier la position de l'aiguille (ou du repère) \emph{entre} deux graduations.

\begin{propriete}[Règle de la demi-graduation]
Pour un appareil analogique utilisé dans de \textbf{bonnes conditions} (œil en face du repère, échelle bien lisible, aiguille stable), l'incertitude de lecture est estimée à la \textbf{demi-graduation} :
\[
\Delta_{\text{lecture}} = \frac{\text{valeur d'une graduation}}{2}.
\]
Si les conditions sont dégradées (échelle fine et dense, aiguille qui oscille, lecture rapide, éclairage médiocre), on retient prudemment une graduation entière : $\Delta_{\text{lecture}} = 1$ graduation.
\end{propriete}

L'idée intuitive est simple : si l'aiguille tombe entre deux graduations, un opérateur soigneux est capable de dire si elle est plus proche de l'une ou de l'autre, voire d'estimer le milieu. Ce qu'il ne peut pas faire, c'est garantir mieux que la moitié de l'écart entre deux graduations.

\begin{popfigure}
\begin{center}
\begin{tikzpicture}[scale=1.0]
% Cadran d'ampèremètre : arc de graduations
\draw[thick, popdark] (-4.6,-0.6) rectangle (4.6,3.4);
\node[popdark, font=\small] at (0,3.0) {\textbf{Ampèremètre — calibre \SI{1}{\ampere}}};
% Arc principal
\draw[thick, popink] (-3.6,0) arc (180:0:3.6);
% Graduations : 0 à 10, de 180° à 0°
\foreach \i in {0,...,10} {
  \pgfmathsetmacro{\ang}{180 - \i*18}
  \draw[thick, popdark] (\ang:3.6) -- (\ang:3.25);
  \pgfmathsetmacro{\lab}{\i/10}
  \node[font=\scriptsize, popdark] at (\ang:2.85) {\pgfmathprintnumber[precision=1,fixed]{\lab}};
}
% Demi-graduations
\foreach \i in {0,...,9} {
  \pgfmathsetmacro{\ang}{180 - \i*18 - 9}
  \draw[popmuted] (\ang:3.6) -- (\ang:3.35);
}
% Aiguille entre 0.6 et 0.7, plus proche de 0.65 (angle pour 0.65 : 180-6.5*18 = 63°)
\draw[very thick, poporange, -{Stealth[length=3mm]}] (0,-0.15) -- (63:3.45);
\fill[popdark] (0,-0.15) circle (0.09);
% Encadrement de la lecture
\draw[thick, poppurple, dashed] (72:3.75) arc (72:54:3.75);
\node[poppurple, font=\scriptsize, align=center] at (90:4.35) {zone d'encadrement\\ $[\SI{0,60}{\ampere}\,;\,\SI{0,70}{\ampere}]$};
\draw[poppurple, dashed, -{Stealth[length=2mm]}] (75:4.0) -- (70:3.8);
% Annotation lecture
\node[poporange, font=\small, align=center] at (0,-1.15) {Lecture : $I = \SI{0,65}{\ampere}$ \quad (graduation : \SI{0,1}{\ampere})};
\end{tikzpicture}
\end{center}
\legende{Cadran d'un ampèremètre analogique : l'aiguille tombe entre les graduations \SI{0,6}{\ampere} et \SI{0,7}{\ampere}. La lecture estimée est $I = \SI{0,65}{\ampere}$, avec $\Delta I = \dfrac{\SI{0,1}{\ampere}}{2} = \SI{0,05}{\ampere}$.}
\end{popfigure}

\begin{exemplebox}[Mesure d'une longueur à la règle graduée au millimètre]
On mesure la longueur d'un livre de physique avec une règle graduée au millimètre, en plaçant soigneusement le zéro en face d'une extrémité. L'autre extrémité tombe entre \SI{23,4}{\centi\meter} et \SI{23,5}{\centi\meter}, plus proche de \SI{23,4}{\centi\meter}. On retient :
\[
L = \SI{23,4}{\centi\meter} \pm \SI{0,05}{\centi\meter}
\qquad\text{car}\qquad
\Delta L = \frac{\SI{1}{\milli\meter}}{2} = \SI{0,5}{\milli\meter} = \SI{0,05}{\centi\meter}.
\]
La valeur vraie de la longueur se trouve donc, avec une confiance d'environ $95\ \%$, dans l'intervalle $[\SI{23,35}{\centi\meter}\,;\,\SI{23,45}{\centi\meter}]$.
\end{exemplebox}

\begin{attention}[Erreur fréquente : la graduation entière]
Beaucoup d'élèves écrivent systématiquement $\Delta = 1$ graduation pour tout appareil analogique. C'est \textbf{faux} pour un appareil bien utilisé : la règle de référence est la \textbf{demi-graduation}. Ne retenir une graduation entière que si l'énoncé signale des conditions de lecture difficiles (aiguille instable, échelle peu lisible). À l'inverse, ne jamais écrire $\Delta = 0$ sous prétexte que « l'aiguille tombe juste sur la graduation » : même dans ce cas, l'œil ne peut garantir la coïncidence à mieux qu'une demi-graduation.
\end{attention}

\courssub{Incertitude de lecture sur un appareil numérique}

Un appareil \cle{numérique} (multimètre digital, balance électronique, chronomètre digital, thermomètre à affichage) affiche directement une valeur avec un nombre fini de chiffres. La tentation est grande de croire que cette valeur est « exacte » : c'est une illusion. L'affichage est \emph{quantifié} : l'appareil arrondit la valeur mesurée au dernier chiffre affiché.

\begin{propriete}[Règle du dernier chiffre affiché]
En l'absence d'information supplémentaire, l'incertitude de lecture sur un appareil numérique est estimée à \textbf{une unité du dernier chiffre affiché} (appelée aussi \cle{digit} ou échelon de quantification) :
\[
\Delta_{\text{lecture}} = 1\ \text{unité du dernier digit affiché}.
\]
Exemple : un affichage $m = \SI{152,37}{\gram}$ correspond à $\Delta m = \SI{0,01}{\gram}$.
\end{propriete}

En effet, tout affichage \num{152,37} peut provenir d'une valeur réelle comprise entre \num{152,365} et \num{152,375} (arrondi au centième) ; l'ambiguïté couvre donc un échelon complet du dernier digit. Cette règle est un \emph{minimum} : la notice du constructeur donne souvent une incertitude plus grande, qu'il faut alors préférer.

\courssub{Incertitude constructeur : lire et interpréter une notice}

Le fabricant d'un instrument de mesure garantit, pour chaque calibre, une \cle{incertitude constructeur} (ou « précision » au sens commercial du terme). Elle est donnée dans la notice sous une forme normalisée, typiquement :
\[
\Delta_{\text{constructeur}} = \pm\bigl(p\ \%\ \text{de la lecture} + n\ \text{digits}\bigr),
\]
où $p\ \%$ est une part proportionnelle à la valeur lue et $n$ digits une part fixe liée à la résolution de l'affichage.

\begin{methode}[Exploiter une notice constructeur]
\begin{enumerate}
\item Repérer, dans la notice, la ligne correspondant à la \textbf{fonction} (V, A, $\Omega$\ldots) et au \textbf{calibre} utilisés.
\item Identifier le pourcentage $p$ et le nombre de digits $n$.
\item Calculer la part proportionnelle : $\dfrac{p}{100}\times(\text{valeur lue})$.
\item Calculer la part fixe : $n \times (\text{valeur d'un digit})$, le digit étant l'unité du dernier chiffre affiché.
\item Additionner les deux contributions : $\Delta = \dfrac{p}{100}\times\text{lecture} + n\times\text{digit}$.
\item Arrondir $\Delta$ à un chiffre significatif (par excès), puis ajuster l'écriture du résultat.
\end{enumerate}
\end{methode}

\begin{exemplebox}[Tension mesurée au multimètre numérique]
Un élève mesure une tension continue : l'afficheur indique $U = \SI{4,52}{\volt}$. La notice du multimètre indique, pour ce calibre : $\pm(\SI{0,5}{\percent}\ \text{de la lecture} + 1\ \text{digit})$. Ici, 1 digit $= \SI{0,01}{\volt}$ (dernier chiffre affiché au centième de volt).
\begin{align*}
\Delta U &= \frac{0,5}{100}\times \SI{4,52}{\volt} + 1\times \SI{0,01}{\volt} \\
&= \SI{0,0226}{\volt} + \SI{0,01}{\volt} \\
&= \SI{0,0326}{\volt} \approx \SI{0,03}{\volt}.
\end{align*}
Le résultat s'exprime donc : $U = \SI{4,52}{\volt} \pm \SI{0,03}{\volt}$, c'est-à-dire $U \in [\SI{4,49}{\volt}\,;\,\SI{4,55}{\volt}]$.
\end{exemplebox}

\begin{attention}[Pièges de la notice]
\begin{itemize}
\item Ne pas oublier la part « digits » : elle est souvent dominante pour les petites lectures.
\item « $\SI{0,5}{\percent}$ » signifie $\dfrac{0,5}{100}$, pas $0,5$ !
\item Le digit dépend du calibre : sur un calibre où l'affichage donne \SI{4,5}{\volt}, un digit vaut \SI{0,1}{\volt}, et non \SI{0,01}{\volt}. Toujours regarder la position du dernier chiffre affiché.
\end{itemize}
\end{attention}

\courssub{La part de l'opérateur : temps de réaction et parallaxe}

L'instrument n'est pas la seule source d'incertitude : l'être humain qui l'utilise en introduit aussi. Deux contributions classiques doivent être connues.

\textbf{Le temps de réaction.} Lorsqu'on déclenche et arrête un chronomètre manuel (mesure de la période d'un pendule, chronométrage d'un sprint au stade de Bafoussam), l'opérateur réagit avec un retard, tant au départ qu'à l'arrêt. Le temps de réaction humain est de l'ordre de \SI{0,1}{\second} à \SI{0,2}{\second}. On retient en général :
\[
\Delta t_{\text{opérateur}} \approx \SI{0,1}{\second} \ \text{à}\ \SI{0,2}{\second},
\]
soit souvent $\Delta t \approx \SI{0,2}{\second}$ pour une mesure manuelle complète (départ + arrêt). C'est pourquoi, pour mesurer la période $T$ d'un pendule, on chronomètre \textbf{plusieurs périodes} (par exemple $10\,T$) : l'incertitude de \SI{0,2}{\second} se répartit alors sur dix périodes, et $\Delta T \approx \SI{0,02}{\second}$.

\textbf{L'erreur de parallaxe.} Lire un cadran à aiguille ou un thermomètre de biais décale la position apparente du repère par rapport à l'échelle : c'est l'\cle{erreur de parallaxe}. Elle peut atteindre une graduation entière. Les bons voltmètres de laboratoire comportent d'ailleurs un \emph{miroir} derrière l'aiguille : la lecture est correcte lorsque l'aiguille et son image se confondent. La règle pratique : placer l'œil \textbf{en face} du repère, à la perpendiculaire du cadran.

\begin{exemplebox}[Chronométrage d'une chute libre]
Pour mesurer la durée de chute d'une bille lâchée depuis une hauteur $h = \SI{1,80}{\meter}$, un élève utilise un chronomètre au centième de seconde et obtient $t = \SI{0,61}{\second}$. L'incertitude d'affichage (\SI{0,01}{\second}) est négligeable devant le temps de réaction : on retient $\Delta t = \SI{0,2}{\second}$. Le résultat est donc $t = \SI{0,6}{\second} \pm \SI{0,2}{\second}$ : c'est l'\textbf{opérateur}, et non l'instrument, qui limite la qualité de la mesure. D'où l'intérêt des capteurs photoélectriques dans les laboratoires modernes.
\end{exemplebox}

\courssub{Choisir la source dominante et exprimer l'incertitude totale}

En pratique, plusieurs contributions coexistent : lecture, constructeur, opérateur, conditions expérimentales. Faut-il toutes les additionner ? Non : si l'une d'elles est nettement plus grande que les autres, elle « écrase » les contributions secondaires.

\begin{methode}[Estimer $\Delta$ pour une mesure unique]
\begin{enumerate}
\item \textbf{Identifier} toutes les sources d'incertitude : lecture (analogique ou numérique), notice constructeur, opérateur (réaction, parallaxe), réglage du zéro, conditions (température, stabilité).
\item \textbf{Chiffrer} chaque contribution simple $\Delta_1, \Delta_2, \ldots$
\item \textbf{Comparer} :
\begin{itemize}
\item si une contribution domine nettement (au moins $3$ à $4$ fois plus grande que les autres), retenir $\Delta = \max(\Delta_i)$ ;
\item si plusieurs contributions sont comparables, les combiner quadratiquement :
\[
\Delta = \sqrt{\Delta_1^{2} + \Delta_2^{2} + \cdots}
\]
(au lycée, on se contente souvent de la somme $\Delta_1 + \Delta_2$, qui majore prudemment l'incertitude).
\end{itemize}
\item \textbf{Arrondir} $\Delta$ à un seul chiffre significatif, par excès, puis exprimer le résultat sous la forme $G = g \pm \Delta g$ avec le même nombre de décimales pour $g$ et $\Delta g$.
\end{enumerate}
\end{methode}

\begin{center}
\begin{tabularx}{\linewidth}{|l|X|X|}
\hline
\rowcolor{popblueL}
\textbf{Critère} & \textbf{Appareil analogique} & \textbf{Appareil numérique} \\
\hline
Indication & Aiguille ou repère sur une échelle graduée & Affichage direct de chiffres \\
\hline
Règle de base & $\Delta = \dfrac{\text{graduation}}{2}$ & $\Delta = 1$ digit (dernier chiffre affiché) \\
\hline
Conditions dégradées & $\Delta = 1$ graduation (aiguille instable, parallaxe) & Notice constructeur : $\pm(p\,\%\,\text{lecture} + n\ \text{digits})$ \\
\hline
Intervention de l'opérateur & Forte (estimation entre graduations, parallaxe) & Faible (mais temps de réaction si chronométrage manuel) \\
\hline
Exemple & Règle au mm : $L = \SI{23,4}{\centi\meter} \pm \SI{0,05}{\centi\meter}$ & Multimètre : $U = \SI{4,52}{\volt} \pm \SI{0,03}{\volt}$ \\
\hline
\end{tabularx}
\end{center}

\begin{exoresolu}[Période d'un pendule simple]
Au laboratoire, on mesure la durée de $10$ oscillations complètes d'un pendule simple avec un chronomètre manuel affichant le centième de seconde. Lecture : $10\,T = \SI{18,42}{\second}$. Le temps de réaction de l'opérateur est estimé à \SI{0,2}{\second} sur la mesure complète.
\begin{enumerate}
\item Déterminer la période $T$ et son incertitude $\Delta T$.
\item Exprimer le résultat sous la forme $T \pm \Delta T$.
\end{enumerate}
\tcblower
\textbf{1.} La période vaut :
\[
T = \frac{\SI{18,42}{\second}}{10} = \SI{1,842}{\second}.
\]
Deux contributions à l'incertitude sur la durée totale : l'affichage (\SI{0,01}{\second}) et le temps de réaction (\SI{0,2}{\second}). Le temps de réaction domine nettement ($0,2 \gg 0,01$) : on retient $\Delta(10T) = \SI{0,2}{\second}$. En divisant par $10$ :
\[
\Delta T = \frac{\SI{0,2}{\second}}{10} = \SI{0,02}{\second}.
\]
\textbf{2.} Résultat final :
\[
T = \SI{1,84}{\second} \pm \SI{0,02}{\second},
\qquad\text{soit}\qquad
T \in [\SI{1,82}{\second}\,;\,\SI{1,86}{\second}].
\]
On remarque que mesurer $10$ périodes a divisé l'incertitude par $10$ : c'est une stratégie expérimentale à retenir pour toute grandeur périodique.
\end{exoresolu}

\begin{aretenir}[L'essentiel de la section]
\begin{itemize}
\item Une mesure unique s'accompagne d'une \cle{incertitude élargie} $\Delta$ définissant un intervalle de confiance : $G = g \pm \Delta g$.
\item Appareil \textbf{analogique} bien utilisé : $\Delta = \dfrac{\text{graduation}}{2}$ ; une graduation entière si les conditions sont mauvaises.
\item Appareil \textbf{numérique} : $\Delta = 1$ digit au minimum ; la \textbf{notice constructeur} $\pm(p\,\%\,\text{lecture} + n\ \text{digits})$ prime si elle est connue.
\item L'\textbf{opérateur} compte : temps de réaction $\approx \SI{0,1}{\second}$ à \SI{0,2}{\second}, erreur de parallaxe évitée en se plaçant face au cadran.
\item On retient la contribution \textbf{dominante}, ou on combine les contributions comparables ; $\Delta$ s'arrondit à un chiffre significatif, par excès.
\end{itemize}
\end{aretenir}

\begin{savaistu}[Et au quotidien au Cameroun ?]
Lorsqu'un technicien d'ENEO relève la tension du réseau (\SI{220}{\volt} nominal), son multimètre n'affiche jamais exactement \SI{220,00}{\volt} : la valeur fluctue et l'appareil lui-même a sa propre incertitude, typiquement $\pm(\SI{1}{\percent} + 2\ \text{digits})$, soit environ $\pm \SI{2}{\volt}$. De même, la balance électronique du marché de Mokolo, certifiée par les services de métrologie du MINMIDT, garantit une incertitude maximale : c'est elle qui protège le client comme le commerçant. La métrologie — la science de la mesure — est partout, silencieuse mais indispensable.
\end{savaistu}

La section suivante montrera comment la \textbf{répétition} des mesures permet d'aller plus loin : moyenne, écart-type et incertitude-type donneront à l'estimation de l'incertitude une base statistique rigoureuse.

\coursec{Série de mesures : moyenne, écart-type, incertitude-type}

Dans la section précédente, nous avons vu comment estimer l'incertitude d'une \cle{mesure unique} : on lit l'appareil une fois, on tient compte de sa résolution ou de sa classe, et l'on écrit le résultat sous la forme $x \pm \Delta x$. Cette démarche est honnête mais souvent pessimiste : une seule mesure peut être affectée par un coup de malchance (réflexe un peu lent au chronomètre, lecture légèrement de biais, vibration de la table au moment de la mesure). Le physicien a une arme redoutable contre ces fluctuations imprévisibles : \textbf{répéter la mesure} et traiter statistiquement la série obtenue. C'est toute l'ambition de cette section.

\courssub{Pourquoi répéter une mesure ?}

\begin{experience}[Le pendule du laboratoire]
Un groupe d'élèves de Terminale C du lycée de Nkolbisson mesure la période $T$ d'un pendule simple (fil de longueur $\ell \approx 1\,\text{m}$, petite sphère métallique). Chaque élève chronomètre une oscillation complète, dans les mêmes conditions. Voici les 20 valeurs relevées (en secondes) :

\begin{center}
\begin{tabular}{|c|cccccccccc|}
\hline
Mesure & 1 & 2 & 3 & 4 & 5 & 6 & 7 & 8 & 9 & 10 \\
\hline
$T$ (s) & 1,97 & 2,01 & 1,99 & 1,95 & 2,00 & 1,98 & 2,03 & 1,96 & 1,99 & 1,97 \\
\hline
Mesure & 11 & 12 & 13 & 14 & 15 & 16 & 17 & 18 & 19 & 20 \\
\hline
$T$ (s) & 2,00 & 1,98 & 1,96 & 2,02 & 1,99 & 1,97 & 2,01 & 1,98 & 1,95 & 2,00 \\
\hline
\end{tabular}
\end{center}

\textbf{Observations :} aucune valeur ne se répète exactement à l'identique ; les résultats s'étalent entre \SI{1,95}{s} et \SI{2,03}{s}, sans tendance systématique (les valeurs trop grandes et trop petites se compensent).

\textbf{Interprétation :} les écarts proviennent de causes \emph{aléatoires} : temps de réaction de l'expérimentateur au départ et à l'arrêt du chronomètre (de l'ordre de \SI{0,1}{s} partagé entre les deux événements), amplitude pas rigoureusement identique, courants d'air. Ces causes changent de signe d'une mesure à l'autre : c'est précisément ce qui rend leur \emph{compensation} possible par moyennage.
\end{experience}

L'idée fondamentale est la suivante : les erreurs \cle{aléatoires} sont tantôt positives, tantôt négatives. En additionnant un grand nombre de mesures, elles se compensent partiellement, alors que la « vraie » valeur, elle, s'additionne toujours dans le même sens. Répéter la mesure ne supprime pas le bruit, mais il le \textbf{dilue}.

\begin{attention}[Ce que la répétition ne peut pas faire]
La répétition ne réduit que les erreurs \textbf{aléatoires}. Une erreur \textbf{systématique} (chronomètre qui retarde, règle mal étalonnée, zéro décalé) affecte toutes les mesures de la même façon : faire la moyenne de 100 mesures biaisées donne une moyenne biaisée ! Contre les erreurs systématiques, seuls l'étalonnage de l'instrument et le changement de méthode sont efficaces.
\end{attention}

\courssub{La moyenne arithmétique : meilleure estimation du mesurande}

Intuitivement, si les écarts sont symétriques autour de la valeur vraie, le « centre de gravité » des mesures est le meilleur candidat pour l'estimer. Ce centre, c'est la moyenne arithmétique.

\begin{definition}[Moyenne d'une série de mesures]
Pour une série de $n$ mesures indépendantes $x_1, x_2, \ldots, x_n$ d'une même grandeur (le \cle{mesurande}), effectuées dans les mêmes conditions, la \textbf{moyenne arithmétique} est :
\[
\bar{x} = \frac{1}{n}\sum_{i=1}^{n} x_i = \frac{x_1 + x_2 + \cdots + x_n}{n}.
\]
$\bar{x}$ constitue la \textbf{meilleure estimation} de la valeur du mesurande.
\end{definition}

Pourquoi la moyenne plutôt que la médiane ou une mesure « qui semble bonne » ? Parce que la moyenne possède une propriété remarquable : la somme des écarts $\sum_i (x_i - \bar{x})$ est \emph{exactement nulle}, et la somme des carrés des écarts $\sum_i (x_i - \bar{x})^2$ est \emph{minimale} lorsqu'on la calcule par rapport à $\bar{x}$ plutôt que par rapport à n'importe quelle autre valeur. La moyenne est donc, au sens des moindres carrés, le point d'équilibre naturel de la série.

\begin{exemplebox}[Calcul de la moyenne des 20 périodes]
En additionnant les 20 valeurs du tableau ci-dessus, on obtient $\sum_i T_i = \SI{39,71}{s}$, d'où :
\[
\bar{T} = \frac{39,71}{20} = \SI{1,9855}{s} \approx \SI{1,986}{s}.
\]
On garde pour l'instant un chiffre de plus que les données ; l'arrondi final sera dicté par l'incertitude (section 4).
\end{exemplebox}

\courssub{L'écart-type expérimental : mesurer la dispersion}

La moyenne seule ne dit pas si les mesures sont resserrées ou très étalées. Deux séries peuvent avoir la même moyenne et des qualités très différentes. Il faut quantifier la \cle{dispersion}.

L'idée : pour chaque mesure, on calcule l'écart $x_i - \bar{x}$. On ne peut pas simplement en faire la moyenne (elle vaut zéro !). On élève donc les écarts au carré — ce qui les rend tous positifs et punit davantage les grands écarts — puis on fait la moyenne de ces carrés, et l'on revient à l'unité d'origine par une racine carrée.

\begin{definition}[Écart-type expérimental]
L'\textbf{écart-type expérimental} d'une série de $n$ mesures est :
\[
s = \sqrt{\dfrac{\displaystyle\sum_{i=1}^{n}\left(x_i - \bar{x}\right)^2}{n-1}}.
\]
Il s'exprime dans la même unité que les mesures et quantifie la dispersion de la série autour de la moyenne : plus $s$ est petit, plus les mesures sont resserrées.
\end{definition}

\begin{attention}[Pourquoi $n-1$ et pas $n$ ?]
C'est \textbf{l'erreur la plus fréquente} au Bac et en TP. On divise par $n-1$ (et non par $n$) parce que les $n$ écarts $x_i - \bar{x}$ ne sont pas indépendants : leur somme est nulle, donc connaître $n-1$ d'entre eux impose le dernier. On dit qu'il n'y a que $n-1$ \emph{degrés de liberté}. Diviser par $n$ sous-estimerait systématiquement la dispersion. Astuce : sur ta calculatrice, choisis bien l'écart-type noté $s$, $\sigma_{n-1}$ ou « écart-type échantillon », et \emph{pas} $\sigma_n$ (écart-type d'une population complète).
\end{attention}

\begin{exemplebox}[Écart-type des 20 périodes]
On calcule les écarts $T_i - \bar{T}$, on les élève au carré, on somme : $\sum_i (T_i - \bar{T})^2 \approx \SI{8,3e-3}{s^2}$. D'où :
\[
s = \sqrt{\frac{8,3\times 10^{-3}}{20-1}} = \sqrt{4,4\times 10^{-4}} \approx \SI{0,021}{s}.
\]
La dispersion typique d'une mesure individuelle est donc d'environ \SI{0,02}{s}, cohérent avec le temps de réaction humain au chronomètre.
\end{exemplebox}

\begin{popfigure}
\begin{center}
\begin{tikzpicture}
\begin{axis}[
    width=0.95\linewidth, height=6.5cm,
    ybar, bar width=10pt,
    xlabel={$T$ (s)}, ylabel={Effectif},
    xmin=1.935, xmax=2.045,
    ymin=0, ymax=8,
    xtick={1.95,1.96,1.97,1.98,1.99,2.00,2.01,2.02,2.03},
    x tick label style={/pgf/number format/fixed, /pgf/number format/precision=2, font=\small},
    ytick={0,2,4,6,8},
    axis lines=left, grid=none,
    every axis plot/.append style={fill=popblue!55, draw=popblue}
]
\addplot coordinates {(1.95,2) (1.96,2) (1.97,3) (1.98,3) (1.99,3) (2.00,3) (2.01,2) (2.02,1) (2.03,1)};
% Bande ± s autour de la moyenne (1.9855 ± 0.021)
\draw[fill=poporange!25, draw=none] (axis cs:1.9645,0) rectangle (axis cs:2.0065,8);
\addplot coordinates {(1.95,2) (1.96,2) (1.97,3) (1.98,3) (1.99,3) (2.00,3) (2.01,2) (2.02,1) (2.03,1)};
\draw[popdark, very thick, dashed] (axis cs:1.9855,0) -- (axis cs:1.9855,8) node[above, font=\small] {$\bar{T} = 1{,}986$ s};
\draw[poporange, thick] (axis cs:1.9645,8) -- (axis cs:2.0065,8) node[midway, above, font=\small, text=poporange] {bande $\bar{T} \pm s$};
\end{axis}
\end{tikzpicture}
\end{center}
\legende{Histogramme des 20 mesures de la période du pendule. La ligne verticale pointillée marque la moyenne $\bar{T}$ ; la bande orange matérialise l'intervalle $[\bar{T}-s\,;\,\bar{T}+s]$, qui contient ici 12 mesures sur 20 (60 \%), compatible avec les 68 \% attendus pour une loi normale.}
\end{popfigure}

\begin{savaistu}[{L'intervalle $[\bar{x}-s\,;\,\bar{x}+s]$ : la règle des 68 \% (complément Bac)}]
Lorsque les erreurs sont purement aléatoires et nombreuses, la distribution des mesures s'approche d'une courbe en cloche appelée \textbf{loi normale} (ou gaussienne). Pour une telle distribution, environ \textbf{68 \%} des mesures tombent dans $[\bar{x}-s\,;\,\bar{x}+s]$, et environ \textbf{95 \%} dans $[\bar{x}-2s\,;\,\bar{x}+2s]$. C'est un moyen rapide de vérifier la cohérence d'une série : dans notre exemple, $12/20 = 60\,\%$, tout à fait compatible avec 68 \% pour un effectif modeste. Ce résultat n'est pas exigible, mais il éclaire le sens physique de $s$.
\end{savaistu}

\courssub{L'incertitude-type sur la moyenne : $u = s/\sqrt{n}$}

Venons-en à la question décisive : de combien peut-on faire confiance à la \emph{moyenne} ?

Reprenons l'intuition de la compensation. Chaque mesure s'écrit $x_i = x_{\text{vrai}} + \varepsilon_i$, où $\varepsilon_i$ est l'erreur aléatoire, de dispersion typique $s$. La moyenne vaut :
\[
\bar{x} = x_{\text{vrai}} + \frac{1}{n}\sum_{i=1}^{n}\varepsilon_i.
\]
Les erreurs $\varepsilon_i$ étant indépendantes et de signes variables, leur somme ne croît pas comme $n$ mais seulement comme $\sqrt{n}$ (c'est le comportement typique d'une « marche au hasard » : en $n$ pas de longueur $s$, on ne s'éloigne en moyenne que de $s\sqrt{n}$). Divisée par $n$, l'erreur résiduelle sur la moyenne décroît donc comme $s/\sqrt{n}$. La démonstration rigoureuse (propagation des variances) est hors programme ; l'essentiel est de retenir le résultat et son sens.

\begin{propriete}[Incertitude-type sur la moyenne (admise)]
Pour une série de $n$ mesures indépendantes d'écart-type expérimental $s$, l'\textbf{incertitude-type} sur la moyenne est :
\[
u(\bar{x}) = \frac{s}{\sqrt{n}}.
\]
C'est l'incertitude (type A) à associer au résultat $\bar{x}$. La démonstration complète est hors programme, mais l'idée directrice (compensation partielle des erreurs aléatoires, croissance en $\sqrt{n}$) a été guidée ci-dessus.
\end{propriete}

\begin{attention}[$s$ et $u$ : deux grandeurs à ne jamais confondre]
\begin{itemize}
\item $s$ mesure la dispersion des mesures \textbf{individuelles} : il ne dépend (presque) pas de $n$ et ne tend \emph{pas} vers zéro quand on accumule les mesures. Refaire 1000 mesures ne rendra pas le chronométrage humain plus précis !
\item $u(\bar{x}) = s/\sqrt{n}$ mesure l'incertitude sur la \textbf{moyenne} : elle diminue quand $n$ augmente. C'est elle qui figure dans le résultat final $\bar{x} \pm u$.
\end{itemize}
Confondre les deux (par exemple écrire le résultat sous la forme $\bar{x} \pm s$) surestime l'incertitude d'un facteur $\sqrt{n}$ ; c'est une faute classique de copie.
\end{attention}

\begin{popfigure}
\begin{center}
\begin{tikzpicture}
\begin{axis}[
    width=0.9\linewidth, height=6.5cm,
    xlabel={Nombre de mesures $n$},
    ylabel={$u(\bar{x}) = s/\sqrt{n}$ (s)},
    domain=1:100, samples=200,
    xmin=0, xmax=100, ymin=0, ymax=0.025,
    xtick={0,10,20,40,60,80,100},
    ytick={0,0.005,0.01,0.015,0.02,0.025},
    y tick label style={/pgf/number format/fixed, /pgf/number format/precision=3, font=\small},
    axis lines=left, grid=major, grid style={popline!40}
]
\addplot[popcurve]{0.021/sqrt(max(0,x))};
\addplot[only marks, mark=*, mark size=2pt, poporange] coordinates {(1,0.021) (4,0.0105) (10,0.0066) (20,0.0047) (100,0.0021)};
\node[font=\small, text=popdark, anchor=west] at (axis cs:30,0.017) {gain de plus en plus faible :};
\node[font=\small, text=popdark, anchor=west] at (axis cs:30,0.0145) {passer de 1 à 4 mesures divise $u$ par 2,};
\node[font=\small, text=popdark, anchor=west] at (axis cs:30,0.012) {il en faut 100 pour le diviser par 10 !};
\end{axis}
\end{tikzpicture}
\end{center}
\legende{Décroissance de l'incertitude-type $u = s/\sqrt{n}$ en fonction du nombre $n$ de mesures (ici $s = \SI{0,021}{s}$). La courbe en $1/\sqrt{n}$ montre un rendement décroissant : au-delà d'une dizaine de mesures, le gain devient faible comparé au temps investi.}
\end{popfigure}

Ce graphique enseigne une leçon de bon sens expérimental : pour diviser l'incertitude par 2, il faut \emph{quadrupler} le nombre de mesures. En pratique, en TP, une dizaine de mesures constitue un excellent compromis entre qualité statistique et temps disponible.

\courssub{Détection d'une valeur aberrante}

Que faire si, au milieu d'une série régulière, apparaît une mesure très éloignée des autres (par exemple $T = \SI{2,25}{s}$ dans notre série) ?

\begin{methode}[Conduite à tenir face à une valeur suspecte]
\begin{enumerate}
\item \textbf{Ne jamais éliminer en silence.} Une valeur écartée ne se supprime pas parce qu'elle « dérange » : ce serait fausser le résultat.
\item \textbf{Chercher une cause identifiable} : mauvais comptage des oscillations, chronomètre relâché trop tôt, lecture erronée. Si une cause expérimentale certaine est identifiée, on peut écarter la mesure \emph{en le signalant dans le compte rendu}.
\item \textbf{Appliquer un critère objectif sinon} : une règle simple (non exigible) consiste à considérer comme aberrante une valeur située à plus de $3s$ de la moyenne (probabilité inférieure à 0,3 \% pour une gaussienne).
\item \textbf{Recalculer} $\bar{x}$, $s$ et $u$ sur la série nettoyée, et comparer avec le résultat initial pour discuter l'influence de la mesure écartée.
\end{enumerate}
\end{methode}

Dans notre exemple, $T = \SI{2,25}{s}$ s'écarterait de $\bar{T}$ de $\SI{0,26}{s} \approx 12\,s$ : c'est typiquement une erreur de comptage (une demi-oscillation en trop vaut environ \SI{1}{s}... ici plutôt un faux départ du chronomètre). On l'élimine en le justifiant.

\courssub{Synthèse : le protocole complet et un exemple résolu}

\begin{methode}[Traiter une série de mesures en 4 étapes]
\begin{enumerate}
\item \textbf{Calculer la moyenne} $\bar{x} = \dfrac{1}{n}\sum_i x_i$ (meilleure estimation du mesurande).
\item \textbf{Calculer l'écart-type expérimental} $s = \sqrt{\dfrac{\sum_i (x_i-\bar{x})^2}{n-1}}$ (dispersion des mesures, diviseur $n-1$ !).
\item \textbf{En déduire l'incertitude-type} $u(\bar{x}) = \dfrac{s}{\sqrt{n}}$ (incertitude sur la moyenne).
\item \textbf{Exprimer le résultat} sous la forme $x = \bar{x} \pm u(\bar{x})$, avec $u$ arrondie à un (ou deux) chiffre(s) significatif(s) et $\bar{x}$ arrondie en cohérence (même position décimale) — règles détaillées à la section 4.
\end{enumerate}
\end{methode}

\begin{exoresolu}[Période d'un pendule : 10 mesures]
\textbf{Énoncé.} Un élève mesure 10 fois la période d'un pendule simple et obtient (en s) : 1,98 ; 1,96 ; 2,01 ; 1,99 ; 1,97 ; 2,00 ; 1,98 ; 1,99 ; 1,96 ; 1,98.
\begin{enumerate}
\item Calculer la moyenne $\bar{T}$.
\item Calculer l'écart-type expérimental $s$.
\item En déduire l'incertitude-type $u(\bar{T})$ et exprimer le résultat.
\item Combien de mesures faudrait-il pour diviser $u$ par 2 ?
\end{enumerate}
\tcblower
\textbf{Solution détaillée.}
\begin{enumerate}
\item Somme des mesures : $\sum_i T_i = \SI{19,82}{s}$. D'où :
\[
\bar{T} = \frac{19,82}{10} = \SI{1,982}{s}.
\]
\item On dresse le tableau des écarts au carré $(T_i - \bar{T})^2$ ; leur somme vaut $\sum_i (T_i - \bar{T})^2 \approx 2,36 \times 10^{-3}\,\text{s}^2$. Donc :
\[
s = \sqrt{\frac{2,36\times 10^{-3}}{10-1}} = \sqrt{2,62\times 10^{-4}} \approx \SI{0,0162}{s}.
\]
\item Incertitude-type sur la moyenne :
\[
u(\bar{T}) = \frac{s}{\sqrt{n}} = \frac{0,0162}{\sqrt{10}} \approx \SI{0,0051}{s}.
\]
On arrondit $u$ à un chiffre significatif : $u = \SI{0,005}{s}$. La moyenne s'arrondit au millième (même rang décimal) : $\bar{T} = \SI{1,982}{s}$.
Résultat : $T = 1{,}982 \pm 0{,}005$ s. Notons que $u$ est environ trois fois plus petite que $s$ : la moyenne est mieux connue que chaque mesure individuelle.
\item Puisque $u \propto 1/\sqrt{n}$, diviser $u$ par 2 exige de multiplier $n$ par 4 : il faudrait $n = 40$ mesures.
\end{enumerate}
\end{exoresolu}

\begin{aretenir}[L'essentiel de la section]
\begin{itemize}
\item Répéter une mesure réduit les effets des erreurs \textbf{aléatoires} (pas des erreurs systématiques).
\item La \textbf{moyenne} $\bar{x} = \frac{1}{n}\sum_i x_i$ est la meilleure estimation du mesurande.
\item L'\textbf{écart-type expérimental} $s = \sqrt{\frac{\sum_i (x_i-\bar{x})^2}{n-1}}$ mesure la dispersion (diviseur $n-1$).
\item L'\textbf{incertitude-type sur la moyenne} est $u(\bar{x}) = \frac{s}{\sqrt{n}}$ : c'est elle qui accompagne le résultat final.
\item Le gain statistique décroît en $1/\sqrt{n}$ : une dizaine de mesures suffit en général au laboratoire.
\end{itemize}
\end{aretenir}

Nous disposons désormais de tous les outils pour quantifier la qualité d'une série de mesures. Reste à apprendre à \emph{présenter} proprement le résultat : combien de chiffres écrire, comment arrondir l'incertitude, comment comparer deux résultats entre eux ? C'est l'objet de la section suivante : l'expression d'un résultat de mesure.

\coursec{Expression d'un résultat de mesure}

\courssub{Transition depuis la section précédente}

Nous venons de voir comment, à partir d'une \emph{série de mesures}, on construit la meilleure estimation de la grandeur recherchée : la \textbf{moyenne} $\bar{x}$, et comment on quantifie sa dispersion grâce à l'\textbf{écart-type} $s$ et à l'\textbf{incertitude-type} $u(\bar{x}) = s/\sqrt{n}$. Nous disposons donc maintenant de deux nombres : une valeur optimale $\bar{x}$ et une incertitude $u(\bar{x})$. Il reste à les \emph{communiquer} correctement. Une mesure sans incertitude n'a pas de sens scientifique ; une incertitude mal arrondie ou mal alignée rend le résultat illisible, voire trompeur. Cette section te donne les règles universelles pour écrire un résultat sous la forme normalisée $X = (x \pm \Delta X)\,\text{unité}$, avec le bon nombre de \textbf{chiffres significatifs}, la bonne \textbf{incertitude relative} et la capacité de \textbf{comparer} deux résultats.

\courssub{Écriture normalisée d'un résultat}

\begin{definition}[Forme normalisée d'un résultat de mesure]
Tout résultat de mesure d'une grandeur $X$ s'écrit sous la forme :
\[
X = (x \pm \Delta X)\; \text{unité}
\]
où :
\begin{itemize}
    \item $x$ est la \textbf{valeur estimée} (moyenne pour une série, lecture unique pour une mesure simple) ;
    \item $\Delta X$ est l'\textbf{incertitude absolue} (notée $u(x)$ ou $U$ selon le contexte) ;
    \item l'unité est celle du Système international (SI) ou un multiple/sous-multiple approprié.
\end{itemize}
L'intervalle de confiance associé (niveau de confiance $\approx 68\,\%$ pour un facteur d'élargissement $k=1$) est $[\,x-\Delta X\,;\, x+\Delta X\,]$.
\end{definition}

\begin{propriete}[Règles des chiffres significatifs]
Un chiffre est \textbf{significatif} s'il porte une information fiable sur la grandeur mesurée.
\begin{enumerate}
    \item Tous les chiffres non nuls sont significatifs.
    \item Les zéros \emph{entre} deux chiffres non nuls sont significatifs (ex. : $10,05$ a 4 chiffres significatifs).
    \item Les zéros \emph{à gauche} du premier chiffre non nul ne sont \textbf{jamais} significatifs (ex. : $0,0042$ a 2 chiffres significatifs).
    \item Les zéros \emph{à droite} du dernier chiffre non nul sont significatifs \textbf{seulement s'ils sont après la virgule} ou si une notation scientifique lève l'ambiguïté (ex. : $1200$ est ambigu ; on écrit $1,20\times10^3$ pour signifier 3 chiffres significatifs).
\end{enumerate}
\end{propriete}

\begin{methode}[Écrire un résultat en 3 temps — La méthode « Arrondir–Aligner–Unité »]
\begin{enumerate}
    \item \textbf{Arrondir l'incertitude $\Delta X$} à \textbf{1 chiffre significatif} (parfois 2 si le premier chiffre est $1$ ou $2$, pour éviter une perte d'information trop brutale). L'arrondi se fait \textbf{par excès} (on majore l'incertitude pour rester prudent).
    \item \textbf{Arrondir la valeur $x$ au même rang décimal} que l'incertitude arrondie (alignement des dernières décimales).
    \item \textbf{Ajouter l'unité} (SI de préférence) : $(x \pm \Delta X)\,\text{unité}$.
\end{enumerate}
\end{methode}

\begin{exemplebox}[Application de la méthode]
Soit une série de mesures de la période d'un pendule : $\bar{T} = 2,3478\,\text{s}$ avec $u(\bar{T}) = 0,0123\,\text{s}$.
\begin{enumerate}
    \item $\Delta T = 0,0123\,\text{s} \xrightarrow{\text{1 ch. sig., par excès}} \mathbf{0,02}\,\text{s}$.
    \item $\bar{T} = 2,3478\,\text{s} \xrightarrow{\text{au centième}} \mathbf{2,35}\,\text{s}$.
    \item Résultat final : $\mathbf{T = (2,35 \pm 0,02)\,\text{s}}$.
\end{enumerate}
\end{exemplebox}

\begin{attention}[Piège classique — Désalignement des décimales]
\emph{Ne jamais écrire} : $g = (9,81234 \pm 0,05)\,\text{m}\cdot\text{s}^{-2}$.
La valeur a 5 décimales, l'incertitude 2 : cela sous-entend une précision illusoire.
\emph{Écriture correcte} : $\Delta g = 0,05\,\text{m}\cdot\text{s}^{-2}$ (déjà à 1 chiffre significatif) ; on aligne $g$ au centième : $g = 9,81\,\text{m}\cdot\text{s}^{-2}$.
Résultat : $\mathbf{g = (9,81 \pm 0,05)\,\text{m}\cdot\text{s}^{-2}}$.
\end{attention}

\begin{popfigure}
\begin{tikzpicture}[scale=1.1, font=\small]
% Arbre de décision
\node[draw, rounded corners, fill=popblueL, thick, align=center] (start) at (0,0) {Valeur brute $x$ et\\ incertitude brute $\Delta X$};
\node[draw, diamond, aspect=2.5, fill=poporangeL, thick, align=center, inner sep=1pt] (q1) at (0,-2) {$\Delta X$ a-t-elle 1 ou 2\\ ch. sig. ?};
\node[draw, rounded corners, fill=popgreenL, thick, align=center] (a1) at (-3.8,-4) {Arrondir $\Delta X$ à 1 ch. sig.\\ \textbf{par excès} (2 ch. sig. si le\\ 1er chiffre est 1 ou 2)};
\node[draw, rounded corners, fill=popgreenL, thick, align=center] (a2) at (3.8,-4) {Conserver $\Delta X$\\ telle quelle};
\node[draw, rounded corners, fill=poppurpleL, thick, align=center] (q2) at (0,-6) {Arrondir $x$ au \textbf{même rang}\\ \textbf{décimal} que $\Delta X$ arrondie};
\node[draw, rounded corners, fill=poppinkL, thick, align=center] (end) at (0,-7.6) {Écrire $X = (x \pm \Delta X)\;\text{unité SI}$};

\draw[->, thick] (start) -- (q1);
\draw[->, thick] (q1) -- node[left, pos=0.5] {Non} (a1);
\draw[->, thick] (q1) -- node[right, pos=0.5] {Oui} (a2);
\draw[->, thick] (a1) |- (q2);
\draw[->, thick] (a2) |- (q2);
\draw[->, thick] (q2) -- (end);
\end{tikzpicture}
\legende{Arbre de décision pour l'écriture correcte d'un résultat de mesure.}
\end{popfigure}

\courssub{Incertitude relative et qualité de la mesure}

\begin{definition}[Incertitude relative]
L'\textbf{incertitude relative} d'une mesure $X = (x \pm \Delta X)$ est le quotient sans dimension :
\[
\frac{\Delta X}{x} \quad \text{souvent exprimé en pourcentage :} \quad \frac{\Delta X}{x} \times 100\,\%.
\]
Elle caractérise la \textbf{qualité} (précision relative) de la mesure, indépendamment de l'unité et de l'ordre de grandeur de la grandeur mesurée.
\end{definition}

\begin{propriete}[Critères de qualité usuels au lycée]
\begin{itemize}
    \item $\dfrac{\Delta X}{x} < 1\,\%$ : mesure \textbf{très précise} (mesures de laboratoire avec du matériel de qualité).
    \item $1\,\% \le \dfrac{\Delta X}{x} \le 5\,\%$ : mesure \textbf{correcte} (typique des TP de lycée).
    \item $\dfrac{\Delta X}{x} > 10\,\%$ : mesure \textbf{peu précise} — il faut améliorer le protocole ou l'instrument.
\end{itemize}
\end{propriete}

\begin{exemplebox}[Calcul d'incertitude relative]
Une masse $m = (25,2 \pm 0,2)\,\text{g}$ (obtenue après arrondi de $25,183 \pm 0,12\,\text{g}$).
\[
\frac{\Delta m}{m} = \frac{0,2}{25,2} \approx 0,00794 = \mathbf{0,79\,\%}.
\]
C'est une mesure de \textbf{très bonne qualité} (incertitude relative $< 1\,\%$).
\end{exemplebox}

\begin{savaistu}[Au Baccalauréat : rigueur = points]
Au Bac camerounais (série C), un résultat numérique \textbf{sans unité} ou avec des \textbf{chiffres significatifs incohérents} (valeur plus précise que l'incertitude, ou incertitude non arrondie par excès) est \textbf{pénalisé} même si le raisonnement physique est correct. Apprends à \emph{toujours} présenter tes résultats sous la forme $(x \pm \Delta X)\,\text{unité}$ avec l'alignement décimal : c'est un réflexe de scientifique qui rapporte des points faciles.
\end{savaistu}

\courssub{Comparaison de deux résultats — Compatibilité}

\begin{definition}[Compatibilité de deux résultats]
Soit deux mesures indépendantes d'une même grandeur $X$ :
\[
X_1 = (x_1 \pm \Delta X_1), \qquad X_2 = (x_2 \pm \Delta X_2).
\]
On dit qu'elles sont \textbf{compatibles} (au niveau de confiance $k=1$) si leurs intervalles de confiance \textbf{se chevauchent} :
\[
[x_1 - \Delta X_1\,;\, x_1 + \Delta X_1] \;\cap\; [x_2 - \Delta X_2\,;\, x_2 + \Delta X_2] \neq \varnothing.
\]
Condition équivalente : $|x_1 - x_2| \le \Delta X_1 + \Delta X_2$.
\end{definition}

\begin{popfigure}
\begin{tikzpicture}[scale=1.2, font=\small]
% Cas compatible
\draw[thick, popmuted] (0,1.5) -- (6,1.5) node[right] {Échelle};
\draw[thick, popblue] (1,1.3) -- (1,1.7);
\draw[thick, popblue] (4,1.3) -- (4,1.7);
\draw[popblue, very thick] (1,1.5) -- (4,1.5) node[midway, above=2pt] {$X_1$};
\fill[popblue] (2.5,1.5) circle (1.5pt) node[below=2pt] {$x_1$};

\draw[thick, poporange] (2,0.4) -- (2,0.8);
\draw[thick, poporange] (5,0.4) -- (5,0.8);
\draw[poporange, very thick] (2,0.6) -- (5,0.6) node[midway, below=2pt] {$X_2$};
\fill[poporange] (3.5,0.6) circle (1.5pt) node[above=2pt] {$x_2$};

\node[draw, rounded corners, fill=popgreenL, align=center] at (3, -0.7) {Intervalles se chevauchent\\ $\Rightarrow$ \textbf{COMPATIBLES}};

% Cas incompatible (décalé vers la droite)
\begin{scope}[xshift=8cm]
\draw[thick, popmuted] (0,1.5) -- (6,1.5) node[right] {Échelle};
\draw[thick, popblue] (1,1.3) -- (1,1.7);
\draw[thick, popblue] (3,1.3) -- (3,1.7);
\draw[popblue, very thick] (1,1.5) -- (3,1.5) node[midway, above=2pt] {$X_1$};
\fill[popblue] (2,1.5) circle (1.5pt) node[below=2pt] {$x_1$};

\draw[thick, poporange] (4,0.4) -- (4,0.8);
\draw[thick, poporange] (6,0.4) -- (6,0.8);
\draw[poporange, very thick] (4,0.6) -- (6,0.6) node[midway, below=2pt] {$X_2$};
\fill[poporange] (5,0.6) circle (1.5pt) node[above=2pt] {$x_2$};

\node[draw, rounded corners, fill=poppinkL, align=center] at (3, -0.7) {Intervalles disjoints\\ $\Rightarrow$ \textbf{INCOMPATIBLES}};
\end{scope}
\end{tikzpicture}
\legende{Comparaison visuelle de deux résultats : chevauchement (compatibles) ou séparation (incompatibles).}
\end{popfigure}

\begin{exemplebox}[Test de compatibilité]
Deux élèves mesurent l'accélération de la pesanteur $g$ au lycée de Ngaoundéré :
\begin{itemize}
    \item Élève A : $g_A = (9,81 \pm 0,05)\,\text{m}\cdot\text{s}^{-2}$ $\to$ intervalle $[9,76\,;\,9,86]$.
    \item Élève B : $g_B = (9,72 \pm 0,04)\,\text{m}\cdot\text{s}^{-2}$ $\to$ intervalle $[9,68\,;\,9,76]$.
\end{itemize}
Les intervalles se touchent en $9,76$ (borne commune) : l'intersection est non vide, les résultats sont \textbf{compatibles}. Si B avait trouvé $g_B = (9,70 \pm 0,04)\,\text{m}\cdot\text{s}^{-2}$, l'intervalle serait $[9,66\,;\,9,74]$ : \textbf{incompatible} avec A.
\end{exemplebox}

\begin{exoresolu}[Traitement complet d'une série de mesures]
\textbf{Énoncé.} Un élève mesure 5 fois la longueur $L$ d'un fil de cuivre avec un pied à coulisse (appréciation $0,01\,\text{mm}$). Résultats en mm : $124,52$ ; $124,55$ ; $124,48$ ; $124,51$ ; $124,54$.
\begin{enumerate}
    \item Calculer la moyenne $\bar{L}$ et l'écart-type $s$.
    \item Déterminer l'incertitude-type $u(\bar{L})$.
    \item Exprimer le résultat final selon les règles.
    \item Calculer l'incertitude relative en $\%$.
    \item La valeur de référence du constructeur est $L_{\text{ref}} = 124,50\,\text{mm}$. Les résultats sont-ils compatibles ?
\end{enumerate}
\tcblower
\textbf{Solution détaillée.}
\begin{enumerate}
    \item \textbf{Moyenne :}
    \[
    \bar{L} = \frac{124,52+124,55+124,48+124,51+124,54}{5} = \frac{622,60}{5} = 124,52\,\text{mm}.
    \]
    \textbf{Écart-type expérimental (formule corrigée, diviseur $n-1$) :}
    \begin{align*}
    s &= \sqrt{\frac{1}{n-1}\sum_{i=1}^n (L_i - \bar{L})^2} \\
      &= \sqrt{\frac{1}{4}\big[(0,00)^2 + (0,03)^2 + (-0,04)^2 + (-0,01)^2 + (0,02)^2\big]} \\
      &= \sqrt{\frac{1}{4}(0 + 0,0009 + 0,0016 + 0,0001 + 0,0004)} \\
      &= \sqrt{\frac{0,0030}{4}} = \sqrt{0,00075} \approx 0,0274\,\text{mm}.
    \end{align*}
    \item \textbf{Incertitude-type sur la moyenne :}
    \[
    u(\bar{L}) = \frac{s}{\sqrt{n}} = \frac{0,0274}{\sqrt{5}} \approx 0,0123\,\text{mm}.
    \]
    \item \textbf{Écriture normalisée (méthode en 3 temps) :}
    \begin{itemize}
        \item $\Delta L = 0,0123\,\text{mm} \xrightarrow{\text{1 ch. sig., par excès}} \mathbf{0,02}\,\text{mm}$.
        \item $\bar{L} = 124,52\,\text{mm}$ : déjà exprimée au centième, alignée sur $\Delta L$.
        \item Résultat : $\mathbf{L = (124,52 \pm 0,02)\,\text{mm}}$.
    \end{itemize}
    \item \textbf{Incertitude relative :}
    \[
    \frac{\Delta L}{\bar{L}} = \frac{0,02}{124,52} \approx 1,6 \times 10^{-4} = \mathbf{0,016\,\%}.
    \]
    Mesure de \textbf{très haute qualité} (instrument précis et répétition des mesures).
    \item \textbf{Compatibilité :} l'intervalle de confiance de l'élève est $[124,50\,;\,124,54]\,\text{mm}$. La valeur constructeur $124,50\,\text{mm}$ (d'incertitude supposée négligeable) appartient à cet intervalle $\Rightarrow$ les résultats sont \textbf{compatibles}.
\end{enumerate}
\end{exoresolu}

\begin{exercice}[Entraînement — Expression et comparaison]
Deux groupes d'élèves mesurent la résistance $R$ d'une thermistance à $25\,^\circ\text{C}$ avec un ohmmètre. Séries obtenues (en $\Omega$) :

Groupe 1 : $985$ ; $992$ ; $988$ ; $990$ ; $987$.

Groupe 2 : $1012$ ; $1008$ ; $1015$ ; $1010$ ; $1005$.
\begin{enumerate}
    \item Pour chaque groupe, donner $R$ sous la forme normalisée $(R \pm \Delta R)\,\Omega$ et l'incertitude relative en $\%$.
    \item Les deux résultats sont-ils compatibles ? Justifier par le test des intervalles.
    \item La notice du fabricant indique $R_{25} = 1000\,\Omega \pm 2\,\%$. Quel groupe est le plus proche de la valeur nominale ?
\end{enumerate}
\end{exercice}

\begin{corrige}[Exercice — Expression et comparaison]
\textbf{Groupe 1 :}
\[
\bar{R}_1 = \frac{985+992+988+990+987}{5} = \frac{4942}{5} = 988,4\,\Omega.
\]
Écarts à la moyenne : $-3,4$ ; $3,6$ ; $-0,4$ ; $1,6$ ; $-1,4$ (en $\Omega$). Carrés : $11,56$ ; $12,96$ ; $0,16$ ; $2,56$ ; $1,96$, de somme $29,2$.
\[
s_1 = \sqrt{\frac{29,2}{4}} = \sqrt{7,3} \approx 2,7\,\Omega ; \qquad u_1 = \frac{s_1}{\sqrt{5}} \approx \frac{2,7}{2,236} \approx 1,2\,\Omega.
\]
Arrondi par excès à 1 chiffre significatif : $\Delta R_1 = 2\,\Omega$ ; on aligne la moyenne à l'unité : $\bar{R}_1 \to 988\,\Omega$.
\[
\boxed{R_1 = (988 \pm 2)\,\Omega} ; \qquad \frac{\Delta R_1}{R_1} = \frac{2}{988} \approx 0,20\,\%.
\]

\textbf{Groupe 2 :}
\[
\bar{R}_2 = \frac{1012+1008+1015+1010+1005}{5} = \frac{5050}{5} = 1010\,\Omega.
\]
Écarts à la moyenne : $2$ ; $-2$ ; $5$ ; $0$ ; $-5$ (en $\Omega$). Carrés : $4$ ; $4$ ; $25$ ; $0$ ; $25$, de somme $58$.
\[
s_2 = \sqrt{\frac{58}{4}} = \sqrt{14,5} \approx 3,8\,\Omega ; \qquad u_2 = \frac{s_2}{\sqrt{5}} \approx \frac{3,8}{2,236} \approx 1,7\,\Omega.
\]
Arrondi par excès : $\Delta R_2 = 2\,\Omega$ ; la moyenne est déjà alignée à l'unité.
\[
\boxed{R_2 = (1010 \pm 2)\,\Omega} ; \qquad \frac{\Delta R_2}{R_2} = \frac{2}{1010} \approx 0,20\,\%.
\]

\textbf{Compatibilité :} intervalle du Groupe 1 : $[986\,;\,990]\,\Omega$ ; intervalle du Groupe 2 : $[1008\,;\,1012]\,\Omega$. Les intervalles sont \textbf{disjoints} $\Rightarrow$ les deux résultats sont \textbf{incompatibles}. Il existe une cause systématique : température réelle différente de $25\,^\circ\text{C}$ pour l'un des groupes, ohmmètre mal étalonné, mauvais contacts, etc.

\textbf{Valeur nominale :} $R_{25} = 1000\,\Omega \pm 2\,\%$ correspond à l'intervalle $[980\,;\,1020]\,\Omega$. Les deux groupes sont compatibles avec la notice du fabricant, mais le Groupe 1 ($988\,\Omega$, écart de $12\,\Omega$) est plus proche de la valeur nominale que le Groupe 2 ($1010\,\Omega$, écart de $10\,\Omega$ en valeur absolue mais de signe opposé — en fait $1010$ est à $10\,\Omega$ de $1000$ et $988$ à $12\,\Omega$ : c'est donc le \textbf{Groupe 2} qui est le plus proche de la valeur nominale).
\end{corrige}

\courssub{Synthèse des règles d'or}

\begin{aretenir}[Check-list avant de rendre une copie]
\begin{enumerate}
    \item L'incertitude a \textbf{1 (ou 2) chiffre(s) significatif(s)} et est arrondie \textbf{par excès}.
    \item La valeur est arrondie au \textbf{même rang décimal} que l'incertitude (alignement).
    \item L'\textbf{unité SI} est présente et correctement écrite ($\text{m}\cdot\text{s}^{-2}$, et non $\text{m/s}^2$).
    \item L'incertitude relative $\Delta X / x$ est calculée et interprétée (qualité de la mesure).
    \item Pour comparer deux résultats : on teste le \textbf{chevauchement des intervalles} $[x-\Delta X\,;\,x+\Delta X]$.
\end{enumerate}
\end{aretenir}

\begin{attention}[Erreurs fatales à ne plus commettre]
\begin{itemize}
    \item Écrire $v = (12,3456 \pm 0,1)\,\text{m}\cdot\text{s}^{-1}$ (désalignement des décimales).
    \item Donner $\Delta X = 0,0123$ sans arrondir (trop de chiffres significatifs).
    \item Oublier l'unité ou écrire $\text{m/s}^2$ au lieu de $\text{m}\cdot\text{s}^{-2}$.
    \item Confondre l'écart-type $s$ (dispersion des mesures) et l'incertitude-type $u(\bar{x}) = s/\sqrt{n}$ (incertitude sur la moyenne).
    \item Affirmer « les résultats sont compatibles car les moyennes sont proches » sans tester le chevauchement des intervalles.
\end{itemize}
\end{attention}

\coursec{Le Système international d'unités (SI)}

Dans la section précédente, nous avons appris à exprimer un résultat de mesure sous la forme $x = \bar{x} \pm \Delta x$, avec le bon nombre de chiffres significatifs. Mais un résultat tel que $L = (2{,}35 \pm 0{,}02)$ ne veut strictement rien dire tant qu'on n'a pas précisé \emph{dans quelle unité} $L$ est exprimé : $2{,}35$~m ? $2{,}35$~cm ? $2{,}35$~km ? L'unité n'est pas un accessoire : elle fait partie intégrante du résultat. Encore faut-il que tout le monde, à Yaoundé comme à Tokyo, parle le même langage. C'est précisément le rôle du \cle{Système international d'unités} (SI), adopté en 1960 et révisé en 2019 : fournir à tous les scientifiques, ingénieurs et techniciens de la planète un socle commun de sept unités de base, à partir desquelles toutes les autres se construisent.

\courssub{Les sept unités de base du SI}

Toute grandeur physique que tu rencontreras en Terminale C — vitesse, force, champ électrique, capacité, activité radioactive, accélération, impédance — peut être exprimée à l'aide de seulement \textbf{sept grandeurs de base}, choisies pour être indépendantes les unes des autres et faciles à mesurer avec précision.

\begin{definition}[Unité de base du SI]
Une \cle{unité de base} est l'unité de mesure de l'une des sept grandeurs fondamentales du Système international : longueur, masse, durée, intensité de courant électrique, température thermodynamique, quantité de matière et intensité lumineuse. Ces unités sont définies de manière exacte à partir de constantes fondamentales de la nature.
\end{definition}

\begin{popfigure}
\begin{center}
\begin{tikzpicture}
\fill[popblueL,rounded corners=2pt] (0,0) rectangle (12.5,1);
\node[text=popdark,font=\bfseries] at (2.2,0.5) {Grandeur de base};
\node[text=popdark,font=\bfseries] at (6.6,0.5) {Nom de l'unité};
\node[text=popdark,font=\bfseries] at (9.6,0.5) {Symbole};
\node[text=popdark,font=\bfseries] at (11.5,0.5) {Dimension};
\foreach \y/\g/\n/\s/\d in {
  -0.6/longueur/mètre/m/L,
  -1.6/masse/kilogramme/kg/M,
  -2.6/durée/seconde/s/T,
  -3.6/intensité de courant/ampère/A/I,
  -4.6/température/kelvin/K/$\Theta$,
  -5.6/quantité de matière/mole/mol/N,
  -6.6/intensité lumineuse/candela/cd/J}{
  \fill[popcream] (0,\y-0.5) rectangle (12.5,\y+0.5);
  \node[anchor=west,text=popdark] at (0.3,\y) {\g};
  \node[text=popblue] at (6.6,\y) {\n};
  \node[text=poporange,font=\bfseries] at (9.6,\y) {\s};
  \node[text=poppurple] at (11.5,\y) {[\d]};
}
\draw[popline] (0,1) rectangle (12.5,-7.1);
\draw[popline] (4.4,1) -- (4.4,-7.1);
\draw[popline] (8.8,1) -- (8.8,-7.1);
\draw[popline] (10.6,1) -- (10.6,-7.1);
\end{tikzpicture}
\end{center}
\legende{Les sept unités de base du Système international, avec la lettre de dimension associée (nous y reviendrons dans la section suivante).}
\end{popfigure}

Quelques remarques essentielles sur ces unités :

\begin{itemize}
\item Le \cle{kilogramme} (symbole kg) est la seule unité de base dont le nom contient un préfixe (« kilo »). On dit donc « un kilogramme », mais le symbole seul est \textbf{kg}, en minuscules.
\item La \cle{seconde} (s) est définie à partir de la fréquence de transition de l'atome de césium 133 : c'est l'unité de base la plus précisément mesurée au monde.
\item L'\cle{ampère} (A) est l'unité de base électrique : le volt, l'ohm, le coulomb en \emph{découlent}, et non l'inverse.
\item Le \cle{kelvin} (K) s'écrit \emph{sans} le mot « degré » et sans le symbole ${}^\circ$ : on dit « 300 kelvins », on écrit $T = 300$~K. La relation avec le degré Celsius est $T(\mathrm{K}) = \theta({}^\circ\mathrm{C}) + 273{,}15$.
\end{itemize}

\begin{savaistu}[2019 : la révolution des constantes fondamentales]
Jusqu'en 2018, le kilogramme était défini par un objet matériel : le « grand K », un cylindre de platine-iridié conservé sous cloche près de Paris. Problème : cet objet perdait ou gagnait quelques microgrammes au fil des décennies ! Depuis le 20 mai 2019, \textbf{les sept unités de base sont toutes définies à partir de constantes fondamentales immuables de la nature} : la vitesse de la lumière $c = \SI{299792458}{m.s^{-1}}$ (exacte), la constante de Planck $h = \num{6,62607015e-34}$~J$\cdot$s (exacte), la charge élémentaire $e = \num{1,602176634e-19}$~C (exacte), la constante de Boltzmann $k$ (exacte), le nombre d'Avogadro $N_A = \num{6,02214076e23}$~mol$^{-1}$ (exact), etc. Le kilogramme n'est donc plus un objet, mais une \emph{loi de la nature}. Aucun laboratoire au monde, y compris au Cameroun, n'a plus besoin de se référer à un étalon matériel lointain.
\end{savaistu}

\courssub{Les unités dérivées}

\begin{definition}[Unité dérivée cohérente]
Une \cle{unité dérivée} est une unité construite comme produit ou quotient de puissances des unités de base. Elle est dite \cle{cohérente} lorsqu'aucun facteur numérique autre que 1 n'intervient dans sa définition : par exemple, l'unité de vitesse cohérente est le mètre par seconde (m$\cdot$s$^{-1}$), et non le kilomètre par heure.
\end{definition}

Beaucoup d'unités dérivées ont reçu un nom propre, en l'honneur de grands scientifiques. En voici les principales au programme de Terminale C, auxquelles on ajoute l'accélération, le champ électrique et l'impédance qui n'ont pas de nom propre mais sont essentielles :

\begin{center}
\begin{tabularx}{\linewidth}{|l|l|X|l|}
\hline
\rowcolor{popblueL} \textbf{Grandeur} & \textbf{Nom} & \textbf{Expression en unités de base} & \textbf{Symbole} \\
\hline
Accélération & -- & m$\cdot$s$^{-2}$ & -- \\
\hline
Force & newton & kg$\cdot$m$\cdot$s$^{-2}$ & N \\
\hline
Énergie, travail & joule & kg$\cdot$m$^{2}\cdot$s$^{-2}$ & J \\
\hline
Puissance & watt & kg$\cdot$m$^{2}\cdot$s$^{-3}$ & W \\
\hline
Charge électrique & coulomb & A$\cdot$s & C \\
\hline
Tension électrique & volt & kg$\cdot$m$^{2}\cdot$s$^{-3}\cdot$A$^{-1}$ & V \\
\hline
Champ électrique & -- & kg$\cdot$m$\cdot$s$^{-3}\cdot$A$^{-1}$ (V$\cdot$m$^{-1}$) & -- \\
\hline
Résistance, Impédance & ohm & kg$\cdot$m$^{2}\cdot$s$^{-3}\cdot$A$^{-2}$ & $\Omega$ \\
\hline
Capacité & farad & kg$^{-1}\cdot$m$^{-2}\cdot$s$^{4}\cdot$A$^{2}$ & F \\
\hline
Inductance & henry & kg$\cdot$m$^{2}\cdot$s$^{-2}\cdot$A$^{-2}$ & H \\
\hline
Champ magnétique & tesla & kg$\cdot$s$^{-2}\cdot$A$^{-1}$ & T \\
\hline
Activité radioactive & becquerel & s$^{-1}$ & Bq \\
\hline
\end{tabularx}
\end{center}

Comment obtient-on ces expressions ? En partant d'une \textbf{loi physique} reliant la grandeur à des grandeurs déjà connues. C'est un excellent entraînement à l'analyse dimensionnelle.

\begin{exemplebox}[Le newton et le volt en unités de base]
\textbf{Le newton.} La deuxième loi de Newton s'écrit $\vec{F} = m\,\vec{a}$. La masse s'exprime en kg et l'accélération en m$\cdot$s$^{-2}$. Donc :
\[ 1~\mathrm{N} = 1~\mathrm{kg} \times 1~\mathrm{m.s^{-2}} = 1~\mathrm{kg.m.s^{-2}}. \]

\textbf{Le volt.} Partons de la puissance électrique $P = U\,I$, avec $P$ en watts et $I$ en ampères. Le watt lui-même vient de $P = \dfrac{E}{t}$ : une énergie (joule, soit kg$\cdot$m$^{2}\cdot$s$^{-2}$) divisée par un temps, donc $1~\mathrm{W} = 1~\mathrm{kg.m^{2}.s^{-3}}$. Alors :
\[ U = \frac{P}{I} \quad\Longrightarrow\quad 1~\mathrm{V} = \frac{1~\mathrm{kg.m^{2}.s^{-3}}}{1~\mathrm{A}} = 1~\mathrm{kg.m^{2}.s^{-3}.A^{-1}}. \]
Tu remarques la méthode : on remonte de loi en loi jusqu'aux unités de base.
\end{exemplebox}

\begin{exoresolu}[Exprimer le farad en unités de base]
Énoncé. La capacité $C$ d'un condensateur est définie par $q = C\,u$, où $q$ est la charge (en coulombs) et $u$ la tension (en volts). Exprimer le farad en fonction des unités de base du SI.
\tcblower
Solution détaillée. De $q = C\,u$, on tire $C = \dfrac{q}{u}$, donc $1~\mathrm{F} = 1~\mathrm{C.V^{-1}}$.
\begin{align*}
\text{Charge : } & 1~\mathrm{C} = 1~\mathrm{A.s} \\
\text{Tension : } & 1~\mathrm{V} = 1~\mathrm{kg.m^{2}.s^{-3}.A^{-1}} \\
\text{Donc : } & 1~\mathrm{F} = \frac{1~\mathrm{A.s}}{1~\mathrm{kg.m^{2}.s^{-3}.A^{-1}}} = 1~\mathrm{kg^{-1}.m^{-2}.s^{4}.A^{2}}.
\end{align*}
Vérification de cohérence : les exposants de A sont $1-(-1) = 2$ et ceux de s sont $1-(-3) = 4$. Le résultat est bien homogène à un quotient de grandeurs SI.
\end{exoresolu}

\courssub{Multiples et sous-multiples : les préfixes du SI}

Les grandeurs physiques s'étendent sur des ordres de grandeur immenses : de la taille d'un noyau atomique ($10^{-15}$~m) à la distance Terre--Soleil ($10^{11}$~m). Pour éviter d'écrire des colonnes de zéros, le SI définit des \cle{préfixes} multiplicateurs.

\begin{center}
\begin{tabularx}{\linewidth}{|l|l|l|X|}
\hline
\rowcolor{popblueL} \textbf{Préfixe} & \textbf{Symbole} & \textbf{Facteur} & \textbf{Exemple} \\
\hline
téra & T & $10^{12}$ & 1,2 TWh : production annuelle d'un grand barrage \\
\hline
giga & G & $10^{9}$ & 4 GHz : fréquence d'un réseau 4G/5G \\
\hline
méga & M & $10^{6}$ & 420 MW : puissance du barrage de Nachtigal \\
\hline
kilo & k & $10^{3}$ & 225 km : distance Yaoundé--Douala par la route \\
\hline
milli & m & $10^{-3}$ & 2 mm : diamètre d'une pluie fine \\
\hline
micro & $\mu$ & $10^{-6}$ & 50 $\mu$m : épaisseur d'un cheveu \\
\hline
nano & n & $10^{-9}$ & 500 nm : longueur d'onde de la lumière verte \\
\hline
pico & p & $10^{-12}$ & 12 pF : capacité d'un petit condensateur \\
\hline
\end{tabularx}
\end{center}

\begin{popfigure}
\begin{center}
\begin{tikzpicture}[scale=1]
\draw[-{Stealth},thick,popdark] (-0.5,0) -- (13.2,0);
\foreach \x/\p/\e in {0/p/$10^{-12}$,2/n/$10^{-9}$,4/$\mu$/$10^{-6}$,6/m/$10^{-3}$,8/k/$10^{3}$,10/M/$10^{6}$,12/G/$10^{9}$}{
  \draw[thick,popdark] (\x,-0.12) -- (\x,0.12);
  \node[below,popdark] at (\x,-0.15) {\e};
  \node[above,popblue,font=\bfseries] at (\x,0.15) {\p};
}
\node[above,poppurple,align=center,font=\small] at (0,0.75) {12 pF\\ condensateur};
\node[above,poppurple,align=center,font=\small] at (4,0.75) {50 $\mu$m\\ cheveu};
\node[above,poppurple,align=center,font=\small] at (6,0.75) {2 mm\\ grain de sable};
\node[above,poppurple,align=center,font=\small] at (8,0.75) {225 km\\ Yaoundé--Douala};
\node[above,poppurple,align=center,font=\small] at (10,0.75) {420 MW\\ Nachtigal};
\node[above,poppurple,align=center,font=\small] at (12,0.75) {4 GHz\\ 4G};
\foreach \x in {0,4,6,8,10,12}{
  \draw[poppurple,dashed] (\x,0.12) -- (\x,0.6);
}
\end{tikzpicture}
\end{center}
\legende{Échelle des préfixes SI de $10^{-12}$ à $10^{9}$, avec des exemples du quotidien camerounais. Chaque graduation correspond à un facteur $10^{3}$ par rapport à la voisine.}
\end{popfigure}

\begin{attention}[Préfixes : les pièges classiques]
\begin{itemize}
\item Le symbole du préfixe est \emph{collé} au symbole de l'unité, sans espace : on écrit \textbf{km}, \textbf{MHz}, \textbf{$\mu$m}.
\item Attention à la casse : \textbf{m} = milli ($10^{-3}$) mais \textbf{M} = méga ($10^{6}$) ! Confondre les deux, c'est se tromper d'un facteur $10^{9}$.
\item Dans une conversion au carré ou au cube, le préfixe est lui aussi élevé à la puissance : $1~\mathrm{cm^{2}} = (10^{-2}~\mathrm{m})^{2} = 10^{-4}~\mathrm{m^{2}}$ et $1~\mathrm{cm^{3}} = 10^{-6}~\mathrm{m^{3}}$. Erreur très fréquente au Bac !
\item Pour la masse, le préfixe s'attache au \emph{gramme} et non au kilogramme : on dit « milligramme » (mg), jamais « microkilogramme ».
\end{itemize}
\end{attention}

\courssub{Règles d'écriture des symboles d'unités}

Le Bureau international des poids et mesures (BIPM) impose des règles typographiques strictes, que les correcteurs du Baccalauréat attendent de toi :

\begin{aretenir}[Les cinq règles d'or]
\begin{enumerate}
\item Les symboles s'écrivent en \textbf{minuscules} (m, s, kg, mol), \emph{sauf} s'ils dérivent d'un nom propre de scientifique, auquel cas la première lettre est une \textbf{capitale} : N (Newton), J (Joule), W (Watt), V (Volta), A (Ampère), K (Kelvin), F (Faraday), H (Henry), T (Tesla), Bq (Becquerel).
\item Un symbole est \textbf{invariable} : jamais de « s » au pluriel. On écrit « 5 kg » et non « 5 kgs ».
\item Jamais de point après un symbole (ce n'est pas une abréviation) : « 3 m » et non « 3 m. ».
\item On laisse une \textbf{espace} (espace insécable de préférence) entre la valeur et le symbole : $U = 220$~V.
\item Le produit d'unités se note avec un point médian ou une espace : kg$\cdot$m$\cdot$s$^{-2}$ ; les exposants négatifs remplacent avantageusement la barre de fraction : m$\cdot$s$^{-1}$ plutôt que m/s dans les expressions composées.
\end{enumerate}
\end{aretenir}

\begin{attention}[Erreurs fréquentes à bannir de tes copies]
\begin{itemize}
\item Écrire « Kg » au lieu de « kg » : le K majuscule est réservé au kelvin ; « Kg » se lirait « kelvin-gramme », ce qui n'a aucun sens.
\item Écrire « M » pour mètre : M est le préfixe méga ! Le mètre se note \textbf{m}.
\item Mélanger les unités dans un même calcul : calculer une vitesse avec une distance en km et un temps en secondes donne des km$\cdot$s$^{-1}$, pas des m$\cdot$s$^{-1}$.
\item Écrire « 50°K » : le kelvin ne prend jamais le symbole degré.
\end{itemize}
\end{attention}

\courssub{La règle d'or du calcul numérique : tout convertir en SI}

\begin{methode}[Convertir systématiquement en unités SI avant tout calcul]
\begin{enumerate}
\item \textbf{Repérer} toutes les données numériques de l'énoncé et leurs unités.
\item \textbf{Convertir} chaque donnée en unité SI cohérente : longueurs en m, masses en kg, durées en s, tensions en V, etc. Écris explicitement la conversion sur ta copie (ex. : $d = 2{,}5~\mathrm{cm} = \num{2,5e-2}$~m).
\item \textbf{Effectuer} le calcul littéral puis l'application numérique avec ces valeurs SI.
\item \textbf{Conclure} : le résultat est alors automatiquement exprimé dans l'unité SI cohérente de la grandeur cherchée. On peut ensuite, si l'énoncé le demande, reconvertir (par exemple en km$\cdot$h$^{-1}$).
\end{enumerate}
\end{methode}

\begin{exoresolu}[Énergie cinétique d'une mototaxi : l'importance des conversions]
Énoncé. Une mototaxi (« benskin ») de masse $m = 150$~kg (avec son conducteur) roule à $v = 54$~km$\cdot$h$^{-1}$ dans les rues de Bafoussam. Calculer son énergie cinétique $E_c = \dfrac{1}{2}mv^2$ en joules.
\tcblower
Solution détaillée. \textbf{Étape 1 — Conversion.} La vitesse est donnée en km$\cdot$h$^{-1}$, qui n'est pas une unité SI cohérente :
\[ v = 54~\mathrm{km.h^{-1}} = \frac{54 \times 10^{3}~\mathrm{m}}{3600~\mathrm{s}} = 15~\mathrm{m.s^{-1}}. \]
\textbf{Étape 2 — Application numérique en SI.}
\[ E_c = \frac{1}{2} \times 150~\mathrm{kg} \times (15~\mathrm{m.s^{-1}})^{2} = 75 \times 225 = \num{16875}~\mathrm{J} \approx \num{1,7e4}~\mathrm{J}. \]
\textbf{Vérification dimensionnelle} : kg$\cdot$(m$\cdot$s$^{-1}$)$^{2}$ = kg$\cdot$m$^{2}\cdot$s$^{-2}$ = J. Le résultat est bien homogène à une énergie.
\textbf{Contre-exemple} : si l'on avait oublié la conversion, on aurait trouvé $\tfrac{1}{2}\times 150 \times 54^{2} = \num{218700}$ « J », un résultat faux d'un facteur $(3,6)^2 \approx 13$ ! C'est l'erreur la plus fréquente — et la plus sanctionnée — au Baccalauréat.
\end{exoresolu}

\begin{exercice}[À toi de jouer]
\begin{enumerate}
\item Exprimer le henry (H) en unités de base, sachant que la tension aux bornes d'une bobine est $u = L\,\dfrac{di}{dt}$.
\item La centrale hydroélectrique de Song Loulou a une puissance de 384 MW. Exprimer cette puissance en watts, puis en unités de base du SI.
\item Convertir en unités SI : $d = 3{,}2$~cm ; $m = 250$~g ; $C = 4{,}7~\mu$F ; $\lambda = 589$~nm.
\end{enumerate}
\end{exercice}

\begin{corrige}[Exercice : À toi de jouer]
\begin{enumerate}
\item De $u = L\,\dfrac{di}{dt}$ on tire $L = \dfrac{u\,t}{i}$, donc $1~\mathrm{H} = 1~\mathrm{V.s.A^{-1}} = 1~\mathrm{kg.m^{2}.s^{-2}.A^{-2}}$.
\item $384~\mathrm{MW} = \num{3,84e8}$~W, soit $\num{3,84e8}$~kg$\cdot$m$^{2}\cdot$s$^{-3}$.
\item $d = \num{3,2e-2}$~m ; $m = \num{2,50e-1}$~kg ; $C = \num{4,7e-6}$~F ; $\lambda = \num{5,89e-7}$~m.
\end{enumerate}
\end{corrige}

Tu maîtrises désormais le langage universel des unités. Dans la section suivante, nous exploiterons pleinement ce socle avec un outil puissant et redoutablement efficace : l'\emph{analyse dimensionnelle}, qui permet de vérifier l'homogénéité de n'importe quelle relation physique — et parfois même de la deviner.

\coursec{Analyse dimensionnelle et homogénéité}

Après avoir maîtrisé le Système international d'unités et les règles d'expression d'un résultat de mesure, nous disposons maintenant du langage universel des physiciens. Mais au-delà des unités, chaque grandeur possède une \textbf{nature intrinsèque} : une longueur reste une longueur, qu'on la mesure en mètres, en kilomètres ou en années-lumière. L'analyse dimensionnelle est l'outil qui permet de manipuler ces natures pour vérifier la cohérence d'une formule, en retrouver une oubliée ou même deviner la forme d'une loi physique avant même de la démontrer. C'est le \emph{premier filtre} qu'un scientifique applique à toute équation.

\courssub{Dimension d'une grandeur physique}

\begin{definition}[Dimension d'une grandeur]
La \textbf{dimension} d'une grandeur physique $G$, notée $[G]$, exprime sa nature physique en fonction des \textbf{dimensions de base} du système d'unités. Elle ne dépend \emph{pas} de l'unité choisie, mais uniquement de la définition de la grandeur.
\end{definition}

Dans le Système international (SI), il y a sept dimensions de base. Nous en utiliserons principalement cinq en Terminale C, plus deux utiles pour la chimie et la thermodynamique :

\begin{popfigure}
\centering
\begin{tabularx}{\linewidth}{|l|c|X|c|}\hline
\rowcolor{popblueL}
\textbf{Dimension de base} & \textbf{Symbole} & \textbf{Grandeur physique associée} & \textbf{Unité SI de base} \\ \hline
Longueur & $[L]$ & Longueur, distance, rayon, position & mètre (\si{\meter}) \\ \hline
Masse & $[M]$ & Masse, inertie & kilogramme (\si{\kilogram}) \\ \hline
Temps & $[T]$ & Durée, période, temps & seconde (\si{\second}) \\ \hline
Intensité électrique & $[I]$ & Courant électrique & ampère (\si{\ampere}) \\ \hline
Température thermodynamique & $[\Theta]$ & Température & kelvin (\si{\kelvin}) \\ \hline
Quantité de matière & $[N]$ & Nombre de moles & mole (\si{\mole}) \\ \hline
Intensité lumineuse & $[J]$ & Flux lumineux & candela (\si{\candela}) \\ \hline
\end{tabularx}
\legende{Les sept dimensions de base du SI. En Terminale C, $[L]$, $[M]$, $[T]$, $[I]$, $[\Theta]$ sont les plus courantes.}
\end{popfigure}

\paragraph{Règles de calcul des dimensions.}
Les dimensions se manipulent comme des grandeurs algébriques (produit, quotient, puissance). Une constante numérique pure (comme $2$, $\pi$, $\frac{1}{2}$, $\sqrt{2}$) est \textbf{sans dimension} : $[k] = 1$.

\begin{propriete}[Calcul des dimensions]
Soient $A$ et $B$ deux grandeurs physiques de dimensions $[A]$ et $[B]$, et $n$ un nombre réel (sans dimension).
\begin{itemize}
    \item Produit : $[A \times B] = [A] \cdot [B]$
    \item Quotient : $\left[\dfrac{A}{B}\right] = \dfrac{[A]}{[B]}$
    \item Puissance : $[A^n] = [A]^n$
\end{itemize}
\end{propriete}

\courssub{Équations aux dimensions des grandeurs du programme}

Toute grandeur dérivée s'exprime comme une combinaison de puissances des dimensions de base. C'est son \textbf{équation aux dimensions}. La connaître par cœur pour les grandeurs du programme est un atout majeur pour le Bac.

\begin{popfigure}
\centering
\begin{tabularx}{\linewidth}{|l|c|l|X|}\hline
\rowcolor{popblueL}
\textbf{Grandeur} & \textbf{Symbole} & \textbf{Équation aux dimensions} & \textbf{Unité SI dérivée (rappel)} \\ \hline
Vitesse & $v$ & $[L][T]^{-1}$ & \si{\meter\per\second} \\ \hline
Accélération & $a$, $\vec{g}$ & $[L][T]^{-2}$ & \si{\meter\per\second\squared} \\ \hline
Force & $F$, $P$ & $[M][L][T]^{-2}$ & newton (\si{\newton}) \\ \hline
Énergie, Travail, Chaleur & $E$, $W$, $Q$ & $[M][L]^2[T]^{-2}$ & joule (\si{\joule}) \\ \hline
Puissance & $P$ & $[M][L]^2[T]^{-3}$ & watt (\si{\watt}) \\ \hline
Tension (Potentiel) & $U$, $V$ & $[M][L]^2[T]^{-3}[I]^{-1}$ & volt (\si{\volt}) \\ \hline
Champ électrique & $\vec{E}$ & $[M][L][T]^{-3}[I]^{-1}$ & \si{\volt\per\meter} = \si{\newton\per\coulomb} \\ \hline
Champ de gravitation & $\vec{g}$ & $[L][T]^{-2}$ & \si{\newton\per\kilogram} = \si{\meter\per\second\squared} \\ \hline
Capacité & $C$ & $[M]^{-1}[L]^{-2}[T]^{4}[I]^{2}$ & farad (\si{\farad}) \\ \hline
Inductance & $L$ & $[M][L]^2[T]^{-2}[I]^{-2}$ & henry (\si{\henry}) \\ \hline
Impédance & $Z$ & $[M][L]^2[T]^{-3}[I]^{-2}$ & ohm (\si{\ohm}) \\ \hline
Activité radioactive & $A$ & $[T]^{-1}$ & becquerel (\si{\becquerel}) \\ \hline
\end{tabularx}
\legende{Récapitulatif des dimensions des grandeurs du programme de Terminale C. Vérifiez chaque ligne à partir des définitions (ex: $U = E/q \Rightarrow [U] = [E]/[Q] = [M][L]^2[T]^{-2}/[I][T] = [M][L]^2[T]^{-3}[I]^{-1}$).}
\end{popfigure}

\begin{exemplebox}[Construction de l'équation aux dimensions de la capacité]
Un condensateur se définit par $C = \dfrac{Q}{U}$.
On connaît $[Q] = [I][T]$ (charge = courant $\times$ temps) et $[U] = [M][L]^2[T]^{-3}[I]^{-1}$.
Donc $[C] = \dfrac{[I][T]}{[M][L]^2[T]^{-3}[I]^{-1}} = [M]^{-1}[L]^{-2}[T]^{4}[I]^{2}$.
C'est bien l'équation du farad (\si{\farad} = \si{\coulomb\per\volt}).
\end{exemplebox}

\courssub{Homogénéité d'une relation physique}

C'est le principe fondateur de l'analyse dimensionnelle.

\begin{propriete}[Principe d'homogénéité dimensionnelle]
Une relation physique correcte est \textbf{nécessairement homogène} : 
\begin{enumerate}
    \item Les deux membres d'une égalité ont \textbf{la même dimension}.
    \item On ne peut \textbf{additionner ou soustraire} que des grandeurs de \textbf{même dimension}.
    \item Les arguments des fonctions transcendantes (log, exp, sin, cos, etc.) sont \textbf{sans dimension}.
\end{enumerate}
\emph{Attention :} L'homogénéité est une condition \textbf{nécessaire} mais \textbf{non suffisante}. Une relation homogène peut être fausse (facteur numérique manquant, mauvais exposant compensé par un autre).
\end{propriete}

\begin{methode}[Vérifier l'homogénéité en 3 étapes]
\begin{enumerate}
    \item \textbf{Écrire} l'équation aux dimensions de chaque terme (membre de gauche et membre de droite, ou chaque terme d'une somme).
    \item \textbf{Simplifier} en utilisant les règles de calcul sur les exposants ($[L]^a [L]^b = [L]^{a+b}$).
    \item \textbf{Comparer} : les dimensions finales doivent être \emph{identiques} terme à terme.
\end{enumerate}
\end{methode}

\begin{exemplebox}[Vérification de $E = \frac{1}{2}mv^2$]
On veut vérifier que cette formule donne bien une énergie.
\begin{itemize}
    \item Membre de gauche : $[E] = [M][L]^2[T]^{-2}$.
    \item Membre de droite : $\left[\frac{1}{2}mv^2\right] = [m] \cdot [v]^2 = [M] \cdot \left([L][T]^{-1}\right)^2 = [M][L]^2[T]^{-2}$.
\end{itemize}
Les deux membres ont la même dimension : \textbf{la relation est homogène} ✓.
\end{exemplebox}

\begin{attention}[Piège classique : l'homogénéité ne prouve pas l'exactitude]
La relation $E = mv^2$ est \textbf{aussi} homogène ($[mv^2] = [M][L]^2[T]^{-2}$), mais elle est \textbf{fausse} (facteur $1/2$ manquant).
La relation $T = 2\pi \sqrt{\frac{m}{k}}$ (période ressort) est homogène. La relation $T = \sqrt{\frac{m}{k}}$ est \textbf{aussi} homogène mais fausse (facteur $2\pi$ manquant).
\textbf{Jamais} on n'écrit : « La formule est homogène, donc elle est juste. » On écrit : « La formule est homogène, elle \emph{peut} être juste. »
\end{attention}

\courssub{Applications : retrouver une loi physique (Pendule simple)}

L'analyse dimensionnelle permet de déterminer la \textbf{forme} d'une loi (les exposants) en supposant que la grandeur cherchée dépend de certaines variables. La constante sans dimension ($k$, $2\pi$, etc.) reste indéterminée.

\begin{experience}[Recherche guidée de la période du pendule simple]
\paragraph{Situation.} Un pendule simple (masse $m$, longueur $\ell$, pesanteur $g$) oscille avec une période $T$. On suppose que $T$ ne dépend \emph{que} de $m$, $\ell$ et $g$ (pas d'amplitude pour de petits angles, pas de frottements).
\paragraph{Modélisation.} On cherche une relation de la forme $T = k \cdot m^{\alpha} \cdot \ell^{\beta} \cdot g^{\gamma}$ avec $k$ sans dimension.
\paragraph{Résolution par les dimensions.}
\[
[T] = [M]^{\alpha} \cdot [L]^{\beta} \cdot \left([L][T]^{-2}\right)^{\gamma} = [M]^{\alpha} [L]^{\beta+\gamma} [T]^{-2\gamma}
\]
On identifie les exposants de chaque dimension de base :
\begin{align*}
    [M] &: \quad 0 = \alpha \\
    [L] &: \quad 0 = \beta + \gamma \\
    [T] &: \quad 1 = -2\gamma
\end{align*}
\paragraph{Résultat.}
$\gamma = -\frac{1}{2}$, $\beta = \frac{1}{2}$, $\alpha = 0$.
La masse \textbf{disparaît} ($\alpha=0$) : la période ne dépend pas de la masse !
\[
T = k \cdot \ell^{1/2} \cdot g^{-1/2} = k \sqrt{\frac{\ell}{g}}
\]
L'expérience ou la résolution de l'équation différentielle donne $k = 2\pi$.
\end{experience}

\begin{popfigure}
\centering
\begin{tikzpicture}[scale=1.1]
    % Pivot
    \fill[popdark] (0,0) circle (3pt);
    \draw[thick, popdark] (-0.5,0.1) -- (0.5,0.1);
    \draw[thick, popdark] (-0.3,0.1) -- (-0.3,0.4) -- (0.3,0.4) -- (0.3,0.1);
    % Fil
    \draw[thick, popmuted] (0,0) -- (2.5,-3.5);
    % Masse
    \shade[ball color=popblue] (2.5,-3.5) circle (0.35);
    % Vecteurs
    \draw[-{Stealth[length=3mm]}, thick, poporange] (2.5,-3.5) -- (2.5,-4.5) node[below right] {$\vec{P}=m\vec{g}$};
    \draw[-{Stealth[length=3mm]}, thick, popgreen] (2.5,-3.5) -- (1.2,-2.5) node[above left] {$\vec{T}$ (tension)};
    % Cotes
    \draw[|<->|, thin, popdark] (0.2,0) -- (2.7,-3.3) node[midway, right=2mm] {$\ell$};
    \draw[dashed, thin, popmuted] (0,-3.5) -- (2.5,-3.5);
    \draw[dashed, thin, popmuted] (0,0) -- (0,-3.5);
    % Angle
    \draw[->, thin, poppurple] (0,-0.5) arc (270:315:0.5) node[midway, right=1mm] {$\theta$};
    % Légende g
    \node[poporange] at (-1.5, -2) {$\vec{g}$};
\end{tikzpicture}
\legende{Pendule simple : forces en jeu ($\vec{P}$, $\vec{T}$) et géométrie ($\ell$, $\theta$). L'analyse dimensionnelle montre que $T$ ne dépend pas de $m$ ni de $\theta$ (pour petits angles), mais seulement de $\sqrt{\ell/g}$.}
\end{popfigure}

\courssub{Limites de l'analyse dimensionnelle}

\begin{aretenir}[Ce que l'analyse dimensionnelle NE PEUT PAS faire]
\begin{itemize}
    \item Déterminer les \textbf{constantes numériques sans dimension} ($2$, $\pi$, $\frac{1}{2}$, $2\pi$, $\sqrt{2}$, etc.).
    \item Différencier des grandeurs de \textbf{même dimension} (ex: travail et couple moteur ont la même dimension $[M][L]^2[T]^{-2}$ mais n'ont pas le même sens physique ; énergie et quantité de chaleur).
    \item Valider un modèle physique complet (elle ne remplace pas l'expérience ni la théorie).
    \item Gérer les sommes de termes de même dimension mais de nature différente (ex: $E = E_p + E_c$ est homogène, mais $E = E_p + F$ ne l'est pas).
\end{itemize}
\end{aretenir}

\begin{savaistu}[Le nombre de Pi de Buckingham]
En modélisation avancée (mécanique des fluides, similitude), le \textbf{théorème $\Pi$ de Buckingham} généralise cette méthode : si un phénomène met en jeu $n$ grandeurs physiques et $k$ dimensions de base indépendantes, la relation peut s'écrire sous la forme d'une fonction de $p = n-k$ nombres \textbf{sans dimension} (les nombres $\Pi$). C'est la base des essais en soufflerie (avions, barrages comme \textbf{Nachtigal} ou \textbf{Lom Pangar} testés sur maquette réduite).
\end{savaistu}

\courssub{Exercice résolu : Vérification de l'impédance d'un circuit RLC série}

\begin{exoresolu}[Impédance complexe d'un circuit RLC série]
\textbf{Énoncé.} On considère un circuit série composé d'une résistance $R$, d'une inductance $L$ et d'une capacité $C$, soumis à une tension sinusoïdale de pulsation $\omega$. L'impédance complexe est donnée par :
\[
\underline{Z} = R + j\left(\omega L - \frac{1}{\omega C}\right)
\]
Vérifier l'homogénéité dimensionnelle de cette expression. Rappel : $j$ est un nombre complexe sans dimension ($j^2 = -1$).

\textbf{Solution détaillée.}
\paragraph{1. Dimensions des grandeurs connues.}
\begin{itemize}
    \item $[R] = [Z] = [M][L]^2[T]^{-3}[I]^{-2}$ (\si{\ohm}).
    \item $[L] = [M][L]^2[T]^{-2}[I]^{-2}$ (\si{\henry}).
    \item $[C] = [M]^{-1}[L]^{-2}[T]^{4}[I]^{2}$ (\si{\farad}).
    \item $[\omega] = [T]^{-1}$ (pulsation, rad/s, radian sans dimension).
    \item $[j] = 1$ (sans dimension).
\end{itemize}

\paragraph{2. Analyse du terme $\omega L$.}
\[
[\omega L] = [T]^{-1} \cdot [M][L]^2[T]^{-2}[I]^{-2} = [M][L]^2[T]^{-3}[I]^{-2}
\]
C'est bien la dimension d'une impédance (ou résistance). ✓

\paragraph{3. Analyse du terme $\frac{1}{\omega C}$.}
\[
\left[\frac{1}{\omega C}\right] = \frac{1}{[T]^{-1} \cdot [M]^{-1}[L]^{-2}[T]^{4}[I]^{2}} = \frac{1}{[M]^{-1}[L]^{-2}[T]^{3}[I]^{2}} = [M][L]^2[T]^{-3}[I]^{-2}
\]
C'est aussi la dimension d'une impédance. ✓

\paragraph{4. Analyse de la somme.}
Les trois termes $R$, $j\omega L$ et $-j\frac{1}{\omega C}$ ont \textbf{exactement la même dimension} : $[M][L]^2[T]^{-3}[I]^{-2}$.
L'addition est donc autorisée dimensionnellement.

\paragraph{Conclusion.}
L'expression $\underline{Z} = R + j\left(\omega L - \frac{1}{\omega C}\right)$ est \textbf{homogène}. Cela ne prouve pas qu'elle est la bonne formule (bien qu'elle le soit), mais elle élimine déjà des erreurs grossières comme $R + \omega L$ (si on oublie le $j$) ou $R + \frac{1}{C}$ (oubli de $\omega$).
\end{exoresolu}

\begin{exercice}[Entraînement : Formule de la fréquence propre d'un circuit LC]
On considère un circuit LC oscillant (pas de résistance). La fréquence propre $f_0$ dépend de l'inductance $L$ et de la capacité $C$.
\begin{enumerate}
    \item Par analyse dimensionnelle, déterminer la forme de la relation $f_0 = k \cdot L^{\alpha} \cdot C^{\beta}$.
    \item Retrouver la valeur de la constante $k$ (rappel de cours ou recherche documentaire).
\end{enumerate}

\end{exercice}

\begin{corrige}[Exercice : Fréquence propre LC]
\paragraph{1. Analyse dimensionnelle.}
$[f_0] = [T]^{-1}$.
$[L] = [M][L]^2[T]^{-2}[I]^{-2}$.
$[C] = [M]^{-1}[L]^{-2}[T]^{4}[I]^{2}$.
On écrit : $[T]^{-1} = \left([M][L]^2[T]^{-2}[I]^{-2}\right)^{\alpha} \left([M]^{-1}[L]^{-2}[T]^{4}[I]^{2}\right)^{\beta}$.
On regroupe :
$[M]^{\alpha-\beta} [L]^{2\alpha-2\beta} [T]^{-2\alpha+4\beta} [I]^{-2\alpha+2\beta}$.
Identification :
\[\begin{cases}
\alpha - \beta = 0 \Rightarrow \alpha = \beta \\
-2\alpha + 4\beta = -1
\end{cases}\]
$\Rightarrow -2\alpha + 4\alpha = -1 \Rightarrow 2\alpha = -1 \Rightarrow \alpha = -\frac{1}{2}$, $\beta = -\frac{1}{2}$.
Donc $f_0 = k \cdot L^{-1/2} C^{-1/2} = \dfrac{k}{\sqrt{LC}}$.
\paragraph{2. Constante $k$.}
Le cours (ou la résolution de l'équation différentielle $L\frac{d^2q}{dt^2} + \frac{q}{C} = 0$) donne $\omega_0 = \frac{1}{\sqrt{LC}}$.
Comme $f_0 = \frac{\omega_0}{2\pi}$, on a $f_0 = \dfrac{1}{2\pi\sqrt{LC}}$.
Donc $k = \dfrac{1}{2\pi}$ (constante sans dimension, indéterminable par l'analyse dimensionnelle seule).
\end{corrige}

\newpage

\coursec{Activité d'intégration}

\coursec{Activité d'intégration : Fiabiliser une mesure de câble sur le pont du Wouri}

\courssub{Objectifs de l'activité}

À l'issue de cette activité d'intégration, tu seras capable de :

\begin{itemize}
\item traiter une \cle{série de mesures} réelle (moyenne, écart-type, incertitude-type) et la comparer à une \cle{mesure unique} ;
\item exprimer correctement un résultat sous la forme $x = (\bar{x} \pm u)$ avec les \cle{chiffres significatifs} adaptés et l'\cle{incertitude relative} ;
\item vérifier l'\cle{homogénéité} d'une relation physique par \cle{analyse dimensionnelle} ;
\item comparer deux résultats expérimentaux et rédiger une conclusion argumentée, comme dans un véritable rapport de contrôle technique.
\end{itemize}

Cette activité mobilise l'ensemble des savoir-faire du module et te place dans la peau d'un technicien chargé de la sécurité d'un ouvrage d'art : la rigueur de tes calculs a ici une conséquence concrète.

\courssub{Situation}

\textbf{Le pont sur le Wouri : fiabiliser une mesure de câble.}

Le pont sur le Wouri, à Douala, est un ouvrage stratégique : il relie les deux rives du fleuve et supporte chaque jour le trafic de dizaines de milliers de véhicules et de trains. Comme tout pont haubané, sa sécurité dépend de la \cle{tension} de ses câbles, qui doit être vérifiée périodiquement lors des campagnes d'entretien.

Plutôt que de mesurer directement la tension (ce qui exigerait de démonter le câble), les ingénieurs utilisent une méthode \emph{vibratoire} : on fait vibrer le câble, on mesure la \cle{période} $T$ de ses oscillations fondamentales, et on en déduit la tension $F$ grâce à la relation physique pour une corde fixe aux deux extrémités :
\[ F = \dfrac{4\,\mu\,L^{2}}{T^{2}} \]
où $\mu$ est la masse linéique du câble et $L$ sa longueur vibrante.

Tu fais partie de l'équipe de contrôle. Voici les données dont tu disposes.

\textbf{Données :}
\begin{itemize}
\item Longueur vibrante du câble, mesurée au décamètre : $L = (48{,}25 \pm 0{,}05)\ \si{m}$ ;
\item Masse linéique du câble (fiche constructeur) : $\mu = 4{,}50\ \si{kg.m^{-1}}$ (valeur exacte pour le calcul) ;
\item Notice du chronomètre : résolution $0{,}01\ \si{s}$, incertitude de mesure unique $\Delta T = \pm 0{,}01\ \si{s}$.
\end{itemize}

\textbf{Série de 12 mesures de la période} (en secondes), relevées au chronomètre par un opérateur :

\begin{center}
\begin{tabular}{|c|cccccccccccc|}\hline
\rowcolor{popblueL} Mesure $n^{\circ}$ & 1 & 2 & 3 & 4 & 5 & 6 & 7 & 8 & 9 & 10 & 11 & 12 \\ \hline
$T$ (s) & 0,42 & 0,45 & 0,41 & 0,43 & 0,44 & 0,40 & 0,42 & 0,46 & 0,43 & 0,41 & 0,44 & 0,42 \\ \hline
\end{tabular}
\end{center}

Un second groupe de contrôle, travaillant sur le même câble avec son propre chronomètre, annonce une tension $F_{B} = (0{,}22 \pm 0{,}02)\ \si{MN}$.

\textbf{Problème :} La tension du câble est-elle correctement déterminée ? Ton résultat est-il compatible avec celui du second groupe ? Rédige la conclusion du rapport de contrôle.

\courssub{Consignes}

\begin{enumerate}
\item Calculer la \cle{moyenne} $\bar{T}$ de la série de 12 mesures, puis l'\cle{écart-type expérimental} $s$ et l'\cle{incertitude-type} $u(\bar{T})$.
\item Comparer $u(\bar{T})$ à l'incertitude $\Delta T = 0{,}01\ \si{s}$ d'une mesure unique. Quel est l'intérêt de répéter les mesures ?
\item Exprimer le résultat sous la forme $T = (\bar{T} \pm u)$ en respectant les règles sur les chiffres significatifs, puis calculer l'\cle{incertitude relative}.
\item Par \cle{analyse dimensionnelle}, vérifier que la relation $F = \dfrac{4\mu L^{2}}{T^{2}}$ est homogène et que $F$ s'exprime bien en newtons.
\item Calculer la tension $F$. En admettant que l'incertitude relative sur $F$ vaut approximativement $\dfrac{\Delta F}{F} \approx \dfrac{\Delta L}{L} + 2\,\dfrac{u}{\bar{T}}$, estimer $\Delta F$ et exprimer $F_{A} = (F \pm \Delta F)$.
\item Les intervalles de tension des groupes A et B se recouvrent-ils ? Les deux résultats sont-ils \cle{compatibles} ?
\item Rédiger la conclusion du rapport de contrôle en cinq lignes maximum.
\end{enumerate}

\courssub{Ressources à mobiliser}

\begin{aretenir}[Boîte à outils de l'activité]
\begin{itemize}
\item Moyenne d'une série de $n$ mesures : $\bar{T} = \dfrac{1}{n}\displaystyle\sum_{i=1}^{n} T_i$.
\item Écart-type expérimental : $s = \sqrt{\dfrac{1}{n-1}\displaystyle\sum_{i=1}^{n}(T_i - \bar{T})^{2}}$.
\item Incertitude-type sur la moyenne : $u(\bar{T}) = \dfrac{s}{\sqrt{n}}$ : répéter les mesures \emph{réduit} l'incertitude.
\item Résultat : $T = (\bar{T} \pm u)$ ; l'incertitude est arrondie à \textbf{un} chiffre significatif (par excès), la moyenne est arrondie au même rang décimal.
\item Incertitude relative : $\dfrac{u}{\bar{T}}$, souvent exprimée en $\%$.
\item Homogénéité : une relation n'est valide que si les deux membres ont la \emph{même dimension} ; $[F] = [M][L][T]^{-2}$ pour une force.
\item Compatibilité de deux mesures : les intervalles $[\bar{x}-u\,;\,\bar{x}+u]$ doivent se \emph{recouper}.
\end{itemize}
\end{aretenir}

\courssub{Démarche et corrigé commenté}

\begin{exoresolu}[Corrigé complet du rapport de contrôle]
Énoncé : traiter la série de 12 mesures de période, exprimer le résultat, valider la formule de tension, calculer $F$ et conclure sur la compatibilité avec le groupe B.
\tcblower
\textbf{Étape 1 : moyenne, écart-type, incertitude-type.}

La somme des 12 mesures vaut $\sum T_i = 5{,}13\ \si{s}$, d'où :
\[ \bar{T} = \dfrac{5{,}13}{12} = 0{,}4275\ \si{s}. \]

On calcule ensuite les écarts $(T_i - \bar{T})$ ; la somme de leurs carrés vaut :
\[ \sum_{i=1}^{12}(T_i - \bar{T})^{2} = 3{,}425\times 10^{-3}\ \si{s^{2}}. \]
L'écart-type expérimental est donc :
\[ s = \sqrt{\dfrac{3{,}425\times 10^{-3}}{11}} = \sqrt{3{,}114\times 10^{-4}} \approx 1{,}76\times 10^{-2}\ \si{s}. \]
L'incertitude-type sur la moyenne :
\[ u(\bar{T}) = \dfrac{s}{\sqrt{12}} = \dfrac{0{,}0176}{3{,}464} \approx 5{,}09\times 10^{-3}\ \si{s}. \]

\textbf{Étape 2 : intérêt de la répétition.}

Pour une mesure unique, l'incertitude du chronomètre est $\Delta T = 0{,}01\ \si{s}$. En répétant 12 fois la mesure et en moyennant, l'incertitude tombe à $u \approx 0{,}005\ \si{s}$, soit une division par environ $2$ (facteur $\sqrt{12} \approx 3,5$ sur l'écart-type, mais l'incertitude instrumentale reste un plancher). La dispersion aléatoire (temps de réaction de l'opérateur) se compense dans la moyenne. La répétition \emph{améliore la fidélité} du résultat.

\textbf{Étape 3 : expression du résultat.}

On arrondit l'incertitude à un chiffre significatif (par excès) : $u = 0{,}005\ \si{s}$. La moyenne est alors arrondie au même rang décimal (millième) : $\bar{T} = 0{,}428\ \si{s}$. D'où :
\[ \boxed{T = (0{,}428 \pm 0{,}005)\ \si{s}} \qquad \dfrac{u}{\bar{T}} = \dfrac{0{,}005}{0{,}428} \approx 1{,}2\ \%. \]

\textbf{Étape 4 : homogénéité de la relation.}

$\mu$ est une masse linéique : $[\mu] = [M][L]^{-1}$. $L$ est une longueur : $[L^{2}] = [L]^{2}$. $T$ est une durée : $[T^{2}] = [T]^{2}$. Le facteur 4 est sans dimension. Donc :
\[ \left[\dfrac{4\mu L^{2}}{T^{2}}\right] = \dfrac{[M][L]^{-1}[L]^{2}}{[T]^{2}} = [M][L][T]^{-2}. \]
Or une force vérifie $[F] = [M][L][T]^{-2}$ (d'après $F = ma$). La relation est donc \textbf{homogène} et $F$ s'exprime bien en newtons (SI). \emph{L'analyse dimensionnelle valide la formule du bureau d'études.}

\textbf{Étape 5 : calcul de la tension.}

Avec $\mu = 4{,}50\ \si{kg.m^{-1}}$, $L = 48{,}25\ \si{m}$ et $T = 0{,}428\ \si{s}$ :
\begin{align*}
F &= \dfrac{4 \times 4{,}50 \times (48{,}25)^{2}}{(0{,}428)^{2}} \\
  &= \dfrac{4 \times 4{,}50 \times 2\,328{,}06}{0{,}18318} \\
  &\approx 2{,}288 \times 10^{5}\ \si{N} = 0{,}229\ \si{MN}.
\end{align*}
Incertitude relative (formule du 1er ordre pour puissances et produits) :
\[ \dfrac{\Delta F}{F} \approx \dfrac{\Delta L}{L} + 2\,\dfrac{u}{\bar{T}} = \dfrac{0{,}05}{48{,}25} + 2\times\dfrac{0{,}005}{0{,}428} \approx 0{,}0010 + 0{,}0234 \approx 0{,}0244 = 2{,}44\ \%. \]
Donc $\Delta F \approx 0{,}0244 \times 0{,}229 \approx 0{,}0056\ \si{MN}$, que l'on arrondit à $0{,}006\ \si{MN}$ (un chiffre significatif). On exprime $F$ au même rang (millième de MN) :
\[ \boxed{F_{A} = (0{,}229 \pm 0{,}006)\ \si{MN}}. \]

\textbf{Étape 6 : compatibilité.}

Intervalle du groupe A (à $1\sigma$) : $[0{,}223\ ;\ 0{,}235]\ \si{MN}$.
Intervalle du groupe B (donné) : $[0{,}20\ ;\ 0{,}24]\ \si{MN}$.
Ces intervalles se \textbf{recouvrent} (par exemple autour de $0{,}225\ \si{MN}$). Les deux résultats sont donc \textbf{compatibles} au sens statistique.

\textbf{Étape 7 : conclusion du rapport (modèle de rédaction).}

\emph{« La période de vibration du câble a été déterminée par une série de 12 mesures : $T = (0{,}428 \pm 0{,}005)\ \si{s}$, soit une incertitude relative de 1,2 \%, deux fois plus faible que celle d'une mesure unique. La relation vibratoire $F = 4\mu L^{2}/T^{2}$, vérifiée homogène, conduit à $F_{A} = (0{,}229 \pm 0{,}006)\ \si{MN}$. Ce résultat étant compatible avec celui du groupe B $(0{,}22 \pm 0{,}02)\ \si{MN}$, la tension du câble est validée dans la plage de tolérance. Aucune intervention corrective n'est requise à ce stade. »}
\end{exoresolu}

\begin{attention}[Pièges fréquents rencontrés dans cette activité]
\begin{itemize}
\item Diviser par $n$ au lieu de $n-1$ dans l'écart-type expérimental (biais sur l'estimation de la variance).
\item Oublier le facteur $\sqrt{n}$ : $u(\bar{T}) = s/\sqrt{n}$, et non $s$.
\item Écrire $T = (0{,}4275 \pm 0{,}005)\ \si{s}$ : la moyenne ne doit pas avoir plus de décimales que l'incertitude.
\item Confondre $\mu$ (masse linéique) et $m$ (masse totale) dans l'analyse dimensionnelle : $[\mu] = [M][L]^{-1}$, $[m] = [M]$.
\item Négliger le facteur 2 devant $\frac{u}{\bar{T}}$ dans la propagation des incertitudes pour $T^2$ au dénominateur.
\item Déclarer deux résultats compatibles sans tracer ou comparer explicitement les intervalles d'incertitude.
\end{itemize}
\end{attention}

\courssub{Bilan de l'activité}

Cette activité t'a fait travailler \emph{comme un scientifique} sur une situation d'ingénierie authentique. Elle a permis de mettre en évidence que :

\begin{itemize}
\item la \cle{répétition des mesures} réduit l'incertitude-type d'un facteur $\sqrt{n}$ : c'est la justification quantitative des protocoles de contrôle qualité industriel ;
\item l'\cle{expression d'un résultat} $(\bar{x} \pm u)$, avec ses chiffres significatifs et son incertitude relative, est un langage universel : sans elle, aucune comparaison entre mesures n'est possible ;
\item l'\cle{analyse dimensionnelle} est un outil de \emph{validation} puissant : elle confirme la cohérence physique d'une formule avant tout calcul numérique ;
\item deux mesures ne sont \cle{compatibles} que si leurs intervalles d'incertitude se recouvrent : critère décisif pour la rédaction d'un rapport technique et la prise de décision.
\end{itemize}

Ces compétences — estimer, exprimer, vérifier, comparer — seront réinvesties dans toutes les leçons de physique de l'année, chaque fois qu'un résultat expérimental devra être exploité avec rigueur.

\newpage

\coursec{Méthodes et exercices résolus}

\courssub{Méthodes et exercices résolus}

\begin{methode}[Estimer l'incertitude d'une mesure unique (Type B)]
    \textbf{Objectif :} Déterminer l'incertitude standard $u(x)$ associée à une \emph{seule} lecture d'un instrument.
    \begin{enumerate}
        \item \textbf{Identifier l'instrument} et sa \textbf{plus petite division} (résolution) $\delta$ (ex : $\delta = \SI{1}{\mm}$ pour une règle graduée, $\delta = \SI{0,01}{\mm}$ pour un pied à coulisse).
        \item \textbf{Estimer l'erreur de lecture} (appréciation de l'opérateur). En l'absence d'information précise, on retient classiquement :
              \[ u_{\text{lecture}} = \frac{\delta}{2} \quad \text{(règle, pied à coulisse)} \qquad \text{ou} \qquad u_{\text{lecture}} = \frac{\delta}{\sqrt{12}} \approx \frac{\delta}{3,5} \quad \text{(répartition uniforme supposée).} \]
        \item \textbf{Consulter la notice du constructeur} (classe de précision, erreur maximale permise $E_{\text{max}}$). Si donnée, $u_{\text{constructeur}} = \frac{E_{\text{max}}}{\sqrt{3}}$ (loi uniforme) ou $\frac{E_{\text{max}}}{2}$ (loi triangulaire).
        \item \textbf{Combiner les composantes} (si plusieurs sources indépendantes) : $u_c = \sqrt{u_1^2 + u_2^2 + \dots}$.
        \item \textbf{Arrondir l'incertitude finale} à \textbf{1 ou 2 chiffres significatifs} (généralement 1 au Bac).
    \end{enumerate}
    \textbf{Formule à retenir :} Pour un appareil numérique à affichage direct (multimètre, balance électronique), l'incertitude est souvent donnée par le constructeur sous la forme $\pm (a\% \text{ de la lecture} + b \text{ digits})$.
    \[ u = \sqrt{ \left( \frac{a}{100} \times x \right)^2 + \left( \frac{b \times \text{résolution}}{\sqrt{3}} \right)^2 } \]
\end{methode}

\begin{exoresolu}[Mesure de la tension secteur avec un multimètre (Contexte ENEO)]
    Un technicien ENEO mesure la tension efficace aux bornes d'une prise domestique à Yaoundé à l'aide d'un multimètre numérique de classe \num{0,5}\%. L'affichage indique \SI{232}{\volt}. La notice précise : « Précision : $\pm(\num{0,5}\% \text{ de la lecture} + 2 \text{ digits}) »$. La résolution (dernier digit) est de \SI{0,1}{\volt}.
    \textbf{1.} Calculer l'incertitude standard $u(U)$.
    \textbf{2.} Exprimer le résultat de la mesure selon la norme.
    \tcblower
    \textbf{Solution détaillée}
    \begin{enumerate}
        \item \textbf{Identification des composantes d'incertitude (Type B) :}
              \begin{itemize}
                  \item Composante liée au pourcentage de lecture : $u_1 = \dfrac{\num{0,5}}{100} \times \SI{232}{\volt} = \SI{1,16}{\volt}$.
                  \item Composante liée à la résolution (quantification) : $u_2 = \dfrac{2 \times \SI{0,1}{\volt}}{\sqrt{3}} = \dfrac{\SI{0,2}{\volt}}{1,732} \approx \SI{0,115}{\volt}$.
                        (On divise par $\sqrt{3}$ car l'erreur de quantification suit une loi uniforme sur $\pm 1$ digit).
              \end{itemize}
        \item \textbf{Incertitude composée (sources indépendantes) :}
              \[ u_c(U) = \sqrt{u_1^2 + u_2^2} = \sqrt{\num{1,16}^2 + \num{0,115}^2} = \sqrt{\num{1,3456} + \num{0,0132}} = \sqrt{\num{1,3588}} \approx \SI{1,166}{\volt}. \]
        \item \textbf{Arrondi de l'incertitude :} On retient \textbf{1 chiffre significatif} (convention Bac) : $u(U) = \mathbf{\SI{1}{\volt}}$.
              (Si on gardait 2 chiffres : \SI{1,2}{\volt}).
        \item \textbf{Arrondi de la valeur mesurée :} La valeur doit avoir la même décimale que l'incertitude. \SI{232}{\volt} est déjà un entier.
        \item \textbf{Expression finale :}
              \[ \boxed{U = (\SI{232}{\volt} \pm \SI{1}{\volt})} \quad \text{ou} \quad \boxed{U = \SI{232}{\volt} \pm \SI{1}{\volt}} \]
              \textbf{Incertitude relative :} $\dfrac{u(U)}{|U|} = \dfrac{1}{232} \approx \num{0,43}\%$.
    \end{enumerate}
    \textbf{Conclusion :} La tension secteur mesurée est compatible avec la norme \SI{230}{\volt} $\pm \num{10}\%$ (soit \SI{207}{\volt} à \SI{253}{\volt}).
\end{exoresolu}

\begin{methode}[Traiter une série de mesures : moyenne, écart-type, incertitude-type (Type A)]
    \textbf{Objectif :} Déterminer la meilleure estimation de la valeur vraie et son incertitude à partir de $n$ mesures répétées ($n \ge 5$).
    \begin{enumerate}
        \item \textbf{Calculer la moyenne arithmétique} (meilleure estimation) :
              \[ \bar{x} = \frac{1}{n} \sum_{i=1}^n x_i \]
        \item \textbf{Calculer l'écart-type expérimental (échantillon)} :
              \[ s = \sqrt{ \frac{1}{n-1} \sum_{i=1}^n (x_i - \bar{x})^2 } \]
              \emph{Attention :} Utiliser $n-1$ au dénominateur (estimateur sans biais), pas $n$.
        \item \textbf{Calculer l'incertitude-type de la moyenne (incertitude Type A)} :
              \[ u_A(\bar{x}) = \frac{s}{\sqrt{n}} \]
        \item \textbf{Vérifier les valeurs aberrantes} (test de Chauvenet ou $3\sigma$ simple) avant de calculer. Si une valeur est suspecte, la noter, la justifier, puis décider de la conserver ou non.
        \item \textbf{Arrondir} $u_A$ à 1 ou 2 chiffres significatifs, puis arrondir $\bar{x}$ au même rang de décimale.
    \end{enumerate}
    \textbf{Réflexe Bac :} La calculatrice donne souvent $\sigma_n$ (écart-type population) et $\sigma_{n-1}$ (échantillon). Choisir \textbf{$\sigma_{n-1}$} (souvent noté $s$ ou $S_x$).
\end{methode}

\begin{exoresolu}[Période d'un pendule simple (Labo Physique Lycée de Ngoa-Ekélé)]
    Un élève de Terminale C effectue 10 mesures de la période $T$ (en secondes) d'un pendule simple de longueur $L \approx \SI{1}{\meter}$ à l'aide d'un chronomètre (résolution \SI{0,01}{\second}).
    Mesures : \num{2,01} ; \num{2,03} ; \num{2,00} ; \num{2,04} ; \num{2,02} ; \num{2,05} ; \num{1,99} ; \num{2,03} ; \num{2,01} ; \num{2,02}.
    \textbf{1.} Calculer la moyenne $\bar{T}$ et l'écart-type $s$.
    \textbf{2.} Déterminer l'incertitude-type $u_A(\bar{T})$.
    \textbf{3.} Exprimer le résultat final. L'incertitude de lecture du chronomètre (\SI{0,005}{\second}) est-elle négligeable ?
    \tcblower
    \textbf{Solution détaillée}
    \begin{enumerate}
        \item \textbf{Moyenne :}
              \[ \bar{T} = \frac{1}{10}(\num{2,01}+\num{2,03}+\num{2,00}+\num{2,04}+\num{2,02}+\num{2,05}+\num{1,99}+\num{2,03}+\num{2,01}+\num{2,02}) = \frac{\num{20,20}}{10} = \SI{2,020}{\second}. \]
        \item \textbf{Écart-type expérimental ($s = \sigma_{n-1}$) :}
              \begin{align*}
                  \sum (T_i - \bar{T})^2 &= (-\num{0,01})^2 + (\num{0,01})^2 + (-\num{0,02})^2 + (\num{0,02})^2 + (\num{0,00})^2 \\
                  &\quad + (\num{0,03})^2 + (-\num{0,03})^2 + (\num{0,01})^2 + (-\num{0,01})^2 + (\num{0,00})^2 \\
                  &= \num{0,0001} + \num{0,0001} + \num{0,0004} + \num{0,0004} + \num{0} + \num{0,0009} + \num{0,0009} + \num{0,0001} + \num{0,0001} + \num{0} \\
                  &= \num{0,0030}\,\si{\second\squared}.
              \end{align*}
              \[ s = \sqrt{ \frac{\num{0,0030}}{10-1} } = \sqrt{ \frac{\num{0,0030}}{9} } = \sqrt{\num{0,000333\dots}} \approx \SI{0,0183}{\second}. \]
        \item \textbf{Incertitude-type de la moyenne :}
              \[ u_A(\bar{T}) = \frac{s}{\sqrt{n}} = \frac{\SI{0,0183}{\second}}{\sqrt{10}} = \frac{\SI{0,0183}{\second}}{3,162} \approx \SI{0,00579}{\second}. \]
        \item \textbf{Arrondis :} $u_A \approx \mathbf{\SI{0,006}{\second}}$ (1 chiffre significatif).
              $\bar{T}$ s'arrondit au millième : $\mathbf{\SI{2,020}{\second}}$.
        \item \textbf{Comparaison avec incertitude de lecture (Type B) :}
              $u_B = \frac{\SI{0,01}{\second}}{2} = \SI{0,005}{\second}$.
              $u_A = \SI{0,006}{\second}$.
              Elles sont du même ordre de grandeur. Aucune n'est négligeable. L'incertitude combinée serait $u_c = \sqrt{\num{0,006}^2 + \num{0,005}^2} \approx \SI{0,008}{\second}$.
              \emph{Au Bac, si on ne demande que l'incertitude Type A, on donne $u_A$.}
        \item \textbf{Résultat final (Type A seul) :}
              \[ \boxed{\bar{T} = (\SI{2,020}{\second} \pm \SI{0,006}{\second})} \quad \text{soit} \quad \boxed{T = \SI{2,020}{\second} \pm \SI{0,006}{\second}} \]
    \end{enumerate}
    \textbf{Conclusion :} La dispersion des mesures (réaction de l'élève) domine légèrement sur la résolution de l'appareil.
\end{exoresolu}

\begin{methode}[Exprimer un résultat de mesure : valeur $\pm$ incertitude, chiffres significatifs, incertitude relative]
    \textbf{Objectif :} Présenter le résultat final selon les conventions internationales (NF ISO/CEI Guide 98-3).
    \begin{enumerate}
        \item \textbf{Arrondir l'incertitude $u$} à \textbf{1 chiffre significatif} (par défaut au Bac) ou 2 si le premier chiffre est 1 ou 2 (pour éviter une perte d'information trop brutale).
              \emph{Ex :} \num{0,0123} $\to$ \num{0,01} ; \num{0,00198} $\to$ \num{0,002} ; \num{123} $\to$ \num{100}.
        \item \textbf{Arrondir la valeur estimée $\bar{x}$} au \textbf{même rang de décimale} (même puissance de 10) que l'incertitude arrondie.
              \emph{Ex :} Si $u = \num{0,01}$, $\bar{x} = \num{12,3456} \to \num{12,35}$.
        \item \textbf{Écrire sous la forme} : $\boxed{x = \bar{x} \pm u}$ (unité une seule fois à la fin, ou factorisée : $(\bar{x} \pm u)\,\text{unité}$).
        \item \textbf{Calculer l'incertitude relative} (sans unité) : $\dfrac{u(x)}{|x|}$, souvent en $\%$.
        \item \textbf{Chiffres significatifs de la valeur :} Le nombre de chiffres significatifs de $\bar{x}$ est imposé par l'arrondi de l'étape 2. Ne jamais écrire plus de chiffres que la précision ne le permet.
    \end{enumerate}
    \textbf{Attention :} Ne jamais écrire $\num{10,5} \pm \num{0,123}$ (décimales non alignées) ni $\num{10,50} \pm \num{0,1}$ (chiffre significatif de trop sur la valeur).
\end{methode}

\begin{exoresolu}[Résistance d'un shunt pour compteur prépayé (SONATREL)]
    On mesure la résistance $R$ d'un shunt manganin par la méthode voltampèremétrique. On obtient : $\bar{R} = \SI{0,01234}{\ohm}$ et $u(R) = \SI{0,000187}{\ohm}$.
    \textbf{1.} Exprimer correctement le résultat.
    \textbf{2.} Calculer l'incertitude relative en $\%$.
    \textbf{3.} On utilise ce shunt pour mesurer un courant $I = \SI{50,0}{\ampere} \pm \SI{0,2}{\ampere}$. La tension aux bornes du shunt est $U = \SI{0,617}{\volt} \pm \SI{0,003}{\volt}$. Vérifier la cohérence avec la loi d'Ohm $U = RI$.
    \tcblower
    \textbf{Solution détaillée}
    \begin{enumerate}
        \item \textbf{Arrondi de l'incertitude :} $u(R) = \SI{0,000187}{\ohm}$. Premier chiffre significatif : 1. On garde \textbf{2 chiffres significatifs} (règle du « 1 ou 2 »).
              $u(R) \approx \mathbf{\SI{0,00019}{\ohm}}$ (soit $\num{1,9} \times 10^{-4}\,\si{\ohm}$).
        \item \textbf{Arrondi de la valeur :} Même rang de décimale : $10^{-5}\,\si{\ohm}$ (cent-millionièmes).
              $\bar{R} = \num{0,01234} \to \mathbf{\SI{0,01234}{\ohm}}$ (déjà au bon rang).
              \emph{Note :} Si $u$ était \SI{0,0002}{\ohm}, $\bar{R}$ serait \SI{0,0123}{\ohm}.
        \item \textbf{Expression finale :}
              \[ \boxed{R = (\SI{0,01234}{\ohm} \pm \SI{0,00019}{\ohm})} \quad \text{ou} \quad \boxed{R = \SI{12,34}{\milli\ohm} \pm \SI{0,19}{\milli\ohm}} \]
        \item \textbf{Incertitude relative :}
              \[ \frac{u(R)}{R} = \frac{\num{0,00019}}{\num{0,01234}} \approx \num{0,0154} = \mathbf{\num{1,54}\%} \]
        \item \textbf{Vérification loi d'Ohm (cohérence des intervalles) :}
              Intervalle pour $U$ : $[\SI{0,614}{\volt} ; \SI{0,620}{\volt}]$.
              Intervalle pour $R \times I$ :
              Valeur nominale : $\num{0,01234} \times \num{50,0} = \SI{0,617}{\volt}$.
              Incertitude relative sur le produit : $\frac{u(RI)}{RI} \approx \sqrt{ \left(\frac{u_R}{R}\right)^2 + \left(\frac{u_I}{I}\right)^2 } = \sqrt{\num{0,0154}^2 + \left(\frac{\num{0,2}}{\num{50}}\right)^2 } = \sqrt{\num{0,000237} + \num{0,000016}} \approx \num{0,0159}$.
              $u(RI) = \num{0,0159} \times \num{0,617} \approx \SI{0,0098}{\volt}$.
              Intervalle pour $RI$ : $[\SI{0,607}{\volt} ; \SI{0,627}{\volt}]$.
              Les intervalles $[\SI{0,614}{\volt} ; \SI{0,620}{\volt}]$ et $[\SI{0,607}{\volt} ; \SI{0,627}{\volt}]$ \textbf{se chevauchent largement}. La mesure est cohérente avec la loi d'Ohm.
    \end{enumerate}
\end{exoresolu}

\begin{methode}[Analyse dimensionnelle des grandeurs physiques (Nouvelles grandeurs du programme)]
    \textbf{Objectif :} Exprimer la dimension d'une grandeur en fonction des dimensions de base $[L], [M], [T], [I], [\Theta]$.
    \begin{enumerate}
        \item \textbf{Écrire la définition physique} ou la loi fondamentale liant la grandeur cherchée aux grandeurs de base.
        \item \textbf{Remplacer chaque grandeur par sa dimension} (utiliser les crochets $[\,]$).
        \item \textbf{Simplifier} les puissances (algèbre des exposants).
        \item \textbf{Vérifier} par analyse dimensionnelle que la relation utilisée est homogène.
    \end{enumerate}
    \textbf{Tableau de référence (à connaître par cœur) :}
    \begin{center}
    \begin{tabular}{|c|c|c|}\hline
    \rowcolor{popblueL}
    \textbf{Grandeur} & \textbf{Relation définissante} & \textbf{Dimension} \\ \hline
    Vitesse $v$ & $v = \frac{dx}{dt}$ & $[L][T]^{-1}$ \\
    Accélération $a$ & $a = \frac{dv}{dt}$ & $[L][T]^{-2}$ \\
    Force $F$ & $F = ma$ & $[M][L][T]^{-2}$ \\
    Énergie/Travail $W$ & $W = F \cdot x$ & $[M][L]^2[T]^{-2}$ \\
    Puissance $P$ & $P = \frac{dW}{dt}$ & $[M][L]^2[T]^{-3}$ \\
    Charge $Q$ & $Q = It$ & $[I][T]$ \\
    Tension $U$ & $U = \frac{W}{Q}$ & $[M][L]^2[T]^{-3}[I]^{-1}$ \\
    Résistance $R$ & $R = \frac{U}{I}$ & $[M][L]^2[T]^{-3}[I]^{-2}$ \\
    \textbf{Capacité $C$} & $C = \frac{Q}{U}$ & $\mathbf{[M]^{-1}[L]^{-2}[T]^4[I]^2}$ \\
    \textbf{Inductance $L$} & $U = L\frac{di}{dt} \Rightarrow L = \frac{U\,t}{I}$ & $\mathbf{[M][L]^2[T]^{-2}[I]^{-2}}$ \\
    \textbf{Impédance $Z$} & $Z = \frac{U}{I}$ & $\mathbf{[M][L]^2[T]^{-3}[I]^{-2}}$ (même dim. que $R$) \\
    \textbf{Champ électrique $E$} & $E = \frac{F}{q} = \frac{U}{d}$ & $\mathbf{[M][L][T]^{-3}[I]^{-1}}$ \\
    \textbf{Activité $A$} & $A = \left|\frac{dN}{dt}\right|$ & $\mathbf{[T]^{-1}}$ (Becquerel) \\ \hline
    \end{tabular}
    \end{center}
\end{methode}

\begin{exoresolu}[Dimensions de la permittivité $\varepsilon_0$ et perméabilité $\mu_0$ (Lignes HT SONATREL)]
    La capacité d'un condensateur plan est $C = \varepsilon_0 \frac{S}{d}$. La force entre deux fils parallèles porteurs de courant est $F = \frac{\mu_0}{2\pi} \frac{I_1 I_2}{d} \ell$.
    \textbf{1.} Déterminer les dimensions de $\varepsilon_0$ (permittivité du vide) et $\mu_0$ (perméabilité du vide).
    \textbf{2.} Vérifier que $\frac{1}{\sqrt{\varepsilon_0 \mu_0}}$ a la dimension d'une vitesse (vitesse de la lumière $c$).
    \tcblower
    \textbf{Solution détaillée}
    \begin{enumerate}
        \item \textbf{Permittivité $\varepsilon_0$ :}
              \[ C = \varepsilon_0 \frac{S}{d} \implies \varepsilon_0 = C \frac{d}{S} \]
              $[C] = [M]^{-1}[L]^{-2}[T]^4[I]^2$ (voir tableau).
              $[d] = [L]$, $[S] = [L]^2$.
              \[ [\varepsilon_0] = [M]^{-1}[L]^{-2}[T]^4[I]^2 \times \frac{[L]}{[L]^2} = \mathbf{[M]^{-1}[L]^{-3}[T]^4[I]^2} \]
        \item \textbf{Perméabilité $\mu_0$ :}
              \[ F = \frac{\mu_0}{2\pi} \frac{I_1 I_2}{d} \ell \implies \mu_0 = \frac{F\,d}{I^2\,\ell} \quad (\text{$2\pi$ sans dimension}) \]
              $[F] = [M][L][T]^{-2}$, $[d] = [L]$, $[I] = [I]$, $[\ell] = [L]$.
              \[ [\mu_0] = \frac{[M][L][T]^{-2} \times [L]}{[I]^2 \times [L]} = \mathbf{[M][L][T]^{-2}[I]^{-2}} \]
        \item \textbf{Vérification pour $c = \frac{1}{\sqrt{\varepsilon_0 \mu_0}}$ :}
              \[ [\varepsilon_0 \mu_0] = [M]^{-1}[L]^{-3}[T]^4[I]^2 \times [M][L][T]^{-2}[I]^{-2} = [L]^{-2}[T]^2 \]
              \[ \left[ \sqrt{\varepsilon_0 \mu_0} \right] = [L]^{-1}[T] \]
              \[ \left[ \frac{1}{\sqrt{\varepsilon_0 \mu_0}} \right] = [L][T]^{-1} \]
              C'est bien la dimension d'une \textbf{vitesse}. \quad \checkmark
    \end{enumerate}
    \textbf{Note :} $\varepsilon_0 \approx \num{8,85e-12}\,\si{\farad\per\meter}$, $\mu_0 = 4\pi \times 10^{-7}\,\si{\henry\per\meter}$, $c \approx \num{3,00e8}\,\si{\meter\per\second}$.
\end{exoresolu}

\begin{methode}[Vérifier l'homogénéité d'une relation physique]
    \textbf{Objectif :} S'assurer qu'une équation est physiquement plausible (condition nécessaire, non suffisante).
    \begin{enumerate}
        \item \textbf{Écrire la dimension de CHAQUE terme} de l'équation (membres de gauche et de droite, et chaque terme d'une somme).
        \item \textbf{Rappel :} Les constantes pures ($2, \pi, \sqrt{3}, \ln(2), \dots$) sont \textbf{sans dimension} ($[1]$).
        \item \textbf{Règle :} Dans une égalité $A = B$, on doit avoir $[A] = [B]$. Dans une somme $A = B + C$, on doit avoir $[A] = [B] = [C]$.
        \item \textbf{Conclusion :} Si les dimensions ne correspondent pas, la relation est \textbf{fausse}. Si elles correspondent, la relation est \textbf{homogène} (peut être juste, mais pas prouvée).
    \end{enumerate}
    \textbf{Astuce :} Pour les arguments de fonctions transcendantes ($\sin, \cos, \exp, \ln$), l'argument est \textbf{toujours sans dimension}.
\end{methode}

\begin{exoresolu}[Détection d'erreurs dans des formules du programme (Bac Blanc)]
    Vérifier l'homogénéité des relations suivantes. Si une relation n'est pas homogène, proposer une correction plausible.
    \begin{enumerate}
        \item $T = 2\pi \sqrt{\frac{m}{k}}$ (Période oscillateur masse-ressort)
        \item $E = \frac{1}{2} C U^2$ (Énergie condensateur)
        \item $v = v_0 + a t^2$ (Cinématique)
        \item $\lambda = \frac{h}{p}$ (Longueur d'onde de Broglie, $h$ constante de Planck, $p$ impulsion)
        \item $I = I_0 e^{-\frac{R}{L}t}$ (Courant circuit RL série)
    \end{enumerate}
    \tcblower
    \textbf{Solution détaillée}
    Rappels dimensions : $[T]=[T]$, $[m]=[M]$, $[k]=[M][T]^{-2}$ (raideur $F=kx$), $[C]=[M]^{-1}[L]^{-2}[T]^4[I]^2$, $[U]=[M][L]^2[T]^{-3}[I]^{-1}$, $[v]=[L][T]^{-1}$, $[a]=[L][T]^{-2}$, $[h]=[M][L]^2[T]^{-1}$ (action/quantité de mouvement $\times$ longueur), $[p]=[M][L][T]^{-1}$, $[R]=[M][L]^2[T]^{-3}[I]^{-2}$, $[L]=[M][L]^2[T]^{-2}[I]^{-2}$.
    \begin{enumerate}
        \item $[2\pi]=1$. $\left[\sqrt{\frac{m}{k}}\right] = \sqrt{\frac{[M]}{[M][T]^{-2}}} = \sqrt{[T]^2} = [T]$. \textbf{Homogène.} \checkmark
        \item $[1/2]=1$. $[C U^2] = [M]^{-1}[L]^{-2}[T]^4[I]^2 \times \left([M][L]^2[T]^{-3}[I]^{-1}\right)^2 = [M]^{-1}[L]^{-2}[T]^4[I]^2 \times [M]^2[L]^4[T]^{-6}[I]^{-2} = [M][L]^2[T]^{-2}$. Dimension de l'énergie. \textbf{Homogène.} \checkmark
        \item Terme 1 : $[v_0] = [L][T]^{-1}$. Terme 2 : $[a t^2] = [L][T]^{-2} \times [T]^2 = [L]$.
              \textbf{NON HOMOGÈNE :} On ne peut pas ajouter une vitesse $[L][T]^{-1}$ et une longueur $[L]$.
              \textbf{Correction :} $v = v_0 + a t$ (terme 2 : $[L][T]^{-2}[T] = [L][T]^{-1}$).
        \item $[h] = [M][L]^2[T]^{-1}$. $[p] = [M][L][T]^{-1}$.
              $\left[\frac{h}{p}\right] = \frac{[M][L]^2[T]^{-1}}{[M][L][T]^{-1}} = [L]$. Dimension de longueur. \textbf{Homogène.} \checkmark
        \item Argument de l'exponentielle : $\left[\frac{R}{L}t\right]$.
              $\left[\frac{R}{L}\right] = \frac{[M][L]^2[T]^{-3}[I]^{-2}}{[M][L]^2[T]^{-2}[I]^{-2}} = [T]^{-1}$.
              $\left[\frac{R}{L}t\right] = [T]^{-1}[T] = [1]$ (sans dimension). \textbf{Homogène.} \checkmark
    \end{enumerate}
\end{exoresolu}

\begin{methode}[Valider un modèle / Trouver une loi par analyse dimensionnelle (Méthode de Vaschy-Buckingham simplifiée)]
    \textbf{Objectif :} Déterminer la forme d'une loi physique $Q = f(X, Y, Z, \dots)$ en supposant une relation de type produit de puissances.
    \begin{enumerate}
        \item \textbf{Lister les grandeurs pertinentes} (variables et paramètres) et leurs dimensions.
        \item \textbf{Écrire l'ansatz (hypothèse de forme)} : $Q = k \, X^a Y^b Z^c \dots$ où $k$ est une constante \textbf{sans dimension}.
        \item \textbf{Écrire l'équation aux dimensions} : $[Q] = [X]^a [Y]^b [Z]^c \dots$
        \item \textbf{Résoudre le système linéaire} pour les exposants $a, b, c \dots$ (un système de 3 à 5 équations selon les dimensions de base présentes).
        \item \textbf{Écrire la relation finale} : $Q = k \dots$ Préciser que $k$ ne peut être déterminé que par l'expérience ou la théorie complète.
    \end{enumerate}
    \textbf{Condition d'application :} Le nombre de grandeurs pertinentes doit être supérieur au nombre de dimensions de base indépendantes (sinon système déterminé ou surdéterminé). Les grandeurs doivent être \emph{indépendantes dimensionnellement}.
\end{methode}

\begin{exoresolu}[Fréquence de résonance d'un circuit RLC série (Radio FM Cameroun)]
    On cherche la fréquence de résonance $f_0$ d'un circuit RLC série. On suppose qu'elle dépend de l'inductance $L$, de la capacité $C$ et de la résistance $R$.
    \textbf{1.} Par analyse dimensionnelle, déterminer la relation entre $f_0$, $L$ et $C$ (montrer que $R$ n'intervient pas).
    \textbf{2.} Une mesure donne $f_0 = \SI{100}{\mega\hertz}$ pour $L = \SI{1}{\micro\henry}$. Calculer $C$.
    \tcblower
    \textbf{Solution détaillée}
    \begin{enumerate}
        \item \textbf{Dimensions :} $[f_0] = [T]^{-1}$, $[L] = [M][L]^2[T]^{-2}[I]^{-2}$, $[C] = [M]^{-1}[L]^{-2}[T]^4[I]^2$, $[R] = [M][L]^2[T]^{-3}[I]^{-2}$.
        \item \textbf{Ansatz :} $f_0 = k \, L^a C^b R^c$ ($k$ sans dimension).
        \item \textbf{Équation aux dimensions :}
              \[ [T]^{-1} = \left([M][L]^2[T]^{-2}[I]^{-2}\right)^a \left([M]^{-1}[L]^{-2}[T]^4[I]^2\right)^b \left([M][L]^2[T]^{-3}[I]^{-2}\right)^c \]
              Regroupement des exposants pour chaque dimension de base :
              \begin{align*}
                  [M] &: 0 = a - b + c \quad (1) \\
                  [L] &: 0 = 2a - 2b + 2c \quad (2) \\
                  [T] &: -1 = -2a + 4b - 3c \quad (3) \\
                  [I] &: 0 = -2a + 2b - 2c \quad (4)
              \end{align*}
        \item \textbf{Résolution :}
              (2) est $2\times(1)$ : redondant. (4) est $-2\times(1)$ : redondant.
              Il reste 2 équations indépendantes (1) et (3) pour 3 inconnues ($a, b, c$).
              \emph{Interprétation :} La résistance $R$ n'est pas indépendante dimensionnellement de $L$ et $C$ pour former une fréquence. On ne peut pas déterminer $c$ de façon unique sans hypothèse physique supplémentaire.
              \textbf{Hypothèse physique :} En résonance série, l'impédance est purement résistive ($Z=R$), la fréquence ne dépend \textbf{pas} de $R$. On pose $c=0$.
              Système réduit :
              \[\begin{cases}
                  a - b = 0 \implies a = b \\
                  -2a + 4b = -1 \implies -2a + 4a = -1 \implies 2a = -1 \implies a = -1/2
              \end{cases}\]
              Donc $a = b = -1/2$.
        \item \textbf{Relation trouvée :}
              \[ f_0 = k \, L^{-1/2} C^{-1/2} = \frac{k}{\sqrt{LC}} \]
              La théorie complète donne $k = \frac{1}{2\pi}$.
              \[ \boxed{f_0 = \frac{1}{2\pi\sqrt{LC}}} \]
        \item \textbf{Application numérique :}
              $f_0 = \SI{100}{\mega\hertz} = \num{1,00e8}\,\si{\hertz}$, $L = \SI{1}{\micro\henry} = \num{1,00e-6}\,\si{\henry}$.
              \[ C = \frac{1}{(2\pi f_0)^2 L} = \frac{1}{(2\pi \times \num{1e8})^2 \times \num{1e-6}} = \frac{1}{4\pi^2 \times \num{1e16} \times \num{1e-6}} = \frac{1}{4\pi^2 \times \num{1e10}} \]
              \[ C \approx \frac{1}{\num{39,48} \times \num{1e10}} \approx \num{2,53e-12}\,\si{\farad} = \boxed{\SI{2,53}{\pico\farad}} \]
    \end{enumerate}
\end{exoresolu}

\begin{attention}[Les 5 erreurs qui coûtent des points au Bac]
    \begin{enumerate}
        \item \textbf{Confondre écart-type ($s$) et incertitude-type ($u_A$).}
              Écrire $u = s$ au lieu de $u = s/\sqrt{n}$. Perte de 2 à 3 points. \emph{Règle :} L'incertitude sur la \emph{moyenne} diminue avec $\sqrt{n}$.
        \item \textbf{Mauvais arrondi de l'incertitude (chiffres significatifs).}
              Donner $u = \SI{0,0123}{\volt}$ (3 ch. sig.) ou $u = \SI{0,1}{\volt}$ pour une valeur \SI{1,23}{\volt} (incohérence décimale).
              \emph{Règle :} 1 ou 2 ch. sig. sur $u$ ; $\bar{x}$ aligné sur la même décimale.
        \item \textbf{Oublier l'unité dans le résultat final} ou l'écrire deux fois : $x = \SI{10,5}{\meter} \pm \SI{0,2}{\meter}$ (correct) vs $x = \SI{10,5}{\meter} \pm \SI{0,2}{\meter}$ (lourd mais accepté) vs $x = \num{10,5} \pm \num{0,2}$ (zéro point).
        \item \textbf{Erreur d'homogénéité : additionner des grandeurs de dimensions différentes.}
              Ex : $v = v_0 + at^2$, $x = vt + \frac{1}{2}at$ (longueur + vitesse). Vérifier \textbf{systématiquement} les dimensions de chaque terme d'une somme.
        \item \textbf{Traiter une constante ($2, \pi, k$) comme ayant une dimension.}
              Dans $T = 2\pi\sqrt{L/g}$, $[2\pi]=1$. Ne jamais écrire $[2\pi]=[T]$. De même, l'argument de $\exp$, $\ln$, $\sin$ est \textbf{toujours sans dimension}. Vérifier l'exposant : $[-\lambda t] = [1]$.
    \end{enumerate}
\end{attention}

\newpage

\coursec{Exercices}

\coursec{Mesures, incertitudes et analyse dimensionnelle}

\courssub{Introduction : travailler comme un scientifique}
En physique, toute mesure est entachée d'imprécisions. Un résultat scientifique ne se résume pas à un nombre ; c'est un intervalle de confiance : \textbf{valeur mesurée $\pm$ incertitude}. Ce module pose les bases du langage métrologique indispensable pour valider un modèle, comparer une théorie à l'expérience ou certifier un dispositif technique (ligne haute tension SONATREL, équipement hospitalier, télécommunications MTN/Orange). Nous allons structurer notre démarche en six étapes, conformément au programme officiel de Terminale C.

\courssub{Mesure, erreur et qualité d'un instrument}

\begin{definition}[Mesure, valeur vraie, erreur]
Une \cle{mesure} est l'ensemble des opérations ayant pour but de déterminer la valeur d'une grandeur physique. La \cle{valeur vraie} (ou valeur conventionnelle) est la valeur idéale, inaccessible en pratique. L'\cle{erreur de mesure} $E$ est la différence entre la valeur mesurée $x_m$ et la valeur vraie $x_v$ : $E = x_m - x_v$. Comme $x_v$ est inconnue, l'erreur ne peut \textbf{jamais} être calculée exactement.
\end{definition}

\begin{definition}[Qualités métrologiques d'un instrument]
Quatre caractéristiques définissent la qualité d'un appareil de mesure :
\begin{itemize}
    \item \cle{Sensibilité} : rapport de la variation de l'indication à la variation correspondante de la grandeur mesurée (pente de la courbe d'étalonnage). Unité : unité de l'indication / unité de la grandeur.
    \item \cle{Précision} (ou \cle{résolution}) : plus petite variation de la grandeur mesurée qui produit un changement perceptible de l'indication (ex: 1 mm sur un pied à coulisse, 0,1 V sur un voltmètre numérique).
    \item \cle{Fidélité} (ou \cle{répétabilité}) : aptitude à donner des indications très proches entre elles lors de mesures répétées dans les mêmes conditions (faible dispersion). Elle se quantifie par l'écart-type.
    \item \cle{Justesse} : aptitude à donner, en moyenne sur une série de mesures, une valeur proche de la valeur vraie (faible erreur systématique). Un instrument \textbf{juste} n'a pas de biais systématique.
\end{itemize}
\end{definition}

\begin{propriete}[Lien entre fidélité et justesse]
Un instrument peut être fidèle sans être juste (mesures groupées mais décalées : erreur systématique). Il ne peut être juste sans être fidèle (moyenne proche de la vraie valeur mais dispersion énorme $\rightarrow$ incertitude trop grande pour affirmer la justesse). \textbf{Justesse = Fidélité + Absence de biais systématique.}
\end{propriete}

\begin{attention}[Pièges de vocabulaire]
\begin{itemize}
    \item Ne confonds pas \textbf{précision} (résolution de l'appareil) et \textbf{fidélité} (dispersion des mesures). Un appareil à haute résolution (ex: 0,001 g) peut donner des mesures très dispersées (faible fidélité) si l'environnement vibre.
    \item L'\textbf{erreur} est une notion théorique (inconnue) ; l'\textbf{incertitude} est une estimation quantitative de la qualité du résultat (intervalle où se trouve la valeur vraie avec une probabilité donnée).
\end{itemize}
\end{attention}

\begin{exemplebox}[Qualités d'un multimètre au laboratoire de Nkongsamba]
On teste un multimètre numérique sur une source de tension étalon $U_{\text{ref}} = 10,00~\text{V}$.
\begin{itemize}
    \item \textbf{Sensibilité} : $10~\text{mV/V}$ (l'affichage change de 10 mV pour 1 V d'entrée).
    \item \textbf{Précision (résolution)} : $1~\text{mV}$ (dernier chiffre affiché).
    \item \textbf{Fidélité} : 10 mesures donnent $10,02~\text{V} \pm 0,01~\text{V}$ (dispersion faible $\rightarrow$ fidèle).
    \item \textbf{Justesse} : La moyenne $10,02~\text{V}$ diffère de $0,02~\text{V}$ de la valeur vraie $\rightarrow$ biais systématique de $+0,02~\text{V}$. L'instrument n'est pas parfaitement juste.
\end{itemize}
\end{exemplebox}

\courssub{Incertitude d'une mesure unique (Évaluation de type B)}

\begin{methode}[Estimation de l'incertitude pour une mesure unique]
\begin{enumerate}
    \item \textbf{Incertitude instrumentale (classe de l'appareil)} : pour un appareil analogique de classe $c$ utilisé sur un calibre $K$ (valeur max de l'échelle), $\Delta U_{\text{instr}} = \dfrac{c \times K}{100}$.
    \item \textbf{Incertitude de lecture} : pour un appareil analogique, on estime généralement $\pm$ demi-graduation ($\frac{1}{2}$ div). Pour un numérique, $\pm 1$ digit (dernier chiffre).
    \item \textbf{Incertitude totale} : on combine les composantes. Si elles sont indépendantes, on fait la \textbf{somme quadratique} (incertitude-type composée) : $u_c = \sqrt{u_1^2 + u_2^2}$. Pour une borne maximale (approche conservatrice), on fait la \textbf{somme algébrique} : $\Delta_{\text{max}} = \Delta_1 + \Delta_2$.
\end{enumerate}
\end{methode}

\begin{exemplebox}[Voltmètre analogique -- Application directe]
Un voltmètre de classe $1,5$, calibre $30~\text{V}$, échelle $100$ divisions. Aiguille sur division $67$.
\begin{enumerate}
    \item Valeur lue : $U = \frac{67}{100} \times 30 = 20,1~\text{V}$.
    \item Incertitude instrumentale : $\Delta U_{\text{instr}} = \frac{1,5 \times 30}{100} = 0,45~\text{V}$.
    \item Incertitude de lecture (1/2 div) : $\Delta U_{\text{lect}} = \frac{0,5}{100} \times 30 = 0,15~\text{V}$.
    \item Incertitude-type composée : $u(U) = \sqrt{0,45^2 + 0,15^2} \approx 0,47~\text{V}$.
    \item Résultat : $U = (20,1 \pm 0,5)~\text{V}$ (incertitude arrondie à 1 ch. sig. $\rightarrow$ 0,5 V ; valeur au même rang : 20,1 V). Incertitude relative : $\frac{0,5}{20,1} \approx 2,5~\%$.
\end{enumerate}
\end{exemplebox}

\courssub{Série de mesures : moyenne, écart-type, incertitude-type (Évaluation de type A)}

\begin{definition}[Grandeurs statistiques]
Pour une série de $n$ mesures $x_1, x_2, \dots, x_n$ :
\begin{itemize}
    \item \cle{Moyenne arithmétique} : $\bar{x} = \frac{1}{n}\sum_{i=1}^n x_i$ (meilleure estimation de la valeur vraie).
    \item \cle{Écart-type expérimental} (ou écart-type échantillon) : $\sigma_{n-1} = \sqrt{\frac{1}{n-1}\sum_{i=1}^n (x_i - \bar{x})^2}$. Il caractérise la dispersion (fidélité).
    \item \cle{Incertitude-type sur la moyenne} (ou erreur standard de la moyenne) : $u(\bar{x}) = \frac{\sigma_{n-1}}{\sqrt{n}}$. C'est l'incertitude sur la meilleure estimation $\bar{x}$.
\end{itemize}
\end{definition}

\begin{propriete}[Propriétés de l'estimateur]
\begin{itemize}
    \item $\bar{x}$ est un estimateur sans biais de la valeur vraie.
    \item $u(\bar{x})$ diminue comme $1/\sqrt{n}$ : pour diviser l'incertitude par 2, il faut multiplier le nombre de mesures par 4.
    \item L'intervalle $[\bar{x} - u(\bar{x})~;~ \bar{x} + u(\bar{x})]$ a une probabilité d'environ $68~\%$ de contenir la valeur vraie (loi normale, $n$ grand).
\end{itemize}
\end{propriete}

\begin{attention}[Règles d'arrondi -- Exigibles au Bac]
\begin{enumerate}
    \item L'incertitude s'arrondit à \textbf{1 chiffre significatif} (parfois 2 si le premier chiffre est 1 ou 2).
    \item La valeur moyenne s'arrondit au \textbf{même rang décimal} que l'incertitude.
    \item Exemple : $\bar{x} = 12,3456$, $u = 0,0234 \rightarrow u = 0,02$, $\bar{x} = 12,35$. Résultat : $x = (12,35 \pm 0,02)~\text{unité}$.
\end{enumerate}
\end{attention}

\courssub{Expression d'un résultat de mesure}

\begin{aretenir}[Forme canonique d'un résultat]
\[ \boxed{X = (x_{\text{best}} \pm u(x))~\text{unité SI}} \]
avec $u(x)$ à 1 ou 2 ch. sig. et $x_{\text{best}}$ aligné sur la même décimale.
\end{aretenir}

\begin{definition}[Incertitude relative]
L'\cle{incertitude relative} (sans unité) est $u_r(x) = \dfrac{u(x)}{|x_{\text{best}}|}$. Elle s'exprime souvent en \textbf{pourcentage} : $u_r(\%) = 100 \times \dfrac{u(x)}{|x_{\text{best}}|}$.
\textbf{Intérêt} : permet de comparer la qualité de mesures de grandeurs différentes (ex: comparer la précision d'une mesure de résistance $1~\%$ et d'une mesure de tension $0,5~\%$).
\end{definition}

\begin{propriete}[Propagation des incertitudes (loi de composition) -- Exigible au Bac]
Pour une fonction $f(x_1, x_2, \dots)$ avec variables indépendantes :
\[ u(f) = \sqrt{ \sum_i \left( \frac{\partial f}{\partial x_i} \right)^2 u^2(x_i) } \]
Cas particuliers usuels :
\begin{itemize}
    \item Somme/Différence $f = x \pm y$ : $u(f) = \sqrt{u^2(x) + u^2(y)}$.
    \item Produit/Quotient $f = x \times y$ ou $f = x / y$ : $\dfrac{u(f)}{|f|} = \sqrt{ \left(\dfrac{u(x)}{x}\right)^2 + \left(\dfrac{u(y)}{y}\right)^2 }$.
    \item Puissance $f = k x^n$ : $\dfrac{u(f)}{|f|} = |n| \dfrac{u(x)}{x}$.
\end{itemize}
\end{propriete}

\courssub{Le Système international d'unités (SI)}

\begin{definition}[Les 7 unités de base du SI (définitions 2019 basées sur constantes fondamentales)]
\begin{center}
\begin{tabularx}{\linewidth}{|l|X|l|}\hline
\rowcolor{popblueL}
\textbf{Grandeur} & \textbf{Unité (symbole)} & \textbf{Constante de définition} \\ \hline
Longueur & mètre (\textbf{m}) & Vitesse de la lumière $c = 299\,792\,458~\text{m/s}$ \\
Masse & kilogramme (\textbf{kg}) & Constante de Planck $h = 6,626\,070\,15 \times 10^{-34}~\text{J}\cdot\text{s}$ \\
Temps & seconde (\textbf{s}) & Fréquence hyperfine Cs $\Delta\nu_{\text{Cs}} = 9\,192\,631\,770~\text{Hz}$ \\
Courant & ampère (\textbf{A}) & Charge élémentaire $e = 1,602\,176\,634 \times 10^{-19}~\text{C}$ \\
Température & kelvin (\textbf{K}) & Constante de Boltzmann $k_B = 1,380\,649 \times 10^{-23}~\text{J/K}$ \\
Quantité de matière & mole (\textbf{mol}) & Nombre d'Avogadro $N_A = 6,022\,140\,76 \times 10^{23}~\text{mol}^{-1}$ \\
Intensité lumineuse & candela (\textbf{cd}) & Efficacité lumineuse $K_{\text{cd}} = 683~\text{lm/W}$ \\ \hline
\end{tabularx}
\end{center}
\end{definition}

\begin{propriete}[Unités dérivées usuelles en Terminale C]
\begin{center}
\begin{tabular}{|l|l|l|}\hline
\rowcolor{popblueL}
\textbf{Grandeur} & \textbf{Unité dérivée (symbole)} & \textbf{Expression en unités de base} \\ \hline
Force & newton (N) & $\text{kg}\cdot\text{m}\cdot\text{s}^{-2}$ \\
Énergie, Travail & joule (J) & $\text{kg}\cdot\text{m}^2\cdot\text{s}^{-2}$ \\
Puissance & watt (W) & $\text{kg}\cdot\text{m}^2\cdot\text{s}^{-3}$ \\
Charge & coulomb (C) & $\text{A}\cdot\text{s}$ \\
Tension, Différence de potentiel & volt (V) & $\text{kg}\cdot\text{m}^2\cdot\text{s}^{-3}\cdot\text{A}^{-1}$ \\
Résistance, Impédance & ohm ($\Omega$) & $\text{kg}\cdot\text{m}^2\cdot\text{s}^{-3}\cdot\text{A}^{-2}$ \\
Capacité & farad (F) & $\text{kg}^{-1}\cdot\text{m}^{-2}\cdot\text{s}^4\cdot\text{A}^2$ \\
Inductance & henry (H) & $\text{kg}\cdot\text{m}^2\cdot\text{s}^{-2}\cdot\text{A}^{-2}$ \\
Champ magnétique & tesla (T) & $\text{kg}\cdot\text{s}^{-2}\cdot\text{A}^{-1}$ \\
Activité radioactive & becquerel (Bq) & $\text{s}^{-1}$ \\ \hline
\end{tabular}
\end{center}
\end{propriete}

\begin{savaistu}[Le SI au Cameroun]
L'Agence des Normes et de la Qualité (ANOR) est l'organisme national de référence pour la métrologie légale. Les compteurs ENEO, les balances des marchés, les instruments des hôpitaux (radiothérapie, imagerie) doivent être vérifiés périodiquement selon les standards du SI pour garantir l'équité commerciale et la sécurité des patients.
\end{savaistu}

\courssub{Analyse dimensionnelle et homogénéité}

\begin{definition}[Dimension d'une grandeur physique]
La \cle{dimension} d'une grandeur exprime sa nature en fonction des grandeurs de base : Longueur $[L]$, Masse $[M]$, Temps $[T]$, Courant $[I]$, Température $[\Theta]$. Une grandeur $G$ s'écrit $[G] = [M]^\alpha [L]^\beta [T]^\gamma [I]^\delta [\Theta]^\epsilon$.
\end{definition}

\begin{aretenir}[Dimensions des nouvelles grandeurs du programme Terminale C]
\begin{center}
\begin{tabular}{|l|c|}\hline
\rowcolor{popblueL}
\textbf{Grandeur} & \textbf{Dimension} \\ \hline
Accélération $\vec{a}$ & $[L][T]^{-2}$ \\
Champ électrique $\vec{E}$ & $[M][L][T]^{-3}[I]^{-1}$ \\
Champ magnétique $\vec{B}$ & $[M][T]^{-2}[I]^{-1}$ \\
Capacité $C$ & $[M]^{-1}[L]^{-2}[T]^{4}[I]^{2}$ \\
Inductance $L$ & $[M][L]^{2}[T]^{-2}[I]^{-2}$ \\
Impédance $Z$ & $[M][L]^{2}[T]^{-3}[I]^{-2}$ (même que résistance) \\
Activité $A$ & $[T]^{-1}$ \\ \hline
\end{tabular}
\end{center}
\end{aretenir}

\begin{methode}[Vérifier l'homogénéité d'une relation / Trouver une loi par analyse dimensionnelle]
\begin{enumerate}
    \item Écrire la dimension de chaque terme de l'équation.
    \item Vérifier que \textbf{tous les termes additifs ont la même dimension} (homogénéité).
    \item Pour trouver une loi $X = k A^\alpha B^\beta \dots$ ($k$ sans dimension) : écrire $[X] = [A]^\alpha [B]^\beta \dots$, identifier les exposants de $[M], [L], [T], [I]$, résoudre le système linéaire.
\end{enumerate}
\end{methode}

\begin{propriete}[Théorème de Vaschy-Buckingham (Pi-théorème) -- Culture]
Si une relation physique lie $n$ grandeurs dont les dimensions dépendent de $k$ dimensions de base indépendantes, la relation peut s'écrire sous la forme d'une relation entre $n-k$ groupes de grandeurs \textbf{sans dimension} (nombres de Reynolds, Mach, etc.). L'analyse dimensionnelle ne détermine \textbf{jamais} les constantes adimensionnelles ($2\pi$, $1/2$, etc.).
\end{propriete}

\begin{attention}[Limites de l'analyse dimensionnelle]
\begin{itemize}
    \item Une relation homogène n'est \textbf{pas nécessairement exacte} (ex: $T = \sqrt{LC}$ au lieu de $2\pi\sqrt{LC}$ ; $E = mc^2$ vs $E = \frac{1}{2}mc^2$).
    \item Elle ne gère pas les sommes de termes de même dimension mais de nature différente (ex: énergie cinétique + énergie potentielle).
    \item Elle ne remplace pas la démonstration physique (lois de Newton, Maxwell, thermodynamique).
\end{itemize}
\end{attention}

\begin{exemplebox}[Période propre d'un circuit LC -- Validation par les dimensions]
On cherche $T(L, C)$. $[T] = [T]$. $[L] = [M][L]^2[T]^{-2}[I]^{-2}$. $[C] = [M]^{-1}[L]^{-2}[T]^4[I]^2$.
Produit $LC$ : $[M]^0[L]^0[T]^2[I]^0 = [T]^2$. Donc $\sqrt{LC}$ a dimension $[T]$.
Toute expression de la forme $T = k \sqrt{LC}$ est homogène ($k$ adimensionnel). L'analyse dimensionnelle ne donne pas $k = 2\pi$.
\end{exemplebox}

\begin{exemplebox}[Célérité d'une onde sur corde -- Recherche de loi]
$v = k F^\alpha \mu^\beta$. $[v]=[L][T]^{-1}$, $[F]=[M][L][T]^{-2}$, $[\mu]=[M][L]^{-1}$.
$[L][T]^{-1} = ([M][L][T]^{-2})^\alpha ([M][L]^{-1})^\beta = [M]^{\alpha+\beta}[L]^{\alpha-\beta}[T]^{-2\alpha}$.
Système : $\begin{cases} \alpha+\beta=0 \\ \alpha-\beta=1 \\ -2\alpha=-1 \end{cases} \Rightarrow \alpha=1/2, \beta=-1/2$.
Loi : $\boxed{v = k \sqrt{\dfrac{F}{\mu}}}$. L'expérience donne $k=1$.
\end{exemplebox}

% ============================================================
% EXERCICES : Je vérifie mes connaissances
% ============================================================
\courssub{Je vérifie mes connaissances}

\begin{exercice}[QCM de synthèse]
Pour chaque question, une seule réponse est exacte. Recopie la lettre correspondante et justifie brièvement.

\textbf{1.} La \cle{fidélité} d'un instrument de mesure caractérise :
\begin{itemize}
    \item[\textbf{a)}] sa capacité à donner des indications proches de la valeur vraie ;
    \item[\textbf{b)}] sa capacité à donner des indications très proches entre elles lors de mesures répétées dans les mêmes conditions ;
    \item[\textbf{c)}] la plus petite variation de la grandeur mesurable ;
    \item[\textbf{d)}] la rapidité de réponse de l'instrument.
\end{itemize}

\textbf{2.} Le nombre de chiffres significatifs de la mesure $L = \num{0,03040}~\text{m}$ est :
\begin{itemize}
    \item[\textbf{a)}] 2 \quad \textbf{b)} 3 \quad \textbf{c)} 4 \quad \textbf{d)} 5
\end{itemize}

\textbf{3.} La dimension de la \cle{capacité} $C$ d'un condensateur ($q = C\,u$, avec $u$ une tension) est :
\begin{itemize}
    \item[\textbf{a)}] $[C] = \text{M}^{-1}\text{L}^{-2}\text{T}^{4}\text{I}^{2}$
    \item[\textbf{b)}] $[C] = \text{M}\,\text{L}^{2}\text{T}^{-3}\text{I}^{-1}$
    \item[\textbf{c)}] $[C] = \text{M}^{-1}\text{L}^{-2}\text{T}^{3}\text{I}$
    \item[\textbf{d)}] $[C] = \text{M}\,\text{L}^{2}\text{T}^{-2}$
\end{itemize}

\textbf{4.} L'\cle{activité} $A$ d'un échantillon radioactif a pour dimension :
\begin{itemize}
    \item[\textbf{a)}] $[A] = \text{T}$ \quad \textbf{b)} $[A] = \text{T}^{-1}$ \quad \textbf{c)} $[A] = \text{T}^{-2}$ \quad \textbf{d)} $[A]$ est sans dimension
\end{itemize}

\textbf{5.} L'\cle{impédance} $Z$ d'un dipôle a la même dimension que :
\begin{itemize}
    \item[\textbf{a)}] une inductance \quad \textbf{b)} une résistance \quad \textbf{c)} une capacité \quad \textbf{d)} une puissance
\end{itemize}
\end{exercice}

\begin{exercice}[Vrai ou faux justifié]
Pour chaque affirmation, précise si elle est vraie ou fausse, puis justifie en une ou deux phrases.
\begin{enumerate}
    \item L'erreur de mesure est la différence entre la valeur mesurée et la valeur vraie ; elle peut toujours être calculée exactement.
    \item Un instrument \cle{juste} donne, en moyenne sur un grand nombre de mesures, une valeur proche de la valeur vraie.
    \item L'incertitude-type $u(x)$ d'une série de $n$ mesures est égale à l'écart-type expérimental $\sigma_{n-1}$.
    \item Une relation physique homogène est nécessairement exacte.
    \item L'unité SI de l'inductance, le henry (H), est une unité dérivée.
\end{enumerate}
\end{exercice}

\begin{exercice}[Définitions et vocabulaire]
\begin{enumerate}
    \item Définis les quatre qualités métrologiques d'un instrument de mesure : \cle{sensibilité}, \cle{précision}, \cle{fidélité}, \cle{justesse}.
    \item Un élève mesure vingt fois la durée de chute d'une bille avec un chronomètre. Les valeurs obtenues sont très dispersées, mais leur moyenne est très proche de la valeur attendue. Quelle qualité fait défaut à cette mesure ? Quelle qualité est au contraire satisfaite ?
    \item Qu'appelle-t-on \cle{incertitude relative} ? Quel est son intérêt pour comparer deux mesures de grandeurs différentes ?
\end{enumerate}
\end{exercice}

\begin{exercice}[Calculs directs]
\begin{enumerate}
    \item Exprime en notation scientifique, avec trois chiffres significatifs, les grandeurs suivantes : $d = \num{0,000452}~\text{m}$ ; $m = \num{12500}~\text{g}$ ; $t = \num{0,0806}~\text{s}$.
    \item Une masse est mesurée : $m = (\num{250,0} \pm \num{0,5})~\text{g}$. Calcule l'incertitude relative en pourcentage.
    \item Un voltmètre numérique affiche $U = \num{12,36}~\text{V}$ avec une incertitude-type $u(U) = \num{0,08}~\text{V}$. Écris le résultat de la mesure sous la forme $U \pm u(U)$ en respectant la règle sur les chiffres significatifs de l'incertitude.
    \item Sachant que $[F] = \text{M}\,\text{L}\,\text{T}^{-2}$ pour une force, détermine la dimension et l'unité SI de la constante de raideur $k$ d'un ressort définie par $F = k\,x$.
\end{enumerate}
\end{exercice}

% ============================================================
% EXERCICES : Je m'entraîne
% ============================================================
\courssub{Je m'entraîne}

\begin{exercice}[Mesure unique au voltmètre analogique]
Au laboratoire du lycée de Nkongsamba, un élève mesure la tension aux bornes d'une lampe avec un voltmètre analogique de \cle{classe} $1,5$, utilisé sur le calibre $30~\text{V}$ avec une échelle de $100$ divisions.
\begin{enumerate}
    \item L'aiguille s'arrête sur la division $67$. Quelle est la valeur lue $U$ de la tension ?
    \item Calcule l'incertitude instrumentale $\Delta U_{\text{instr}} = \dfrac{\text{classe} \times \text{calibre}}{100}$.
    \item En ajoutant une incertitude de lecture estimée à une demi-graduation, exprime le résultat de la mesure sous la forme $U \pm \Delta U_{\text{tot}}$, puis donne l'incertitude relative en pourcentage. (On combinera les incertitudes par somme quadratique pour obtenir l'incertitude-type composée).
\end{enumerate}
\end{exercice}

\begin{exercice}[Homogénéité : période propre d'un circuit LC]
Trois élèves proposent chacun une expression de la période propre d'oscillation d'un circuit $(L, C)$ :
\[ T_1 = 2\pi\sqrt{LC} \qquad T_2 = 2\pi\sqrt{\dfrac{L}{C}} \qquad T_3 = 2\pi\,L\,C \]
\begin{enumerate}
    \item Rappelle les dimensions de l'inductance $L$ et de la capacité $C$ (on pourra utiliser $u = L\dfrac{di}{dt}$ et $q = C\,u$).
    \item Détermine la dimension de chacun des produits $LC$, $\dfrac{L}{C}$ et des trois expressions proposées.
    \item Élimine les expressions non homogènes et conclus. L'expression retenue est-elle pour autant démontrée ? Commente.
\end{enumerate}
\end{exercice}

\begin{exercice}[Compatibilité de deux mesures de g]
Deux groupes d'élèves de Terminale C, l'un au lycée de Yaoundé et l'autre au lycée de Bafoussam, déterminent l'intensité de la pesanteur à l'aide d'un pendule simple. Ils annoncent :
\[ g_1 = (\num{9,78} \pm \num{0,03})~\text{m}\cdot\text{s}^{-2} \qquad g_2 = (\num{9,83} \pm \num{0,02})~\text{m}\cdot\text{s}^{-2} \]
\begin{enumerate}
    \item Représente sur un axe gradué les intervalles $[g_1 - u(g_1)\,;\,g_1 + u(g_1)]$ et $[g_2 - u(g_2)\,;\,g_2 + u(g_2)]$.
    \item Deux mesures sont dites \cle{compatibles} si leurs intervalles d'incertitude se recoupent. Ces deux mesures sont-elles compatibles ? Justifie.
    \item Calcule l'écart relatif $\dfrac{|g_2 - g_1|}{\bar{g}}$ en pourcentage, où $\bar{g}$ est la moyenne des deux valeurs. Propose une cause possible du désaccord observé.
\end{enumerate}
\end{exercice}

\begin{exercice}[Analyse dimensionnelle guidée : onde sur une corde]
Une mototaxi de livraison fait vibrer une corde tendue entre deux poteaux. On cherche l'expression de la célérité $v$ de l'onde transversale le long de la corde sous la forme :
\[ v = k\,F^{\alpha}\,\mu^{\beta} \]
où $F$ est la tension de la corde (une force), $\mu$ sa masse linéique (masse par unité de longueur) et $k$ une constante sans dimension.
\begin{enumerate}
    \item Donne les dimensions de $v$, $F$ et $\mu$.
    \item Écris l'équation aux dimensions de la relation proposée.
    \item En identifiant les exposants de M, L et T, établis un système de trois équations d'inconnues $\alpha$ et $\beta$.
    \item Résous ce système et donne l'expression finale de $v$. Vérifie son homogénéité.
\end{enumerate}
\end{exercice}

% ============================================================
% EXERCICES : Je me prépare au Bac
% ============================================================
\courssub{Je me prépare au Bac}

\begin{exercice}[Type Bac — Série de mesures : période d'un pendule simple]
Dans le cadre des travaux pratiques, un groupe d'élèves de Terminale C du lycée de Douala étudie un pendule simple de longueur $\ell = (\num{1,000} \pm \num{0,002})~\text{m}$. Pour déterminer la période $T$ des petites oscillations, un élève chronomètre huit fois la durée de $10$ oscillations complètes et en déduit huit valeurs de la période (en secondes) :
\begin{center}
\begin{tabular}{|l|cccccccc|}
\hline
Mesure $n^{\circ}$ & 1 & 2 & 3 & 4 & 5 & 6 & 7 & 8 \\
\hline
$T$ (s) & 2,01 & 1,99 & 2,02 & 2,00 & 1,98 & 2,01 & 2,03 & 2,00 \\
\hline
\end{tabular}
\end{center}
On donne : $\displaystyle\sum_{i=1}^{8} T_i = \num{16,04}~\text{s}$ et $\displaystyle\sum_{i=1}^{8} (T_i - \bar{T})^2 = \num{1,75e-3}~\text{s}^2$.
\begin{enumerate}
    \item \textbf{Exploitation statistique.}
    \begin{enumerate}
        \item Calcule la valeur moyenne $\bar{T}$ de la période.
        \item Calcule l'écart-type expérimental $\sigma_{n-1} = \sqrt{\dfrac{1}{n-1}\displaystyle\sum_{i=1}^{n}(T_i - \bar{T})^2}$.
        \item En déduis l'incertitude-type sur la moyenne : $u(\bar{T}) = \dfrac{\sigma_{n-1}}{\sqrt{n}}$.
        \item Exprime le résultat de la mesure sous la forme $T = \bar{T} \pm u(\bar{T})$, en arrondissant l'incertitude à un chiffre significatif. Calcule l'incertitude relative en pourcentage.
    \end{enumerate}
    \item \textbf{Détermination de g.} La période du pendule simple est $T = 2\pi\sqrt{\dfrac{\ell}{g}}$.
    \begin{enumerate}
        \item Vérifie par analyse dimensionnelle que cette relation est homogène.
        \item Exprime $g$ en fonction de $T$ et $\ell$, puis calcule sa valeur numérique.
        \item L'incertitude relative sur $g$ vérifie $\dfrac{u(g)}{g} = \sqrt{\left(\dfrac{u(\ell)}{\ell}\right)^2 + \left(2\,\dfrac{u(T)}{T}\right)^2}$. Calcule $u(g)$ et exprime le résultat final $g \pm u(g)$.
        \item La valeur tabulée à Douala est $g_{\text{ref}} = \num{9,78}~\text{m}\cdot\text{s}^{-2}$. La mesure est-elle compatible avec cette valeur ?
    \end{enumerate}
\end{enumerate}
\end{exercice}

\begin{exercice}[Type Bac — Mesure d'une résistance et incertitudes]
Au laboratoire du collège de Garoua, on souhaite déterminer la résistance $R$ d'un conducteur ohmique par deux méthodes.

\textbf{Partie A : mesure au ohmmètre.} L'ohmmètre numérique utilisé affiche $R = \num{46,7}~\Omega$. La notice indique : « précision : $\pm(1\,\%$ de la lecture $+ 2$ digits$)$ », le digit valant $\num{0,1}~\Omega$ sur ce calibre.
\begin{enumerate}
    \item Calcule l'incertitude $\Delta R$ de cette mesure.
    \item Exprime le résultat sous la forme $R \pm \Delta R$ avec un nombre correct de chiffres significatifs, puis donne l'incertitude relative en pourcentage.
\end{enumerate}

\textbf{Partie B : méthode volt-ampèremétrique.} On mesure la tension $U$ aux bornes du résistor et l'intensité $I$ qui le traverse :
\[ U = (\num{4,52} \pm \num{0,05})~\text{V} \qquad I = (\num{96} \pm \num{2})~\text{mA} \]
\begin{enumerate}[resume]
    \item Calcule la valeur de $R = \dfrac{U}{I}$.
    \item Sachant que $\dfrac{u(R)}{R} = \sqrt{\left(\dfrac{u(U)}{U}\right)^2 + \left(\dfrac{u(I)}{I}\right)^2}$, calcule l'incertitude relative puis l'incertitude-type $u(R)$. Exprime le résultat $R \pm u(R)$.
    \item Les résultats des deux parties sont-ils compatibles ? Justifie.
    \item \textbf{Analyse dimensionnelle.} La puissance dissipée par effet Joule est $P = R\,I^2$. Vérifie que $R\,I^2$ est homogène à une puissance, c'est-à-dire de dimension $\text{M}\,\text{L}^2\,\text{T}^{-3}$.
\end{enumerate}
\end{exercice}

\begin{exercice}[Type Bac — Analyse dimensionnelle et validation de modèles]
Lors de la conception d'une ligne de transport d'électricité entre le barrage de Nachtigal et Yaoundé, les ingénieurs modélisent plusieurs grandeurs physiques. On te propose d'utiliser l'analyse dimensionnelle pour valider leurs expressions. On rappelle : $[F] = \text{M}\,\text{L}\,\text{T}^{-2}$ (force), $[U] = \text{M}\,\text{L}^2\,\text{T}^{-3}\,\text{I}^{-1}$ (tension), $[E] = \text{M}\,\text{L}\,\text{T}^{-3}\,\text{I}^{-1}$ (champ électrique), $[B] = \text{M}\,\text{T}^{-2}\,\text{I}^{-1}$ (champ magnétique).
\begin{enumerate}
    \item \textbf{Champ électrique.} Deux expressions du champ électrique entre les armatures d'un condensateur plan sont proposées : $E = \dfrac{U}{d}$ et $E = U \times d$, où $d$ est la distance entre les armatures. En utilisant les dimensions rappelées, détermine laquelle est homogène.
    \item \textbf{Capacité d'un condensateur plan.} On propose $C = \varepsilon\,\dfrac{S}{d}$, où $S$ est la surface des armatures, $d$ leur distance et $\varepsilon$ la permittivité du milieu.
    \begin{enumerate}
        \item Sachant que $[C] = \text{M}^{-1}\,\text{L}^{-2}\,\text{T}^{4}\,\text{I}^{2}$, détermine la dimension de $\varepsilon$.
        \item L'unité SI de $\varepsilon$ est le farad par mètre ($\text{F}\cdot\text{m}^{-1}$). Ce résultat est-il cohérent avec ta réponse précédente ?
    \end{enumerate}
    \item \textbf{Force de Laplace sur la ligne.} Un tronçon de ligne de longueur $\ell$, parcouru par un courant $I$ et placé dans un champ magnétique $B$, subit une force dont on propose trois expressions :
    \[ F_1 = B\,I\,\ell \qquad F_2 = \dfrac{B\,I}{\ell} \qquad F_3 = B\,I\,\ell^2 \]
    Par analyse dimensionnelle, élimine les expressions impossibles.
    \item \textbf{Validation d'un modèle.} Un technicien affirme que la période d'oscillation d'un circuit $(R, L, C)$ série faiblement amorti est $T = 2\pi\sqrt{LC}$, indépendante de $R$.
    \begin{enumerate}
        \item Vérifie l'homogénéité de cette expression.
        \item Explique pourquoi l'homogénéité seule ne suffit pas à valider le modèle, et propose une démarche expérimentale permettant de tester l'hypothèse « $T$ ne dépend pas de $R$ ».
    \end{enumerate}
\end{enumerate}
\end{exercice}

\newpage

\coursec{Corrigés des exercices}

\begin{corrige}[Exercice 1 — QCM de synthèse]
\textbf{1. Réponse b).} La fidélité (ou répétabilité) caractérise l'aptitude d'un instrument à donner des indications très proches entre elles lors de mesures répétées dans les mêmes conditions. La réponse a) définit la justesse, la réponse c) définit la précision (résolution).

\textbf{2. Réponse c) : 4 chiffres significatifs.} Dans $\num{0,03040}$, les zéros à gauche ne comptent pas ; les chiffres significatifs sont $3$, $0$, $4$, $0$ (le zéro final, placé après la virgule, est significatif car il traduit la précision de la mesure).

\textbf{3. Réponse a).} De $q = Cu$ : $[C] = \dfrac{[q]}{[u]} = \dfrac{\text{I}\,\text{T}}{\text{M}\,\text{L}^2\,\text{T}^{-3}\,\text{I}^{-1}} = \text{M}^{-1}\,\text{L}^{-2}\,\text{T}^{4}\,\text{I}^{2}$.

\textbf{4. Réponse b).} L'activité est un nombre de désintégrations par unité de temps : $[A] = \text{T}^{-1}$ (unité SI : le becquerel, $\text{Bq} = \text{s}^{-1}$).

\textbf{5. Réponse b).} L'impédance est le rapport d'une tension par une intensité ($Z = U/I$ en valeur efficace), comme la résistance : $[Z] = [R] = \text{M}\,\text{L}^2\,\text{T}^{-3}\,\text{I}^{-2}$ (unité : l'ohm).
\end{corrige}

\begin{corrige}[Exercice 2 — Vrai ou faux justifié]
\begin{enumerate}
    \item \textbf{Faux.} La définition $E = x_m - x_v$ est correcte, mais la valeur vraie $x_v$ est par nature inaccessible : l'erreur ne peut \textbf{jamais} être calculée exactement. On ne peut qu'estimer une \emph{incertitude}.
    \item \textbf{Vrai.} C'est la définition de la justesse : absence d'erreur systématique, c'est-à-dire moyenne des mesures proche de la valeur vraie.
    \item \textbf{Faux.} L'écart-type expérimental $\sigma_{n-1}$ caractérise la dispersion des mesures individuelles ; l'incertitude-type sur la \emph{moyenne} est $u(\bar{x}) = \dfrac{\sigma_{n-1}}{\sqrt{n}}$. Les deux ne sont égales que pour $n = 1$ (auquel cas $\sigma_{n-1}$ n'est d'ailleurs pas défini).
    \item \textbf{Faux.} L'homogénéité est une condition \emph{nécessaire} mais non \emph{suffisante} : $T = \sqrt{LC}$ et $T = 2\pi\sqrt{LC}$ sont toutes deux homogènes, mais seule la seconde est exacte. L'analyse dimensionnelle ne détermine pas les constantes sans dimension.
    \item \textbf{Vrai.} Le henry s'exprime en unités de base : $\text{H} = \text{kg}\cdot\text{m}^2\cdot\text{s}^{-2}\cdot\text{A}^{-2}$ ; ce n'est pas l'une des sept unités de base du SI.
\end{enumerate}
\end{corrige}

\begin{corrige}[Exercice 3 — Définitions et vocabulaire]
\begin{enumerate}
    \item \textbf{Sensibilité} : rapport de la variation de l'indication de l'instrument à la variation correspondante de la grandeur mesurée (pente de la courbe d'étalonnage). \textbf{Précision} (résolution) : plus petite variation de la grandeur mesurée qui provoque un changement perceptible de l'indication. \textbf{Fidélité} : aptitude à donner des indications très proches entre elles lors de mesures répétées dans les mêmes conditions (faible dispersion). \textbf{Justesse} : aptitude à donner, en moyenne sur une série de mesures, une valeur proche de la valeur vraie (absence de biais systématique).
    \item Les valeurs sont très dispersées : la \cle{fidélité} fait défaut. La moyenne est proche de la valeur attendue : la \cle{justesse} est satisfaite (pas d'erreur systématique, par exemple pas de décalage de zéro du chronomètre ; la dispersion vient du temps de réaction de l'opérateur, erreur aléatoire).
    \item L'\cle{incertitude relative} est le rapport $u_r(x) = \dfrac{u(x)}{|x_{\text{best}}|}$, souvent exprimée en pourcentage. Étant sans dimension, elle permet de comparer la \emph{qualité} de mesures portant sur des grandeurs de natures ou d'ordres de grandeur différents : une mesure à $0,5\,\%$ est de meilleure qualité relative qu'une mesure à $2\,\%$, quelles que soient les unités.
\end{enumerate}
\end{corrige}

\begin{corrige}[Exercice 4 — Calculs directs]
\begin{enumerate}
    \item $d = \num{4,52e-4}~\text{m}$ ; $m = \num{1,25e4}~\text{g}$ (soit $\SI{12,5}{kg}$) ; $t = \num{8,06e-2}~\text{s}$.
    \item Incertitude relative : $u_r(m) = \dfrac{0,5}{250,0} = \num{2,0e-3}$, soit $\boxed{0,20\,\%}$.
    \item L'incertitude $u(U) = \SI{0,08}{V}$ a déjà un chiffre significatif et s'exprime au centième de volt ; la valeur doit donc être arrondie au centième : $U = \SI{12,36}{V}$ convient. Résultat : $\boxed{U = (\num{12,36} \pm \num{0,08})~\text{V}}$.
    \item De $F = kx$ : $[k] = \dfrac{[F]}{[x]} = \dfrac{\text{M}\,\text{L}\,\text{T}^{-2}}{\text{L}} = \boxed{\text{M}\,\text{T}^{-2}}$. Unité SI : $\text{kg}\cdot\text{s}^{-2}$, c'est-à-dire le newton par mètre ($\text{N}\cdot\text{m}^{-1}$), cohérent avec $k = F/x$.
\end{enumerate}
\end{corrige}

\begin{corrige}[Exercice 5 — Mesure unique au voltmètre analogique]
\begin{enumerate}
    \item Proportionnalité lecture/échelle : $U = \dfrac{67}{100} \times 30 = \boxed{\num{20,1}~\text{V}}$.
    \item Incertitude instrumentale : $\Delta U_{\text{instr}} = \dfrac{1,5 \times 30}{100} = \boxed{\num{0,45}~\text{V}}$.
    \item Incertitude de lecture (demi-graduation) : une graduation vaut $\dfrac{30}{100} = \SI{0,30}{V}$, donc $\Delta U_{\text{lect}} = \dfrac{0,30}{2} = \SI{0,15}{V}$.

    Combinaison quadratique (sources indépendantes) :
    \[ u(U) = \sqrt{0,45^2 + 0,15^2} = \sqrt{0,2025 + 0,0225} = \sqrt{0,225} \approx \SI{0,474}{V}. \]
    On arrondit l'incertitude à un chiffre significatif : $u(U) = \SI{0,5}{V}$, et la valeur au même rang (dixième) : $\boxed{U = (\num{20,1} \pm \num{0,5})~\text{V}}$.

    Incertitude relative : $\dfrac{0,5}{20,1} \approx 0,025$, soit $\boxed{2,5\,\%}$.
\end{enumerate}
\textbf{Point de vigilance :} ne pas oublier de convertir la demi-graduation en volts (via le calibre) avant de la combiner avec l'incertitude instrumentale ; ne jamais additionner des divisions et des volts.
\end{corrige}

\begin{corrige}[Exercice 6 — Homogénéité : période propre d'un circuit LC]
\begin{enumerate}
    \item De $u = L\dfrac{di}{dt}$ : $[L] = \dfrac{[u]\,[t]}{[i]} = \dfrac{\text{M}\,\text{L}^2\,\text{T}^{-3}\,\text{I}^{-1} \times \text{T}}{\text{I}} = \text{M}\,\text{L}^2\,\text{T}^{-2}\,\text{I}^{-2}$.

    De $q = Cu$ : $[C] = \dfrac{[q]}{[u]} = \dfrac{\text{I}\,\text{T}}{\text{M}\,\text{L}^2\,\text{T}^{-3}\,\text{I}^{-1}} = \text{M}^{-1}\,\text{L}^{-2}\,\text{T}^{4}\,\text{I}^{2}$.
    \item Produit $LC$ :
    \[ [LC] = \text{M}^{1-1}\,\text{L}^{2-2}\,\text{T}^{-2+4}\,\text{I}^{-2+2} = \text{T}^2. \]
    Rapport $\dfrac{L}{C}$ :
    \[ \left[\dfrac{L}{C}\right] = \text{M}^{1-(-1)}\,\text{L}^{2-(-2)}\,\text{T}^{-2-4}\,\text{I}^{-2-2} = \text{M}^2\,\text{L}^4\,\text{T}^{-6}\,\text{I}^{-4}. \]
    Dimensions des trois expressions proposées ($2\pi$ est sans dimension) :
    \begin{align*}
    [T_1] &= \left[\sqrt{LC}\right] = \text{T} \quad \text{(homogène à une durée)} ;\\
    [T_2] &= \left[\sqrt{L/C}\right] = \text{M}\,\text{L}^2\,\text{T}^{-3}\,\text{I}^{-2} \quad \text{(dimension d'une résistance)} ;\\
    [T_3] &= [LC] = \text{T}^2 \quad \text{(carré d'une durée)}.
    \end{align*}
    \item Seule $T_1 = 2\pi\sqrt{LC}$ est homogène à une période ; $T_2$ et $T_3$ sont éliminées. \textbf{Conclusion prudente :} l'homogénéité ne \emph{démontre} pas l'exactitude de $T_1$ : toute expression $k\sqrt{LC}$ avec $k$ sans dimension serait aussi homogène. La valeur $k = 2\pi$ s'obtient uniquement par la résolution de l'équation différentielle $L\ddot{q} + q/C = 0$ (ou par l'expérience).
\end{enumerate}
\end{corrige}

\begin{corrige}[Exercice 7 — Compatibilité de deux mesures de g]
\begin{enumerate}
    \item Intervalles : $I_1 = [9,75\,;\,9,81]~\text{m}\cdot\text{s}^{-2}$ et $I_2 = [9,81\,;\,9,85]~\text{m}\cdot\text{s}^{-2}$.
    \begin{popfigure}
    \begin{tikzpicture}[x=8cm]
        \draw[->, thick] (9.70,0) -- (9.90,0) ;
        \foreach \x in {9.70,9.75,9.80,9.85,9.90}
            \draw (\x,0.03) -- (\x,-0.03) node[below, font=\small] {\num{\x}};
        \draw[line width=2.5pt, popblue] (9.75,0.25) -- (9.81,0.25);
        \fill[popblue] (9.78,0.25) circle (1.5pt);
        \node[popblue, above, font=\small] at (9.78,0.25) {$g_1$};
        \draw[line width=2.5pt, poporange] (9.81,0.55) -- (9.85,0.55);
        \fill[poporange] (9.83,0.55) circle (1.5pt);
        \node[poporange, above, font=\small] at (9.83,0.55) {$g_2$};
    \end{tikzpicture}
    \legende{Intervalles d'incertitude des deux mesures de $g$ (en $\text{m}\cdot\text{s}^{-2}$).}
    \end{popfigure}
    \item Les intervalles se touchent en $g = \SI{9,81}{m.s^{-2}}$ : $I_1 \cap I_2 = \{9,81\} \neq \varnothing$. Les deux mesures sont donc \textbf{compatibles}, mais de justesse : l'intersection est réduite à un point, ce qui invite à la prudence.
    \item Moyenne : $\bar{g} = \dfrac{9,78 + 9,83}{2} = \SI{9,805}{m.s^{-2}}$. Écart relatif :
    \[ \dfrac{|g_2 - g_1|}{\bar{g}} = \dfrac{0,05}{9,805} \approx \num{5,1e-3} \approx \boxed{0,5\,\%}. \]
    Causes possibles du désaccord : erreur systématique sur la mesure de la longueur du pendule (position du centre de la bille), amplitude des oscillations trop grande (la formule $T = 2\pi\sqrt{\ell/g}$ suppose les petits angles), temps de réaction différent des opérateurs au chronomètre, ou différence réelle (faible) de $g$ entre Yaoundé et Bafoussam (altitude).
\end{enumerate}
\end{corrige}

\begin{corrige}[Exercice 8 — Analyse dimensionnelle guidée : onde sur une corde]
\begin{enumerate}
    \item $[v] = \text{L}\,\text{T}^{-1}$ ; $[F] = \text{M}\,\text{L}\,\text{T}^{-2}$ (force) ; $[\mu] = \dfrac{[m]}{[\ell]} = \text{M}\,\text{L}^{-1}$.
    \item Équation aux dimensions :
    \[ \text{L}\,\text{T}^{-1} = \left(\text{M}\,\text{L}\,\text{T}^{-2}\right)^{\alpha}\left(\text{M}\,\text{L}^{-1}\right)^{\beta} = \text{M}^{\alpha+\beta}\,\text{L}^{\alpha-\beta}\,\text{T}^{-2\alpha}. \]
    \item Identification des exposants :
    \[ \begin{cases} \alpha + \beta = 0 & \text{(M)} \\ \alpha - \beta = 1 & \text{(L)} \\ -2\alpha = -1 & \text{(T)} \end{cases} \]
    \item De (T) : $\alpha = \dfrac{1}{2}$. De (M) : $\beta = -\alpha = -\dfrac{1}{2}$. Vérification dans (L) : $\alpha - \beta = \dfrac{1}{2} + \dfrac{1}{2} = 1$ : cohérent.

    Expression finale : $\boxed{v = k\sqrt{\dfrac{F}{\mu}}}$.

    Vérification d'homogénéité : $\left[\sqrt{F/\mu}\right] = \left(\dfrac{\text{M}\,\text{L}\,\text{T}^{-2}}{\text{M}\,\text{L}^{-1}}\right)^{1/2} = \left(\text{L}^2\,\text{T}^{-2}\right)^{1/2} = \text{L}\,\text{T}^{-1} = [v]$. La relation est homogène. (L'expérience ou la mécanique des cordes montre que $k = 1$.)
\end{enumerate}
\end{corrige}

\begin{corrige}[Exercice 9 — Type Bac : série de mesures, période d'un pendule simple]
\textbf{Barème indicatif (sur 10) :} 1.a) 1 pt ; 1.b) 1 pt ; 1.c) 1 pt ; 1.d) 1,5 pt ; 2.a) 1 pt ; 2.b) 1,5 pt ; 2.c) 1,5 pt ; 2.d) 1,5 pt.

\begin{enumerate}
    \item \textbf{Exploitation statistique.}
    \begin{enumerate}
        \item $\bar{T} = \dfrac{1}{8}\displaystyle\sum_{i=1}^{8} T_i = \dfrac{16,04}{8} = \boxed{\num{2,005}~\text{s}}$.
        \item $\sigma_{n-1} = \sqrt{\dfrac{1}{7}\times \num{1,75e-3}} = \sqrt{\num{2,5e-4}} \approx \boxed{\num{0,0158}~\text{s}} \approx \num{0,016}~\text{s}$.
        \item $u(\bar{T}) = \dfrac{\sigma_{n-1}}{\sqrt{n}} = \dfrac{0,0158}{\sqrt{8}} \approx \dfrac{0,0158}{2,828} \approx \boxed{\num{5,6e-3}~\text{s}}$.
        \item Incertitude arrondie à un chiffre significatif : $u(\bar{T}) = \SI{0,006}{s}$ ; la moyenne est arrondie au millième : $\bar{T} = \SI{2,005}{s}$.
        \[ \boxed{T = (\num{2,005} \pm \num{0,006})~\text{s}} \qquad u_r(T) = \dfrac{0,006}{2,005} \approx \boxed{0,30\,\%}. \]
    \end{enumerate}
    \item \textbf{Détermination de g.}
    \begin{enumerate}
        \item $[T] = \text{T}$ ; $\left[2\pi\sqrt{\ell/g}\right] = \sqrt{\dfrac{\text{L}}{\text{L}\,\text{T}^{-2}}} = \sqrt{\text{T}^2} = \text{T}$. La relation est homogène. \emph{(Point de vigilance : citer $[g] = \text{L}\,\text{T}^{-2}$ et rappeler que $2\pi$ est sans dimension.)}
        \item $T^2 = 4\pi^2\dfrac{\ell}{g} \Rightarrow g = \dfrac{4\pi^2\ell}{T^2}$.
        \[ g = \dfrac{4\pi^2 \times 1,000}{(2,005)^2} = \dfrac{39,478}{4,0200} \approx \boxed{\num{9,82}~\text{m}\cdot\text{s}^{-2}}. \]
        \item Incertitude relative :
        \[ \dfrac{u(g)}{g} = \sqrt{\left(\dfrac{0,002}{1,000}\right)^2 + \left(2\times\dfrac{0,006}{2,005}\right)^2} = \sqrt{(0,002)^2 + (0,00599)^2} \approx \sqrt{\num{4,0e-6} + \num{3,58e-5}} \approx \num{6,3e-3}. \]
        Donc $u(g) = 0,0063 \times 9,82 \approx \SI{0,062}{m.s^{-2}}$, arrondi à $\SI{0,06}{m.s^{-2}}$ :
        \[ \boxed{g = (\num{9,82} \pm \num{0,06})~\text{m}\cdot\text{s}^{-2}}. \]
        \item Intervalle de la mesure : $[9,76\,;\,9,88]~\text{m}\cdot\text{s}^{-2}$. La valeur de référence $g_{\text{ref}} = \SI{9,78}{m.s^{-2}}$ appartient à cet intervalle : \textbf{la mesure est compatible} avec la valeur tabulée.
    \end{enumerate}
\end{enumerate}
\textbf{Points de vigilance :} diviser par $\sqrt{n}$ et non par $n$ pour l'incertitude-type sur la moyenne ; ne pas oublier le facteur 2 devant l'incertitude relative de $T$ (car $g \propto T^{-2}$) ; aligner le rang décimal de la valeur sur celui de l'incertitude ; pour la compatibilité, comparer la valeur de référence à l'\emph{intervalle}, pas à la seule valeur centrale.
\end{corrige}

\begin{corrige}[Exercice 10 — Type Bac : mesure d'une résistance et incertitudes]
\textbf{Barème indicatif (sur 10) :} A.1) 1,5 pt ; A.2) 1,5 pt ; B.3) 1 pt ; B.4) 2 pts ; B.5) 1,5 pt ; B.6) 2,5 pts.

\textbf{Partie A.}
\begin{enumerate}
    \item Composantes de l'incertitude : $1\,\%$ de la lecture $= 0,01 \times 46,7 = \SI{0,467}{\Omega}$ ; $2$ digits $= 2 \times 0,1 = \SI{0,2}{\Omega}$. Notice « $\pm(\dots + \dots)$ » : on additionne (borne constructeur) :
    \[ \Delta R = 0,467 + 0,2 = 0,667 \approx \boxed{\SI{0,7}{\Omega}}. \]
    \item Incertitude à un chiffre significatif ($\SI{0,7}{\Omega}$, rang du dixième) : $\boxed{R = (\num{46,7} \pm \num{0,7})~\Omega}$. Incertitude relative : $\dfrac{0,7}{46,7} \approx \boxed{1,5\,\%}$.
\end{enumerate}

\textbf{Partie B.}
\begin{enumerate}[resume]
    \item $R = \dfrac{U}{I} = \dfrac{4,52}{0,096} \approx \boxed{\num{47,1}~\Omega}$. \emph{(Convertir $I$ en ampères : $96~\text{mA} = \SI{0,096}{A}$.)}
    \item Incertitudes relatives : $\dfrac{u(U)}{U} = \dfrac{0,05}{4,52} \approx 0,0111$ ; $\dfrac{u(I)}{I} = \dfrac{2}{96} \approx 0,0208$.
    \[ \dfrac{u(R)}{R} = \sqrt{0,0111^2 + 0,0208^2} = \sqrt{\num{1,23e-4} + \num{4,34e-4}} \approx \num{0,0236} \approx 2,4\,\%. \]
    $u(R) = 0,0236 \times 47,1 \approx \SI{1,11}{\Omega}$, arrondi à $\SI{1}{\Omega}$ (rang de l'unité) : $\boxed{R = (\num{47} \pm \num{1})~\Omega}$.
    \item Intervalles : ohmmètre $[46,0\,;\,47,4]~\Omega$ ; volt-ampèremétrique $[46\,;\,48]~\Omega$. Ils se recoupent largement : \textbf{les deux mesures sont compatibles}.
    \item $[R\,I^2] = [R]\,[I]^2 = \left(\text{M}\,\text{L}^2\,\text{T}^{-3}\,\text{I}^{-2}\right)\times\text{I}^2 = \text{M}\,\text{L}^2\,\text{T}^{-3} = [P]$. L'expression $P = RI^2$ est homogène à une puissance. \emph{(On pouvait aussi utiliser $[R] = [U]/[I]$ : $[RI^2] = [U][I]$, dimension d'une puissance.)}
\end{enumerate}
\textbf{Points de vigilance :} bien convertir les mA en A ; pour un quotient, composer les incertitudes \emph{relatives} (somme quadratique) et non les incertitudes absolues ; la compatibilité se conclut sur le recouvrement des intervalles, pas sur l'égalité des valeurs centrales.
\end{corrige}

\begin{corrige}[Exercice 11 — Type Bac : analyse dimensionnelle et validation de modèles]
\textbf{Barème indicatif (sur 10) :} 1) 1,5 pt ; 2.a) 2 pts ; 2.b) 1 pt ; 3) 2,5 pts ; 4.a) 1 pt ; 4.b) 2 pts.

\begin{enumerate}
    \item \textbf{Champ électrique.} $[U] = \text{M}\,\text{L}^2\,\text{T}^{-3}\,\text{I}^{-1}$, $[d] = \text{L}$.
    \[ \left[\dfrac{U}{d}\right] = \text{M}\,\text{L}\,\text{T}^{-3}\,\text{I}^{-1} = [E] \quad \text{(homogène)} ; \qquad [U \times d] = \text{M}\,\text{L}^3\,\text{T}^{-3}\,\text{I}^{-1} \neq [E]. \]
    Seule l'expression $\boxed{E = \dfrac{U}{d}}$ est homogène.
    \item \textbf{Capacité d'un condensateur plan.}
    \begin{enumerate}
        \item De $C = \varepsilon\dfrac{S}{d}$ : $[\varepsilon] = [C]\times\dfrac{[d]}{[S]} = \text{M}^{-1}\,\text{L}^{-2}\,\text{T}^{4}\,\text{I}^{2} \times \dfrac{\text{L}}{\text{L}^2} = \boxed{\text{M}^{-1}\,\text{L}^{-3}\,\text{T}^{4}\,\text{I}^{2}}$.
        \item Le farad a pour dimension $[C] = \text{M}^{-1}\,\text{L}^{-2}\,\text{T}^{4}\,\text{I}^{2}$, donc le farad par mètre a pour dimension $\text{M}^{-1}\,\text{L}^{-3}\,\text{T}^{4}\,\text{I}^{2}$ : \textbf{cohérent} avec le résultat précédent.
    \end{enumerate}
    \item \textbf{Force de Laplace.} $[B] = \text{M}\,\text{T}^{-2}\,\text{I}^{-1}$, $[I] = \text{I}$, $[\ell] = \text{L}$, cible : $[F] = \text{M}\,\text{L}\,\text{T}^{-2}$.
    \begin{align*}
    [F_1] &= \text{M}\,\text{T}^{-2}\,\text{I}^{-1} \times \text{I} \times \text{L} = \text{M}\,\text{L}\,\text{T}^{-2} = [F] \quad \checkmark \\
    [F_2] &= \text{M}\,\text{L}^{-1}\,\text{T}^{-2} \neq [F] \quad \text{(rejetée)} \\
    [F_3] &= \text{M}\,\text{L}^{2}\,\text{T}^{-2} \neq [F] \quad \text{(rejetée)}
    \end{align*}
    Seule $\boxed{F_1 = B\,I\,\ell}$ est dimensionnellement possible.
    \item \textbf{Validation d'un modèle.}
    \begin{enumerate}
        \item $[LC] = \text{T}^2$ (voir exercice 6), donc $\left[2\pi\sqrt{LC}\right] = \text{T}$ : l'expression est homogène à une période.
        \item L'homogénéité est une condition \emph{nécessaire mais non suffisante} : elle ne peut ni fixer la constante $2\pi$, ni garantir que $R$ n'intervient pas (on pourrait avoir $T = 2\pi\sqrt{LC}\,f\!\left(R\sqrt{C/L}\right)$ avec $f$ fonction sans dimension, parfaitement homogène). \textbf{Démarche expérimentale :} réaliser un circuit $(R, L, C)$ série avec $L$ et $C$ fixés, faire varier $R$ (boîte à décades) en restant en régime pseudo-périodique, mesurer la pseudo-période sur l'oscillogramme de $u_C(t)$ pour chaque valeur de $R$, puis tracer $T$ en fonction de $R$ : si $T$ reste constante aux incertitudes près, l'hypothèse est validée dans le domaine exploré ; sinon elle est réfutée (en réalité $T = \dfrac{2\pi}{\sqrt{\frac{1}{LC} - \frac{R^2}{4L^2}}}$, qui tend vers $2\pi\sqrt{LC}$ quand $R$ est faible).
    \end{enumerate}
\end{enumerate}
\textbf{Points de vigilance :} écrire explicitement l'équation aux dimensions avant de conclure ; ne jamais affirmer qu'une expression homogène est « démontrée » ; pour la question 4.b), le jury attend une démarche \emph{expérimentale} précise (grandeur fixée, grandeur variée, grandeur mesurée, critère de validation avec incertitudes).
\end{corrige}

\newpage

\coursec{Fiche bilan}

\begin{popfigure}
\centering
\begin{tikzpicture}[
    mindmap,
    grow cyclic,
    every node/.style={concept, execute at begin node=\hskip0pt},
    concept color=popblue,
    level 1/.append style={level distance=4.5cm, sibling angle=60, font=\small\bfseries},
    level 2/.append style={level distance=3cm, sibling angle=45, font=\footnotesize},
    branch1/.style={concept color=poporange},
    branch2/.style={concept color=popgreen},
    branch3/.style={concept color=poppurple},
    branch4/.style={concept color=poppink},
    branch5/.style={concept color=popgold},
    branch6/.style={concept color=popblue}
]
\node[root concept, text width=3.5cm, align=center] {Mesures, Incertitudes\\ Analyse Dimensionnelle}
    child[branch1] { node[concept, text width=2.8cm, align=center] {Erreur \& Instrument}
        child { node {Erreur absolue/relative} }
        child { node {Sensibilité} }
        child { node {Précision / Fidélité} }
        child { node {Justesse} }
    }
    child[branch2] { node[concept, text width=2.8cm, align=center] {Mesure Unique}
        child { node {$u = \frac{\text{classe}}{2}$} }
        child { node {$u = \frac{\text{appréciation}}{2}$} }
        child { node {Incertitude relative $u_r$} }
    }
    child[branch3] { node[concept, text width=2.8cm, align=center] {Série de Mesures}
        child { node {Moyenne $\bar{x} = \frac{1}{n}\sum x_i$} }
        child { node {Écart-type $\sigma_{n-1}$} }
        child { node {Incertitude-type $u(\bar{x}) = \frac{\sigma}{\sqrt{n}}$} }
    }
    child[branch4] { node[concept, text width=2.8cm, align=center] {Expression Résultat}
        child { node {$x = \bar{x} \pm u$} }
        child { node {Chiffres significatifs} }
        child { node {Arrondi cohérent} }
        child { node {Notation scientifique} }
    }
    child[branch5] { node[concept, text width=2.8cm, align=center] {Système International (SI)}
        child { node {7 unités de base} }
        child { node {Unités dérivées} }
        child { node {Préfixes (p, n, $\mu$, m, k, M, G)} }
        child { node {Nouvelles grandeurs : $a, E, C, L, Z, A$} }
    }
    child[branch6] { node[concept, text width=2.8cm, align=center] {Analyse Dimensionnelle}
        child { node {Dimensions $[L],[M],[T],[I],[\Theta],[N],[J]$} }
        child { node {Homogénéité équation} }
        child { node {Vérification / Recherche formule} }
        child { node {Analyse dimensionnelle complète} }
    };
\end{tikzpicture}
\legende{Carte mentale de la leçon P01 : Mesures, incertitudes et analyse dimensionnelle.}
\end{popfigure}

\begin{bilan}
\courssub{Définitions clés}
\begin{itemize}
    \item \textbf{Erreur} $E = x_{\text{mesuré}} - x_{\text{vrai}}$ (inconnue car $x_{\text{vrai}}$ inaccessible).
    \item \textbf{Incertitude} $u$ : intervalle centré sur la valeur mesurée où se trouve la valeur vraie avec une probabilité donnée (souvent 68\% pour $k=1$).
    \item \textbf{Qualités d'un instrument} :
    \begin{itemize}
        \item \textit{Sensibilité} : rapport $\frac{\text{variation indication}}{\text{variation grandeur}}$.
        \item \textit{Précision (Fidélité)} : dispersion des indications pour une même valeur vraie (liée à l'aléatoire).
        \item \textit{Justesse} : proximité de la moyenne des indications de la valeur vraie (liée au systématique).
    \end{itemize}
    \item \textbf{Mesure unique} : $u$ estimée par la demi-classe d'appareil ($\frac{\text{classe}}{2}$) ou la demi-appréciation de lecture ($\frac{\text{lecture}}{2}$), on retient la plus grande.
    \item \textbf{Série de $n$ mesures} $x_1, \dots, x_n$ :
    \begin{itemize}
        \item Moyenne $\bar{x} = \frac{1}{n}\sum_{i=1}^n x_i$ (meilleure estimation de la valeur vraie).
        \item Écart-type échantillon $\sigma_{n-1} = \sqrt{\frac{1}{n-1}\sum_{i=1}^n (x_i - \bar{x})^2}$ (dispersion).
        \item \textbf{Incertitude-type} $u(\bar{x}) = \frac{\sigma_{n-1}}{\sqrt{n}}$ (incertitude sur la moyenne).
    \end{itemize}
    \item \textbf{Incertitude relative} $u_r = \frac{u}{|x|}$ (souvent en \%). Sans dimension.
    \item \textbf{Chiffres significatifs (CS)} : tous les chiffres certains + le premier incertain. L'incertitude s'arrondit à 1 ou 2 CS ; la valeur mesurée s'arrondit au même rang décimal que $u$.
    \item \textbf{Dimension} : exposition d'une grandeur en produit de puissances des 7 grandeurs de base : $[L], [M], [T], [I], [\Theta], [N], [J]$.
    \item \textbf{Homogénéité} : une équation physique est correcte seulement si les deux membres ont la \emph{même dimension} (et la même unité).
\end{itemize}

\courssub{Formules à connaître par cœur}
\begin{center}
\begin{tabularx}{\linewidth}{|l|X|X|}\hline
\rowcolor{popblueL}
\textbf{Grandeur} & \textbf{Formule} & \textbf{Unité / Dimension} \\ \hline
Incertitude mesure unique & $u = \max\left(\frac{\text{classe}}{2}, \frac{\text{appréciation}}{2}\right)$ & Même unité que $x$ \\ \hline
Moyenne (série) & $\bar{x} = \frac{1}{n}\sum_{i=1}^n x_i$ & Même unité que $x$ \\ \hline
Écart-type (série) & $\sigma_{n-1} = \sqrt{\frac{1}{n-1}\sum_{i=1}^n (x_i - \bar{x})^2}$ & Même unité que $x$ \\ \hline
Incertitude-type (moyenne) & $u(\bar{x}) = \frac{\sigma_{n-1}}{\sqrt{n}}$ & Même unité que $x$ \\ \hline
Incertitude relative & $u_r = \frac{u}{|x|}$ & Sans dimension (\% ou $10^{-2}$) \\ \hline
Résultat final & $x = \bar{x} \pm u(\bar{x})$ & Arrondi cohérent \\ \hline
Dimensions de base & Longueur $[L]$, Masse $[M]$, Temps $[T]$, Courant $[I]$, Tempé. $[\Theta]$, Quantité $[N]$, Luminosité $[J]$ & --- \\ \hline
Exemples nouvelles grandeurs & Accélération $[L][T]^{-2}$, Champ électrique $[M][L][T]^{-3}[I]^{-1}$, Capacité $[M]^{-1}[L]^{-2}[T]^4[I]^2$, Inductance $[M][L]^2[T]^{-2}[I]^{-2}$, Impédance $[M][L]^2[T]^{-3}[I]^{-2}$, Activité $[T]^{-1}$ & Vérifier par analyse dim. \\ \hline
\end{tabularx}
\end{center}

\courssub{Méthodes express}
\begin{enumerate}
    \item \textbf{Estimer $u$ (mesure unique)} : Lire la classe de l'appareil (ex: multimètre classe 1\% $\to$ $u = 1\% \times \text{valeur lue}$) OU estimer l'appréciation de lecture (ex: règle graduée 1 mm $\to$ $u = 0,5$ mm). Prendre le max.
    \item \textbf{Traiter une série} : Calculer $\bar{x}$ (calculatrice mode STAT), $\sigma_{n-1}$ (touche $\sigma_n$ ou $s_x$), puis $u(\bar{x}) = \sigma_{n-1}/\sqrt{n}$. Vérifier valeurs aberrantes (critère de Chauvenet ou $3\sigma$).
    \item \textbf{Exprimer le résultat} : Arrondir $u$ à 1 CS (ou 2 si le 1er est 1 ou 2). Arrondir $\bar{x}$ au même rang décimal. Écrire $\bar{x} \pm u$ (unité SI). Calculer $u_r$ en \%.
    \item \textbf{Analyse dimensionnelle} : Remplacer chaque grandeur par sa dimension. Simplifier. Vérifier $[Gauche] = [Droite]$. Permet de retrouver l'exposant manquant dans une loi de type $y = k x^a$ ou de détecter une erreur de formule.
\end{enumerate}
\end{bilan}

\begin{aretenir}[Checklist avant le Bac]
\begin{itemize}
    \item $\square$ Je distingue \textbf{erreur} (inconnue) et \textbf{incertitude} (estimée).
    \item $\square$ Je sais estimer $u$ pour une mesure unique (classe/2 vs lecture/2).
    \item $\square$ Je calcule $\bar{x}$, $\sigma_{n-1}$ et $u(\bar{x}) = \sigma/\sqrt{n}$ pour une série.
    \item $\square$ J'arrondis correctement $u$ (1 ou 2 ch. sign.) et $\bar{x}$ au même rang décimal.
    \item $\square$ J'écris le résultat final sous la forme $x = \bar{x} \pm u$ avec l'unité SI adaptée.
    \item $\square$ Je calcule l'incertitude relative $u_r$ et l'exprime en \%.
    \item $\square$ Je cite les 7 unités de base du SI et leurs symboles (m, kg, s, A, K, mol, cd).
    \item $\square$ Je détermine la dimension des grandeurs du programme (accélération, champ, capacité, inductance, impédance, activité...).
    \item $\square$ Je vérifie l'homogénéité dimensionnelle d'une équation physique.
    \item $\square$ J'utilise l'analyse dimensionnelle pour retrouver une relation ou valider une hypothèse.
    \item $\square$ Je connais les préfixes SI (p, n, $\mu$, m, k, M, G) et convertis correctement.
\end{itemize}
\end{aretenir}

\findecours

\end{document}
